\documentclass{article}
\usepackage{graphicx} % Required for inserting images
\usepackage{James}
\usepackage{dynkin-diagrams}
\title{Forcing and the Continuum Hypothesis}
\author{James Hindmarch}
\setcounter{section}{0}
\date{January 2024}
\begin{document}

\maketitle
\tableofcontents

\newpage
\newlec
\section{Introduction}
\subsection{Independence}
\begin{ex*}
    In Euclidean geometry, Euclid lays out five postulates, and in particular he lays out the fifth postulate: \begin{enumerate}
        \item[(v)] Parallel postulate: For any given line, and any given point not on that line, there is exactly one line through the given point which does not meet the given line.
    \end{enumerate}
    People kind of hated this axiom, it was significantly more complicated than the other axioms, so people wanted to prove it using the other axioms. However, in the 19th Century Gauss, Lobachevsky, and others proved that one can consistently have axioms (i)-(iv), and not (v).\gap So (v) is independent of (i)-(iv).
\end{ex*}
\begin{ex*}
    Let $\varphi := \exists x \, (x^{2}=2)$. We have that $\bbQ$ is a field satisfying $\neg \varphi$, but that $\bbQ[\sqrt{2}]$ is a field satisfying $\varphi$. So $\varphi$ is independent of the field axioms.
\end{ex*}
\begin{ex*}
    The number of roots of $x^{4}=2$, in $\bbQ$ the are no roots to this. In $\bbR$ there are two roots to this, and in $\bbC$ there are 4 roots to this.
\end{ex*}
However, set theory was set up to encapsulate \textit{all} mathematics. So perhaps set theory is immune to this type of problem, and everything can either be proved or disproved.\gap
G\"odel's incompleteness theorems showed that set theory is \textit{not} immune to this type of problem, and the consistency of $\mathsf{ZFC}$ is independent of $\mathsf{ZFC}$ (assuming that $\mathsf{ZFC}$ is in fact consistent).
\subsection{Continuum problem}
Now cast your mind back to Cantor, and recall \begin{fancythm*}{Cantor}
    \[|\bbN| < \calP(\bbN).\]
\end{fancythm*}
Then we have a natural question, is there a set $X$ such that $|\bbN| < |X| < |\calP(\bbN)|$?
The \textbf{Continuum hypothesis} is the claim that if $X \subseteq \calP(\bbN)$ is infinite, then $|X| = |\bbN|$ or $|X| = |\calP(\bbN)|$. That is to say, $2^{\aleph_{0}} = \aleph_{1}$. This question inspired a lot of maths, with some key results being:
\begin{fancythm*}{Cantor, 1883}
    Any closed subset of $\bbR$ satisfies $\sfCH$.
\end{fancythm*}
\begin{fancythm*}{Alexandrov/Hausdorff, 1916}
    Any Borel set of $\bbR$ satisfies $\sfCH$.
\end{fancythm*}
\begin{fancythm*}{Suslin, 1930}
    Any analytic subset of $\bbR$ satisfies $\sfCH$.
\end{fancythm*}
\begin{fancythm*}{G\"odel, 1938}
    $\Con(\sfZF) \Rightarrow \Con(\sfZFC + \sfCH)$.
\end{fancythm*}
\begin{fancythm*}{Cohen, 1963}
    $\Con(\sfZF) \Rightarrow \Con(\sfZFC + \neg \sfCH)$.
\end{fancythm*}
\subsection{Systems of Set Theory}
The language of set theory, $\calL := \calL_{\in}$, consists of: \begin{itemize}
    \item First-order predicate logic
    \item Two binary symbols: $\in$ and $=$.
    \item Set variables: $\nu_{1}, \dots, \nu_{n}, \dots$
    \item Logical connectives: $\vee$, $\neg$.
    \item Brackets: $($, and $)$.
    \item Existential quantifier: $\exists$
\end{itemize}
\begin{remark*}
    $\wedge$, $\rightarrow$, and $\forall$ can be defined in terms of $\vee$, $\neg$, and $\exists$.
\end{remark*}
\begin{fancydefine}{Free variable/Bound variable}
    A variable is \textbf{bound} in $\varphi$ if it is the occurrence of an $x$ which is in a quantifier, or if it is within the scope of such a quantifier.\gap
    Otherwise, we say a variable is \textbf{free}, or occurs freely. We write $\Fr(\varphi)$ to denote the set of free variable in $\varphi$.
\end{fancydefine}
\begin{remark*}
    \begin{itemize}
        \item We write $\varphi(u_{0}, \dots, u_{n})$ to emphasise the dependence of $\varphi$ on $u_{1}, \dots, u_{n}$.
        \item We allow ourselves to freely chage variables, i.e. $\varphi(v_{0}, \dots, v_{n})$ denotes the same formula with different variables.
        \item We assume the substituted variables are free.
        \item However, writing $\varphi(u_{0}, \dots, u_{n})$ will not imply that $u$ occurs freely, or even occurs at all.
    \end{itemize}
\end{remark*}
\subsection{Theories}
The theories we consider are \begin{itemize}
    \item $\sfZF$ consisting of the axioms: Extension, Pairing, Unions, Empty set, Foundation, Separation, Replacement, Power set, and Infinity.
    \item $\sfZFC$ is $\sfZF$ with the Axiom of Choice.
    \item $\sfZF^{-}$ is $\sfZF$ without the power set axiom. (When we say $\sfZF$ without the power set axiom, we have to change replacement to the axiom scheme of collection, since otherwise we can prove the power set axiom).
    \item $\sfZFC^{-}$ is $\sfZF^{-}$ + Well-ordering principle.
    \item $\sfZ$ is $\sfZF \, \setminus$ Replacement.
\end{itemize}
\begin{remark*}
    Our main background theory will be $\sfZF$.
\end{remark*}
\section{Classes}
\begin{fancydefine*}{Definable class/Proper Class}
    A class $X$ is \textbf{definable} over $\calM$ if there exists a formula $\varphi$, and sets $a_{1}, \dots, a_{n} \in M$ such that \[ \forall z \in \calM \,\, ( z \in X \Leftrightarrow \varphi(z, a_{1}, \dots, a_{n})).\]
    A class $X$ is \textbf{proper} (over $\calM$) if $X \notin \calM$.
\end{fancydefine*}
We will assume that all classes are definable.
\begin{ex*}
    \begin{itemize}\hspace{0.1cm}
        \item \begin{itemize}
                  \item $V = \{x \, : \, x = x\}$
                  \item $R = \{ x \, : \, x \notin x\}$
                  \item $\Ord$
              \end{itemize}
              are all definable proper classes
        \item Any set is a class.
        \item Classes are heavily dependent on the model. If $\calM = \sfZ$ then $\Ord = \calM = \sfZ$.
    \end{itemize}
\end{ex*}
\subsection{Adding defined functions}
The objects $0$, $\perp$, $\subseteq$, $\cap$ all do not exist in our language.
\begin{fancydefine*}{defined $n$-ary predicate/defined $n$-ary function symbol}
    Assume that $\calL \subseteq \calL'$, and $T$ is a set of sentences of $L$. Then \begin{itemize}
        \item $p$ is a \textbf{defined $n$-ary predicate} symbol over $T$ if there is a formula $\varphi$ in $\calL$ such that \[ T \vdash \forall x_{1}, \dots, x_{n} \, ( p(x_{1}, \dots, x_{n} \leftrightarrow \varphi(x_{1}, \dots, x_{n}))).\]
        \item $f$ is a \textbf{defined $n$-ary function symbol} over $T$ if there exists a formula $\varphi$ in $\calL$ such that for any $x_{1}, \dots, x_{n}$, \[ p(x_{1}, \dots, x_{n}) = y\] iff $T \proves \varphi(x_{1},\dots, x_{n},y)$ or $T \proves \forall x_{1}, \dots, x_{n} \exists ! y, \,\,\, \varphi(x_{1} ,\dots, x_{n},y)$.
    \end{itemize}
\end{fancydefine*}
\begin{fancydefine*}{Extension by definitions}
    A set of sentences $T'$ of $calL'$ is an extension by definitions of $T$ over $\calL$ iff $T' = T \cup S$ when \[ S = \{ \varphi_{s} \, | \, s \in \calL' \setminus \calL\}\] and each $\varphi_{s}$ is a definition of $s$ in the language $\calL$ over $T$.
\end{fancydefine*}
\begin{fancyprop}{}
    The following are defined over $\sfZF$: $0$, $1$, $\subseteq$, $\cap$, $\calP$, $\cup$.
\end{fancyprop}
\begin{fancythm*}{}
    Suppose $\calL \subseteq \calL'$, and $T$ is a set of sentences of $\calL$, and $T'$ is an extension by definitions of $T$ over $\calL$. Then we have \begin{enumerate}
        \item[(1)] \underline{Conservativity:} If $\varphi$ is a sentence of $\calL$. Then $T \proves \varphi$ iff $T' \proves \varphi$.
        \item[(2)] \underline{Abbreviations}: If $\varphi$ is a formula of $\calL'$, then there exists a formula $\hat{\varphi}$ of $\calL$, with \[ \Fr(\varphi) = \Fr(\hat{\varphi}),\] and $T \proves \forall x \, ( \varphi \leftrightarrow \hat{\varphi})$.
    \end{enumerate}
\end{fancythm*}
\begin{ex*}
    The intersection of two sets, $a \cap b$, can be defined as \textit{the} set $c$ such that \[ \forall x \, (x \in c \leftrightarrow x \in a \cap x \in b)\]
    This only makes sense if there is a \textit{unique} $c$ satisfying this.\gap
    For example, let \[ \calM = \{ a,c,d, \{a\}, \{a,b\},\{a,b,c\},\{a,b,d\}\}.\] Both $\{a\}$ and $\{a,b\}$ satisfy $\varphi(c)$.
\end{ex*}

\newlec

Assume that there is a countable transitive model of $\sfZFC$, $\calM$. Then $|\bbR \cap M|$ is countable, so $\exists v \in \bbR \setminus M$. So $v$ is a \textit{proper} class over $\calM$, but $v$ is not \textit{definable} over $\calM$ (if it were, then replacement would but it in $\calM$).\gap
\begin{remark*}
    \begin{itemize}
        \item If it's not definable, then we \say{can't talk about it} in $\calL_{\in}$.
        \item So the only proper classes that affect our theory are the definables ones, $\Ord$, $V$, etc.
        \item We can use formulae of the form: \[ \exists C\, (C \text{ is a class} \wedge \forall x ( x \in C \rightarrow \dots)).\] This is an abbreviation for \say{There is a formula $\theta$ such that $\forall x,\, (\theta(x) \rightarrow \dots)$}.
    \end{itemize}
\end{remark*}
Suppose $\calM$ is a model of $\sfZF$. Let $\calD$ be the collection of definable classes over $\calM$. Then $\calD$ is a set in $V$, and $(\calM, \calD)$ is a model of a second-order version of $\sfZF$, which is known as G\"odel-Bernays Set theory.
\section{Absoluteness}
Observe that definitions often appear to give the same set, regardless of which model of $\sfZFC$ we are working in. \begin{exs*}
    \begin{enumerate}
        \item $\{x \, : \, x \neq x\}$ gives the empty set.
        \item $\{ x \, : \, x = a \vee x = b\}$ gives the pair set.
    \end{enumerate}
\end{exs*}
However, other definitions do not give the same set, for example $\calP(\bbN)$ is not the same in the countable transitive model as it is in a more \say{normal} model.\gap
Therefore, to discuss which definitions are definite, we need to define what it means for $\varphi$ to hold in a structure.
\subsection{Relativisation}
\begin{fancydefine}{Bounded Quantifier}
    $\forall x \in a (\dots)$ is an abbreviation for $\forall x (x \in a \rightarrow \dots)$.\gap
    $\exists x \in a, \, \varphi(x)$ is an abbreviation for $\exists x \, (x \in a \wedge \varphi(x))$.
\end{fancydefine}
\begin{fancydefine}{Relativisation}
    Let $W$ be a class, then we define $\varphi^{W}$ by recursion as \begin{itemize}
        \item $(x \in y)^{W} \equiv (x \in y)$
        \item $(x = y)^{W} \equiv (x=y)$
        \item $(\varphi \vee \psi)^{W} \equiv \varphi^{W} \vee \psi^{W}$
        \item $(\neg \varphi)^{W} \equiv \neg (\varphi)^{W}$
        \item $(\exists x \, \varphi)^{W} \equiv \exists x \in W \, \varphi^{W}$.
    \end{itemize}
\end{fancydefine}

\begin{fancyprop}{}
    \begin{align*}
        (\varphi \wedge \psi)^{W}      & \equiv \varphi^{W} \wedge \psi^{W}      \\
        (\varphi \rightarrow \psi)^{W} & \equiv \varphi^{W} \rightarrow \psi^{W} \\
        (\forall x, \varphi(x))^{W}    & \equiv \forall x \in W \varphi^{W}(x).
    \end{align*}
\end{fancyprop}
\proof Exercise.
\endproof
\begin{fancyprop}{}
    Suppose $\calM \subseteq \calN$ and $\calM$ is a definable class over $\calN$, then the relation $\calM \models \varphi$ if first-order expressible in $\calN$
\end{fancyprop}
\proof Suppose $\calM \subseteq \calN$, $\calM$ is definable by $\theta$ (so $\forall z \in \calN, \, \theta(z) \Leftrightarrow z \in \calM).$ \gap
Then we claim that $(\calN, \in) \models \varphi^{\calM}$ iff $(\calM, \in) \models \varphi$. $\calN \models (x \in y)^{\calM}$ iff $\calN \models x \in y$ (and $x,y \in \calM$) iff $\theta(x), \, \theta(y), \langle x, y\rangle \in (\in \cap \calM^{2})$ i.e. $\calM \models x \in y$.\gap
Suppose the claim holds for $\varphi$ and $\psi$, then \begin{align*}
    N \models (\varphi \vee \psi)^{M} & \text{ iff } N \models \varphi^{M} \vee \psi^{M}                          \\
                                      & \text{ iff } N \models \varphi^{M} \text{ or } N \models \psi^{M}         \\
                                      & \text{ iff } \calM \models \varphi^{M} \text{ or } M \models \psi^{\calM}
\end{align*}
Also have that $\calN \models (\exists x \varphi(x))^{\calM}$ iff $\calN \models \exists x \, (x \in \calM \wedge \varphi^{M}(x))$ iff there is some $x \in \calN$ such that $\calN \models x \in \calM$ and $\calN \models \varphi^{M}(x)$ iff there is some $x \in \calN$ such that $\theta(x) \wedge \calM \models \varphi(x)$.
So $\calM \models \exists x \varphi(x)$.
\endproof
\begin{fancydefine}{Upwards/Downwards absolute, absolute}
    Suppose that $\calM \subseteq \calN$ are classes, and $\varphi(u_{1}, \dots, u_{n})$ isa  formula, then
    \begin{itemize}
        \item $\varphi$ is upwards absolute for $\calM$, $\calN$ iff \[ \forall x_{1}, \dots, \forall x_{n} \in M ( \varphi^{M}(x_{1}, \dots, x_{n}) \rightarrow \varphi^{\calN}(x_{1}, \dots, x_{n}))\]
        \item $\varphi$ is downwards absolute for $\calM$, $\calN$ iff \[ \forall x_{1},\dots,x_{n} \in M\, (\varphi^{\calN}(x_{1}, \dots, x_{n}) \rightarrow \varphi^{\calM}(x_{1}, \dots, x_{n}))\]
        \item $\varphi$ is absolute if it is both upwards and downwards absolute.
    \end{itemize}
\end{fancydefine}
\begin{exs*}
    \begin{itemize}
        \item If $\calN = V$, ust say $\varphi$ is absolute for $\calM$
        \item If $\Gamma$ is a set of formulae, then $\Gamma$ is absolute for $\calM$, $\calN$ iff $\varphi$ is absolute for $\calM$, $\calN$.
    \end{itemize}
\end{exs*}
\begin{ex*}
    \begin{itemize}
        \item If $\calM \subseteq \calN$ both satisfy extensionality, then $\emptyset$ is absolute via the formula $\forall x \in a (x \neq x)$.
        \item $\calP(2)$ is not absolute between $4$ and $V$. In $4$ it will not have $\{0,1,2,3\}$ for example.
    \end{itemize}
\end{ex*}
\begin{ex*}
    $\varphi \leftrightarrow \psi$ does not imply $\varphi^{M} \leftrightarrow \psi^{M}$. For example, let $\varphi(v) = \forall x (x \notin v)$ in $\sfZF$ this defines $\emptyset$. Now the following expresses $0 \in z$: \[ \psi(z) \equiv \exists y ( \varphi(y) \wedge y \in z)\]
    \[\theta(z) \equiv \forall y (\varphi(y)) \rightarrow y \in z.\]
    Note that $\exists ! y \, \varphi(y)$, then these are equivalent. However $a=0$, $b=\{0\}$, $c=\{\{\{0\}\}\}$, then let $\calM = \{a,b,c\}$. Then $\varphi^{\calM}(a)$ holds, so $\psi^{M}(b)$, but $\varphi^{M}(c)$ also holds, so $\theta^{\calM}(b)$ fails.
\end{ex*}
The common theme amongst our counterexamples so far has been the lack of transitivity. \begin{fancydefine}{Transitive Class}
    Given classes $\calM \subseteq \calN$, $\calM$ is \textbf{transitive} in $\calN$ if \[ \forall x, y \, (x,y \in N \wedge x \in M \wedge y \in x \rightarrow y \in \calM).\]

\end{fancydefine}
\section{The Levy Hierarchy}
\begin{fancydefine}{$\Delta_{0}$}
    The class $\Delta_{0}$ of formulae is the smallest class $\Gamma$ which is closed under the following: \begin{itemize}
        \item Closure under atomic formulae: \[ \forall i,\,j \, (v_{i} \in v_{j}) \in \Gamma,\]
              and \[ (v_{i} = v_{j}) \in \Gamma.\]
        \item Closure under propositional connectives, i.e. if $\varphi, \psi \in \Gamma$, then $\varphi \vee \psi \in \Gamma$, $\neg \varphi \in \Gamma$.
        \item Closure under \textit{bounded} quantifiers, i.e. for all $i,j$, \[ \varphi \in \Gamma \Rightarrow (\forall v_{i} \in v_{j} \varphi) \in \Gamma,\,\,\,\,\,\,\, \exists v_{i} \in v_{j}. \, \varphi \in \Gamma.\]
    \end{itemize}
\end{fancydefine}
\begin{ex*}
    \item $v_{1} \in v_{2} \vee \forall v_{3} \in v_{4} (v_{4} = v_{5}) \rightarrow \exists v_{1} \in v_{1} (v_{1} = v_{2}) \in \Gamma$ (even though it's meaningless nonsense)
\end{ex*}
Then for the rest of the Levy Hierarchy, we proceed by induction to define:
\begin{fancydefine}{Levy Hierarchy} We set $\Sigma_{0} = \Pi_{0} = \Delta_{0}$.
    \begin{itemize}
        \item If $\varphi$ is $\Pi_{n-1}$, then $\exists v_{i} \varphi$ is $\Sigma_{n}$.
        \item If $\varphi$ is $\Sigma_{n-1}$, then $\forall v_{i} \varphi$ is $\Pi_{n}$.
    \end{itemize}
\end{fancydefine}
\begin{ex*}
    $\forall v_{1} \exists v_{2} \forall v_{3} (v_{4}=v_{7})$ is $\Pi_{3}$. But $\forall v_{1}(v_{1}=v_{2}) \wedge v_{3} \in v_{4}$, or $\forall v_{1}(v_{1} = v_{2} \wedge \exists v_{3}(v_{3} \in v_{4}))$ are not in $\Pi_{n}$ or $\Sigma_{n}$ for any $n$.
\end{ex*}
\begin{fancydefine}{The Levy Hierarchy for a theory}
    Given a $\calL$-theory $\calT$, we let $\Sigma_{n}^{\calT}$ be the class of formulae $\Gamma$ such that for any $\varphi \in \Gamma$, there exists $\psi \in \Sigma_{n}$ such that $\calT \proves \varphi \leftrightarrow \psi$. $\Pi_{n}^{\calT}$ is defined the same way. \gap
    Finally, we say that a formula is $\Delta_{n}^{T}$ if there is a $\psi$ in $\Sigma_{n}$ and a $\theta \in \Pi_{n}$ such that $\calT \proves \varphi \leftrightarrow \psi \wedge \varphi \leftrightarrow \theta$.
\end{fancydefine}
\newlec
\textbf{Warning:} \begin{itemize}
    \item $\exists x_{1} \forall x_{2} \exists x_{3} \forall y ( y \in v \rightarrow v \neq v)$ is $\Sigma_{4}$ as written, but it is clearly logically equivalent to $\forall y \in v (v \neq v)$ which is $\Sigma_{0}$.
    \item In $\sfZ$, it can be the case that you have $\varphi$ being $\Sigma_{1}^{\sfZ}$ but $\forall x \in a.\, \varphi$ is not $\Sigma_{1}^{\sfZ}$.
    \item $\Delta_{n}$ only ever makes sense in the context of a theory $\calT$ for $n>0$.
    \item We only work in classical logic in this course, but in intuitionistic logic, these classes are very badly behaved, since you can have $\varphi$ being $\Sigma_{1}^{\calT}$ but $\neg \varphi$ not being $\Pi_{1}^{\calT}$, and in fact these classes will not cover the whole universe.
\end{itemize}
\begin{fancylem}{}\hspace{0.1cm}

    \begin{itemize}
        \item If $\varphi$ and $\psi$ are in $\Sigma_{n}^{\sfZF}$ then so are $\exists v_{i} \varphi,\, \varphi \wedge \psi, \varphi \vee \psi$, as well as $\exists v_{i} \in v_{j} \varphi$, and $\forall v_{i} \in v_{j} \varphi$.
        \item If $\varphi$ is in $\Sigma_{n}^{\sfZF}$, then $\neg \varphi$ is in $\Pi_{n}^{\sfZF}$.
        \item For every $\varphi$ there exists $n$ such that $\varphi$ is in $\Sigma_{n}^{\sfZF}$.
        \item If $\varphi$ is in $\Sigma_{n}^{\sfZF}$ and $m \geq n$ then $\varphi$ is in $\Sigma_{m}^{\sfZF}$.
    \end{itemize}
\end{fancylem}
\proof Example sheet 1. \endproof
\subsection{Absoluteness of $\Delta_{0}$}
\begin{fancythm}{}
    Suppose that $\calM$ is transitive in $\calN$ and $\varphi( \bar{u})$ is a $\Delta_{0}$ formula, then for any $\bar{a} \in \calM$, \[ M \models \varphi(\bar{a}) \text{ iff } \calN \models \varphi(\bar{a}).\]
\end{fancythm}
\proof By induction on the class $\Delta_{0}$. \begin{itemize}
    \item Atomic formulas, propositional connectives are immediate, so we only care about bounded quantifiers (and the only hard one is existential, since the universal quantifier is also immediate).
    \item Suffices to consider $\exists x \in a. \, \varphi$ where $\varphi$ is absolute between $\calM$ and $\calN$. \begin{enumerate}
              \item[$\Rightarrow$] If $\calM \models \exists x \in a.\, \varphi$, then by definition \[ \calM \models \exists x (x \in a \wedge \varphi(x)).\]
                    Fix $b \in \calM$ such that $\calM \models b \in a \wedge \varphi(b)$.  Then $a,b \in \calN$. So $\calN \models b \in a \wedge \varphi(b)$. I.e. $\calN \models \exists x \in a. \, \varphi(x)$.
              \item[$\Leftarrow$] Suppose $\calN \models \exists x \in a. \, \varphi(x)$, where $a \in M$. Then fix $b \in \calN$ such that \[ \calN \models b \in a \wedge \varphi(b).\]
                    Since $\calM$ is transitive in $\calN$, $b \in \calM$. Therefore $\calM \models b \in a \wedge \varphi(b)$, and so $\calM \models \exists x \in a. \, \varphi(x)$, as required.
          \end{enumerate}
\end{itemize}
\endproof
\begin{fancyprop}{}
    The following are $\Delta_{0}^{\sfZF}$ and therefore absolute between transitive models: \begin{enumerate}
        \item $x \subseteq y$.
        \item $a = \{x,y\}$.
        \item $a = (x,y)$.
        \item $a = x \times y$.
        \item $a = \cup b$.
        \item $\Tc(a)$, the transitive closure of $a$.
        \item $x = \emptyset$.
        \item $r$ is a relation.
        \item $r$ is a function.
        \item $r$ is a relation with domain $a$ and range $b$.
        \item $r"a$, where $r"a := \{ y \, : \, \exists x \in a.\, (x,y) \in r\}$.
        \item $r\upharpoonright a$.
    \end{enumerate}
\end{fancyprop}
\proof I think this is laid out well in Jech. \endproof
\begin{remark*}
    $\cf(a)$, \say{$a$ is a cardinal}, $\omega_{1}$ and $y = \calP(x)$ are \textit{not} absolute.
\end{remark*}
\begin{fancylem}{}
    \say{$a$ is finite} is $\Delta_{1}^{\sfZF}$.
\end{fancylem}
\proof Example sheet 1. \endproof
\begin{fancyprop}{}
    Suppose $\calM \subseteq \calN$, then \begin{itemize}
        \item $\Sigma_{1}$ formulae are upwards absolute.
        \item $\Pi_{1}$ formulae are downwards absolute.
    \end{itemize}
\end{fancyprop}
\proof Easy, think about what \say{$\forall$} and \say{$\exists$} mean.
\begin{corr}{}
    $\Delta_{1}^{\sfZF}$ formulae are absolute between transitive models.
\end{corr}
\proof A $\Delta_{1}^{\sfZF}$ statement is $\Pi_{1}^{\sfZF}$, so downwards absolute, and $\Sigma_{1}^{\sfZF}$, so upwards absolute. \endproof
\begin{fancylem}{$\sfZF$}
    The statement \say{$\alpha$ is an ordinal} is absolute.
\end{fancylem}
\proof Exercise. $\alpha$ is an ordinal if it is a transitive set of transitive sets. The latter can be written as \[ \forall \beta \in \alpha\,\, \forall \gamma \in \beta\,\,(\gamma \in \alpha) \wedge \forall \beta \in \alpha \,\, \forall \gamma \in \beta \,\,\ \forall \delta \in \gamma ( \delta \in \beta).\]
That wasn't much of an exercise.
\endproof
\begin{fancylem}{}
    The statement \say{$r$ is a strict total ordering of $\alpha$} is $\Delta_{0}$.
\end{fancylem}
\proof We need to state \begin{itemize}
    \item $r$ is transitive on $a$.
    \item $r$ satisfies trichotomy (on $a$).
    \item $r$ is irreflexive (on $a$).
\end{itemize}
These all use bounded quantifiers.
\endproof
\begin{corr}{}
    The statement \say{$x$ is a transitive set, totally ordered by $\in$} is $\Delta_{0}$.
\end{corr}
\begin{fancylem}{$\sfZF$}
    The statement \say{$r$ is a well ordering on $a$} is $\Delta_{1}^{\sfZF}$.
\end{fancylem}
\proof \begin{itemize}
    \item $\Pi_{1}$: We can say: $r$ is a relation on $a$ and, \[ \forall X (\exists z \in X)(z=z) \wedge X \subseteq a) \rightarrow \exists z \in X \, \forall y \in X \, (y,z) \notin r.\]
    \item For $\Sigma_{1}$, we first need the following claim: a relation is well-founded iff there exists a function \[ f: a \to \Ord \text{ s.t. } (y,x) \in r \Rightarrow f(y) < f(x).\]
          Suppose $f$ is well-founded, then define $f: \alpha \to \Ord$ by $f(x)= \sup\{f(y)+1 \, : \, (y,x) \in r\}$.
    \item For the other direction, take $X \subseteq a$ non-empty, and look at $f"X \subseteq \Ord$. This has a minimal element, which we'll call $\alpha$. Then for any $z \in X$, if $f(z)= \alpha$, then \[ \forall y \in X.\, f(y) \geq \alpha.\] So by definition $(y,z) \notin r$. (Not claiming that $z$ is unique.)\gap
          Then to show $\Sigma_{1}$, note we have \[ \exists f (f \text{ is a function } \wedge \forall u \in \ran(f) (a \in \Ord) \wedge \forall x,y \in a ( (y,z) \in r \rightarrow f(y) \inf(x))).\]
\end{itemize}
\endproof
\begin{fancyprop}{}
    The following are $\Delta_{0}^{\sfZF}$.
    \begin{itemize}
        \item $x$ is a limit ordinal.
        \item $x$ is a successor ordinal.
        \item $x$ is a finite ordinal.
        \item $\omega$.
        \item $0,1,2,\dots,37,\dots$.
    \end{itemize}
\end{fancyprop}
\endproof
\begin{fancyprop}{}
    The following are $\Pi_{1}^{\sfZF}$ and hence downwards absolute: \begin{itemize}
        \item $\kappa$ is a cardinal.
        \item $\kappa$ is regular.
        \item $\kappa$ is a limit cardinal.
        \item $\kappa$ is a strong limit cardinal.
    \end{itemize}
\end{fancyprop}
\endproof

\section{Relativizing Axioms}
\begin{fancylem}{$\sfZF$}
    Suppose $W$ is a transitive class, and $W \neq \emptyset$. Then we have $(\text{Extensionality})^{W}$, $(\text{Empty set})^{W}$, $(\text{Foundation})^{W}$.
\end{fancylem}
\proof \hspace{0.1cm} \begin{itemize}
    \item Extensionality says that \[ \forall x( \forall y ( \forall z ( z \in x \leftrightarrow z \in y) \rightarrow x = y)).\]
          Then we want to relativise this to $W$, so we get: \[\forall x \in W( \forall y \in W ( \forall z \in W ( z \in x \leftrightarrow z \in y)^{W} \rightarrow (x = y)^{W})).\]
          Then we note that the two formulae in the bracket are absolute, so they hold in $W$ if and only if they \textit{actually} hold. Then since $W$ is transitive, if $x, y \in W$, then also $x,y \subseteq W$. Suppose $x \neq y$, then by extensionality in the universe, we can fix some $z \in x$, $z \notin y$. But then $z \in W$, so we have $\exists z \in W ( z \in x \wedge z \notin y)$ which contradicts our assumption.
    \item Foundation is a very similar argument, write it all out, then use foundation in the background universe to show that foundation holds in the inner model.
    \item For empty set, it is even easier. We  first fix $x \in W$ such that $x \cap W = \emptyset$ using the fact that $W$ is non-empty, and that foundation holds in $V$. Since $W$ is transitive, $x \subseteq W$. Therefore $x = \emptyset \in W$. Moreover $(x=\emptyset)$ is $\Delta_{0}$, so it is absolute between transitive models.
\end{itemize}

\newlec

\begin{fancylem}{}\hspace{0.1cm}
    \begin{itemize}
        \item Suppose $W$ is a transitive class, then if \[ \forall x,y \in W, \, \{x,y\} \in W\] then $(\text{Pairing})^{W}$ holds.
        \item If for any $x \in W$, $\cup x \in W$, then $(\text{Union})^{W}$ holds.
        \item If $\omega \in W$, then $(\text{Infinity})^{W}$ holds.
        \item If for every $\varphi$, with $\Fr(\varphi) \subseteq \{x,a, v_{1}, \dots, v_{n}\}$, we have that $\forall a,x_{1}, \dots x_{n} \in W (\{ x \in a | \varphi^{W}(x,a,v_{1}, \dots, x_{n}) \} \in W)$, then separation must hold in $W$.
        \item If for every formula $\varphi$ with $\Fr(\varphi) \subseteq \{x,z,a,v_{1}, \dots, v_{n}\}$, we have that for any $a,v_{1},\dots, v_{n} \in W$. If \[\forall x \in a \exists ! y \in W\,  \varphi^{W}(x,y,a,v_{1}, \dots, v_{n}).\] we have that \[ \exists b \in W \left( \{ y | \exists x \in a. \, \varphi^{W}(x,y,a,v_{1},\dots,v_{n})\} \subseteq \right)\] then replacement holds in $W$.
        \item If $\forall a \in W$, $\exists b \in W$ such that ($\calP(a) \cap W = b$), then $(\text{Power Set})^{W}$ holds.
    \end{itemize}
\end{fancylem}
\begin{corr}{}
    If $W$ is a transitive class satisfying the conditions of the previous lemma, then $(\sfZF)^{W}$ i.e. $W$ is a model of $\sfZF$.
\end{corr}
\section{Transfinite Recursion}
\begin{fancydefine}{Set-like relation}
    A relation $R$ is \textbf{set-like} on a class $A$ iff $\forall x \in A. \, \{ y \in A\, |\, yRx\}$ is a set.
\end{fancydefine}
\begin{ex*}\hspace{0.1cm}
    \begin{itemize}
        \item $\in$ is set-like.
        \item Any relation on a set will be set-like.
    \end{itemize}
\end{ex*}
\begin{fancydefine}{Absolute class}
    Let $A$ be a class and fix $\varphi$ such that $A = \{ x \,| \, \varphi(x)\}$. Then $A^{W} = \{x \, | \, \varphi^{W}(x)\}$. We will say that $A$ is \textbf{absolute} if $A^{W} = A \cap W$.
\end{fancydefine}
\begin{fancydefine}{Absolute class-relation}
    View a class relation $R \subseteq V \times V$ as a collection of ordered pairs $\{ (x,y) \, | \, \psi(x,y)\}$. Then $R^{W} = \{(x,y) \, | \, \psi^{W}(x,y)\}$. Say $R$ is \textbf{absolute} for $W$ iff $R^{W} = R \cap (W \times W)$.
\end{fancydefine}
Observe that if $R$ is a function, we can only refer to the function $R^{W}$ if we first check $(\forall x \exists! y. \, \psi(x,y))^{W}$. If we relativise a function, we will assume that we have already checked this. In this case $R^{W} : W \to W$ and $R$ is an absolute function for $W$ iff $R^{W} = R \upharpoonright W$.
\begin{fancythm}{Transfinite Recursion}
    Let $R$ be a relation which is well-founded and set-like on a class $A$, and let \[ F : A \times V \to V\] be a function. Given $x \in A$, let $\pred(A,x,R) = \{ y \in A \, | \, yRx\}$. Then there is a unique function $G: A \to V$ such that \[ \forall x \in A. \, G(x)=F(x,G \upharpoonright \pred(A,x,R)).\]
\end{fancythm}
\begin{fancythm}{Absoluteness of transfinite recursion}
    Let $R$ be a relation which is well-founded and set-like on a class $A$. Let $F : A \times V \to V$ be a class function and let $G : A \to V$ be the unique function given by transfinite recursion. Let $W$ be a transitive model of $\sfZF$, and suppose that \begin{enumerate}
        \item $A$ and $F$ are absolute between $W$ and the universe $V$.
        \item $R$ is also absolute for $W$ $(R \text{ is set-like on } A)^{W}$.
        \item $\forall x \in W$, $\pred(A,x,R) \subseteq W$.
    \end{enumerate}
    Then $G$ is absolute for $W$.
\end{fancythm}
\proof By absoluteness, we have that $A^{W} = A \cap W$, and $R^{W} = R \cap (W \times W)$. Therefore every non-empty subset of $A^{W}$ has an $R^{W}$ minimal element, so $(R \text{ is well-founded on } A)^{W}$.\gap
So apply transfinite recurison in $W$ to define a (unique) function \[G^{W} : A^{W} \to W\] such that \[ \forall x \in A^{W}. \, G^{W}(x) = F^{W}(x, G^{W} \upharpoonright \pred^{W}(A^{W},x,R^{W})).\]
To prove $G$ is absolute, it suffices to show $G^{W} = G \upharpoonright A^{W}$. We prove this by transfinite induction in $W$. Suppose that for all $yRx$, $G^{W}(y) = G(y)$, then we can write \begin{align*}G^{W}(x) & = F^{W}(x,G^{W} \upharpoonright \pred^{W}(A^{W},x,R^{w})) \\
                       & = F(x, G \upharpoonright \pred(A,x,R))
\end{align*}
\endproof
\begin{corr}{}
    The following are absolute for transitive models of $\sfZFC$: \begin{itemize}
        \item $\rank(x)$.
        \item $\Tc(x)$.
        \item Ordinal arithmetic operations, $+$ and $\cdot$, since these are both defined via transfinite recursion.
    \end{itemize}
\end{corr}
\section{The Reflection Thoerem}
Recall the Tarski-Vaught test from model theory, which says:
\begin{fancythm*}{Tarski-Vaught Test}
    Let $\calM \subseteq \calN$ as structures, with universes $M,N$. Then TFAE: \begin{enumerate}
        \item[(i)] $\calM \preceq \calN$ ($\calM$ is an elementary substructure of $\calN$).
        \item[(ii)] For any formula $\varphi(v, \bar{w})$ and $\bar{a} \in M$, if there exists $b \in N$ such that \[ \calN \models \varphi(b, \bar{a}),\] then there exists $c \in M$ such that $\calN \models \varphi(c, \bar{a})$.
    \end{enumerate}
\end{fancythm*}
\proof In Model Theory. \endproof
\begin{fancydefine}{Subformula closed}
    A finite list of formulae \[ \bar{\varphi} = \varphi_{1}, \dots, \varphi_{n}\] is said to be \textbf{subformula closed} if every subformula of a formula is on the list.
\end{fancydefine}
\begin{ex*}
    $\varphi_{i} = \exists x.\, \psi_{1}$ and $\varphi_{j} = \psi_{2} \wedge \psi_{3}$, then $\psi_{1},\psi_{2},\psi_{3}$ also appear on the list.
\end{ex*}
\begin{fancylem}{}
    Let $\bar{\varphi} = \varphi_{1}, \dots, \varphi_{n}$ be a subformula closed list, and $W \subseteq Z$ are two non-empty classes, then TFAE: \begin{itemize}
        \item[(i)] $\bar{\varphi}$ are absolute for $W, Z$
        \item[(ii)] Whenever $\varphi_{i}$ is of the form $\exists x.\, \varphi_{j}(x,\bar{y})$ where the free variables of $\varphi_{j}$ are contained in $\{x, \bar{y}\}$, then \[\forall \bar{y} \in W \left( \exists x \in Z. \, \varphi_{j}(x, \bar{y}) \rightarrow \exists x \in W. \, \varphi_{j}(x,\bar{y}) \right).\]
    \end{itemize}
\end{fancylem}
\proof \hspace{0.1cm}
\begin{enumerate}
    \item[(i) $\Rightarrow$ (ii)] Suppose $\varphi_{i} \equiv \exists x. \, \varphi_{j}(x,\bar{y})$, fix $\bar{y} \in W$. Then $\varphi_{i}^{Z}(\bar{y}) \equiv \exists x \in Z. \, \varphi_{j}^{Z}(x, \bar{y})$. So if $\exists x \in Z. \, \varphi_{j}^{Z}(x, \bar{y})$, then, by absoluteness, $\varphi_{1}^{W}(\bar{y})$ holds, i.e. $\exists x \in W. \, \varphi_{j}^{W}(x, \bar{y})$.\gap
          Now, $W \subseteq Z$, and absoluteness of $\varphi$ means that $ \exists x \in W. \, \varphi_{j}(x, \bar{y})$.
    \item[(ii) $\Rightarrow$ (i)] We will prove this by induction on the length of $\varphi_{i}$. \begin{itemize}
              \item Base case is when $\varphi_{i}$ is atomic or of the form $\varphi_{j} \vee \varphi_{k}$, or $\varphi_{i} \equiv \neg \varphi_{j}$ are immediate.
              \item Suppose $\varphi_{i} \equiv \exists x \varphi_{j}(x, \bar{y})$, and fix $\bar{y} \in W$. Then \begin{align*}
                        \varphi_{i}^{Z}(\bar{y}) & \Leftrightarrow \exists x \in Z. \, \varphi_{j}^{Z}(x, \bar{y})                                                             \\
                                                 & \Leftrightarrow \exists x \in W.\, \varphi_{j}^{Z}(x, \bar{y}) \tag{$\Rightarrow$ assumption, $\Leftarrow$, $W\subseteq Z$} \\
                                                 & \Leftrightarrow \varphi_{i}^{W}(\bar{y})
                    \end{align*}
          \end{itemize}
\end{enumerate}
\endproof
\newlec
\begin{fancythm}{Reflection Theorem}
    Let $W$ be a non-empty class, and suppose there is a class function $F_{W}$ such that for any $\alpha \in \Ord$, $F_{W}(\alpha) = W_{\alpha} \in V$. Assume that \begin{enumerate}
        \item[(i)] If $\alpha < \beta$, then $W_{\alpha} \subseteq W_{\beta}$.
        \item[(ii)] If $\lambda$ is a limit ordinal, then $W_{\lambda} = \cup_{\alpha<\lambda} W_{\alpha}$
        \item[(iii)] $W = \cup_{\alpha \in \Ord} W_{\alpha}$
    \end{enumerate}
    Then for any finite collection of formulae $\bar{\varphi} = \varphi_{1}, \dots, \varphi_{n}$. $\sfZF \proves \forall \alpha \exists \beta>\alpha (\beta \text{ is a limit ordinal} \wedge \bar{\varphi} \text{ are absolute for } W_{\beta}, \,W)$.
\end{fancythm}
\proof Let $\bar{\varphi}$ be a finite list of formulae, then WLOG, we can assume $\bar{\varphi}$ is subformula closed and there are no universal quantifications (i.e. we always phrase $\forall x.\, \psi$ as $\neg \exists x. \, \neg \psi$). \gap
Then if $\varphi_{i} \equiv \exists x.\, \varphi_{j} (x,\bar{y})$ where $\bar{y}$ is a tuple of length $k_{i}$, we define \[F_{i} : W^{k_{i}} \to \Ord\] by \[ F_{i}(\bar{y}) =
    \begin{cases}
        0    & \text{ if $\neg \exists x \in W. \, \varphi_{j}^{W}(x, \bar{y})$}                                        \\
        \eta & \text{ when $\eta$ is the least ordinal such that $\exists x \in W_{\eta}\, \varphi_{j}^{W}(x,\bar{y})$}
    \end{cases}\]
Write $G_{i}(\delta) = \sup \{ F_{i}(\bar{y}) \, | \, \bar{y} \in W_{\delta}^{k_{i}}\}$. We have the following: \begin{itemize}
    \item If $\varphi_{i}$ is not of the above form, then $G_{i}(\delta) = 0$ for all $\delta$.
    \item Finally, let $K(\delta) = \max\{\delta + 1,\, G_{1}(\delta), \dots, G_{n}(\delta)\}$.
\end{itemize}
Note that $F_{i}$ work in an analogous way to Skolem functions, and that $F_{i}$ is a well-defined function and, by replacement, since $W_{\delta}$ is a set, $F_{i}"W_{\delta}^{k_{i}} \in V$.
$G_{i}$ is monotone, so if $\delta < \delta'$, then $G_{i}(\delta) \leq G_{i}(\delta')$. Then we claim that \[ \forall \alpha \, \exists \beta > \alpha \, (\beta \text{ is a limit ordinal} \wedge \forall \delta < \beta, \, \forall i \leq n.\,\, G_{i}(\delta) < \beta).\]
To prove this, set $\lambda_{0} = \alpha$ and $\lambda_{t+1} = K(\lambda_{t})$. Then $\beta = \sup_{t \in \omega} \lambda_{t}$. If $\delta < \beta$, then $\delta < \lambda_{t}$ for some $t$, so $G_{i}(\delta) \leq G_{i}(\lambda_{t}) \leq K(\lambda_{t}) = \lambda_{t+1} < \beta$. \gap
To complete the theorem, suppose $\varphi_{i} \equiv \exists x.\, \varphi_{j}(x,\bar{y})$. Now fix $\bar{y} \in W_{\beta}$, Now suppose that $\exists x \in W. \, \varphi_{j}^{W}(x, \bar{y})$, since $\beta$ is a limit, $\bar{y} \in W_{\gamma}$ for some $\gamma < \beta$. Thus $0 < F_{i}(\bar{y}) \leq G_{i}(\gamma)< \beta$, so by construction $\exists x \in W_{\beta}.\, \varphi_{j}(x, \bar{y})$, hence $\bar{\varphi}$ is absolute for $W_{\beta},\,W$.
\endproof

\begin{corr}{Montague-L\'evy Reflection}
    For any finite list of formulae $\bar{\varphi}$: \[ \sfZF \proves \forall \alpha \exists \beta > \alpha (\bar{\varphi}\text{ are absolute for }V_{\beta})\]
\end{corr}

Warning: Reflection is a theorem scheme, so for any choice of $\varphi_{1}, \dots, \varphi_{n}$, it is a theorem of $\sfZF$ that $\varphi_{1},\dots, \varphi_{n}$ are absolute for some $V_{\beta}$. We do not and can not have $\sfZF \proves \forall \bar{\varphi} \forall \alpha \exists \beta > \alpha. \, (\bar{\varphi} \text{ is  absolute for $W_{\beta}$, $W$.})$. \gap
Also, $\bar{\varphi}$ is absolute for $W_{\beta}, W$ does not imply that $(\bar{\varphi})^{W_{\beta}}$, i.e. it doesn't mean that they actually hold.
\begin{remark}
    If $\bar{\varphi}$ is any finite list of axioms of $\sfZF$ then there are arbitrarily large $\beta$ such that $\bar{\varphi}$ holds in $V_{\beta}$. But, if $\beta$ is a limit cardinal, then $V_{\beta} \models \sfZ (+ \text{Choice})$, so we can restrict $\bar{\varphi}$ to instances of replacement.
\end{remark}
\begin{corr}{$\sfZF$}
    Let $\calT \supseteq \sfZF$ be a set of axioms in $\calL_{\in}$, and $\varphi_{1}, \dots, \varphi_{n}$ a finite list of axioms from $\calT$. Then \[\calT \proves \forall \alpha \exists \beta > \alpha. \, (\bigwedge_{i=1}^{n} \varphi_{i})^V_{\beta}.\]
\end{corr}
\begin{corr}{$\sfZFC$}
    Let $W$ be a class, $\bar{\varphi}$ a finite list of formulae in $\calL_{\in}$. Then \begin{align*} \sfZFC \proves \forall x \subseteq W ( & \Trans(x) \rightarrow \exists y. (x \subseteq y \wedge \Trans(y)                           \\
                                                      & \wedge \bar{\varphi} \text{ are absolute for $y,W$} \wedge |y| \leq \max\{ \omega ,|x|\}))\end{align*}
    Taking $x=\omega$, $W=V$.
\end{corr}
\begin{corr}{}
    Let $\calT$ be any set of sentences in $\calL_{\in}$ such that $T \proves \sfZFC$, and let $\varphi_{1}, \dots, \varphi_{n} \in \calT$. Then \[\calT \proves \exists y. \left(\Trans(y) \wedge |y| = \omega \wedge \left( \bigwedge_{i=1}^{n} \right)^{y}\right).\]
\end{corr}
\begin{corr}{}
    Let $\calT$ be any consistent set of sentences in $\calL_{\in}$ such that $\calT \proves \sfZF$. Then $\calT$ is not finitely axiomatizable. That is, for any finite set of sentences $\Gamma$ in $\calL_{\in}$ such that $\calT \proves \Gamma$ there exists a sentence $\varphi$ such that $T \proves \varphi$ but $\Gamma \notproves \varphi$.
    This is only true for first-order theories, e.g. G\"odel-Bernays set theory (second-order set theory) is finitely axiomatizable.
\end{corr}
\proof Let $\varphi_{1}, \dots, \varphi_{n}$ be a set of sentences such that \[ \calT \proves \bigwedge_{i=1}^{n} \varphi_{i}.\] Suppose for a contradiction that from $\varphi_{1}, \dots, \varphi_{n}$, one can prove every axiom of $\calT$. Then by reflection $\calT \proves \forall \alpha. \, \exists \beta > \alpha. \, \left( \left(\bigwedge_{i=1}^{n} \varphi_{i}\right)^{V_{\beta}} \leftrightarrow \left(\bigwedge_{i=1}^{n} \varphi_{i} \right) \right)$.
Fix $\beta_{0}$ to be the least ordinal such that $\bigwedge_{i=1}^{n} \varphi_{i}^{V_{\beta_{0}}}$. Then all axioms of $\calT$ hold in $V_{\beta_{0}}$. (I.e. $V_{\beta_{0}} \models \calT$). Since $\calT$ extends $\sfZF$, basic absoluteness results hold. So if $\alpha \in V_{\beta_{0}}$, then $V_{\alpha}^{V_{\beta}} - V_{\alpha} \wedge V_{\beta_{0}} = V_{\alpha}$. SO $V_{\alpha}$ is absolute for $V_{\beta_{0}}$. Since $\calT$ proves $\exists \alpha \bigwedge_{i=1}^{n} \varphi_{i}^{V_{\alpha}}$. Since $V_{\beta_{0}}$ satisfies every axiom of $\calT$, this must be true in $V_{\beta_{0}}$. So $\exists \alpha < \beta_{0}. \, \bigwedge_{i=1}^{n} \varphi_{i}^{V_{\alpha}}$. Contradicting the minimality of $\beta_{0}$.\endproof
\newlec
\section{Cardinal Arithmetic}
For this section, we will work in $\sfZFC$, since in $\sfZF$ it's quite hard to define what a cardinal even is.
\begin{fancydefine}{Cardinality}
    The \textbf{cardinality} of $x$, written $|x|$ is the least ordinal $\alpha$ such that there is a bijection between $\alpha$ and $x$.
\end{fancydefine}
\begin{fancydefine}{Cardinal operations}
    Let $\kappa$ and $\lambda$ be cardinals, then \begin{itemize}
        \item $\kappa + \lambda = |\{0\} \times \kappa \cup \lambda \times \{1\}|$.
        \item $\kappa \cdot \lambda = |\kappa \times \lambda|$
        \item $\kappa^{\lambda} = |^{\lambda}\kappa| = |\{f \, : \, f: \lambda \to \kappa\}|$
        \item $\kappa^{<\lambda} = \sup\{\kappa^{\alpha} \, : \, \alpha \in \Card, \, \alpha< \lambda\}$.
    \end{itemize}
\end{fancydefine}
\begin{fancythm}{Hessenberg}
    If $\kappa$ and $\lambda$ are infinite, then $\kappa+ \lambda = \kappa \cdot \lambda = \max\{\kappa, \, \lambda\}$.
\end{fancythm}
\begin{fancylem}{}
    If $\kappa, \lambda, \mu$ are infinite cardinals, then \begin{itemize}
        \item $\kappa^{\lambda + \mu} = \kappa^{\lambda} \cdot \kappa^{\mu}$.
        \item $\left(\kappa^{\lambda}\right)^{\mu} = \kappa^{\lambda \mu}$.
    \end{itemize}
\end{fancylem}
\begin{fancydefine}{Cofinal/Cofinality}
    A map $f: \alpha \to \beta$ is \textbf{cofinal} if $\sup(\ran(f)) =\beta$. The cofinality of a limit ordinal $\gamma$, written $\cf(\gamma)$ is the least ordinal $\alpha$ such that there is a cofinal map $f: \alpha \to \gamma$.
\end{fancydefine}
\begin{remark*}\hspace{0.1cm}
    \begin{itemize}
        \item $\cf(\gamma) \leq \gamma$.
        \item $\omega = \cf(\omega) = \cf(\omega + \omega) = \cf(\aleph_{omega})$.
        \item $\cf(\gamma) \leq |\gamma|$.
    \end{itemize}
\end{remark*}
\begin{fancydefine}{Singular/Regular ordinals}
    A limit ordinal is \textbf{singular} if $\cf(\gamma)<\gamma$. If $\cf(\gamma)=\gamma$, then we call the ordinal \textbf{regular}.
\end{fancydefine}
\begin{fancythm}
    Let $\gamma$ be a limit ordinal, then: \begin{itemize}
        \item If $\gamma$ is regular, then $\gamma$ is a cardinal.
        \item $\gamma^{+}$ is a regular cardinal. (The cardinal successor, not the ordinal successor.)
        \item $\cf(\cf(\alpha))=\cf(\alpha)$.
        \item $\aleph_{\alpha}$ is regular whenever $\alpha = 0$ or $\alpha$ is a successor ordinal.
        \item If $\lambda$ is a limit ordinal, then $\cf(\aleph_{\lambda}) = \cf(\lambda)$.
    \end{itemize}
\end{fancythm}
\begin{fancythm}{}
    If $\kappa$ is a regular cardinal, and $\calF$ is a family of sets, with $|\calF|<\kappa$ and $|X|<\kappa$ for all $X \in \calF$, then it will be the case that $|\cup \calF| < \kappa$.
\end{fancythm}
\proof We prove this by induction on $|\calF| = \gamma < \kappa$. Suppose that the claim holds for $\gamma$, and let $\calF = \langle X_{\alpha} \, | \, \alpha < \gamma + 1 \rangle$. Then \begin{align*}
    |\cup F| & = |\cup_{\alpha < \gamma} X_{\alpha} \cup X_{\gamma}           \\
             & = |\cup_{\alpha<\gamma} X_{\alpha}| + | X_{\gamma}|            \\
             & = \max \{ | \cup_{\alpha < \gamma} X_{\alpha}|, |X_{\gamma}|\} \\
             & <\kappa
\end{align*}
Now suppose $\gamma$ is a limit, and the claim holds for all $\beta<\kappa$. Let $\calF = \langle X_{\alpha} \, | \, \alpha<\gamma \rangle$. Define $g : \gamma \to \kappa$ by $g(\beta) = |\cup_{\alpha<\gamma} X_{\alpha}|$, since $\kappa$ is regular, $g(\gamma) = |\cup \calF| < \kappa$.
\endproof.
\begin{fancydefine}
    Let $\kappa_{i}$ for $i \in I$ be an index sequence of cardinal numbers and let $\langle X_{i} \, | i \in I\rangle$ be a sequence of pairwise disjoint sets with $|X_{i}| = \kappa_{i}$ for all $i \in I$. Then the \textbf{cardinal sum} of $\langle \kappa_{i} \, | \, i \in I \rangle$ is $\sum_{i \in I} |X_{i}|$.\gap
    The \textbf{cardinal product} of $\langle \kappa_{i} \, | \, i \in I\rangle$ is \[\prod_{i \in I}\kappa_{i} = |\prod_{i \in I} X_{i}|\] where \[|\prod_{i \in I} X_{i} = \{ f \, | \, f \text{ is a function, } \dom(f) = I, \text{ and } \forall i \in I. \, f(i) \in X_{i}\}.\]
\end{fancydefine}
\begin{fancythm}{K\"onig}
    Let $I$ be an indexing set and suppose that $\kappa_{i} < \lambda_{i}$ for all $i \in I$. Then \[ \sum_{i \in I} \kappa_{i} < \prod_{i \in I} \lambda_{i}.\]
\end{fancythm}
\proof Let $\langle B \, | \, i \in I \rangle$ be a sequence of disjoint sets with $|B_{i}| = \lambda_{i},$ and let $B = \prod_{i \in I} B_{i}$. It will suffice to show that if $\langle A_{i} \, | \, i \in I \rangle$ is a sequence of subsets of $B$ such that for all $i \in I$, $|A_{i}| = \kappa_{i}$, then $\cup_{i \in I} A_{i} \neq B_{i}$.\gap
Given such a sequence, let $S_{i}$ be the projection of $A_{i}$ onto its $i^{\text{th}}$ co-ordinate, \[ S_{i} = \{ f(i) \, | \, f \in A_{i}\}.\]
Then $S_{i} \subseteq B_{i}$ and \[|S_{i}| \leq |A_{i}| = \kappa_{i} < \lambda_{i} = |B_{i}|.\]
So fix a $t_{i} \in B_{i}\setminus S_{i}$. Finally, define $g \in B$, $g(i)=t_{i}$. By construction $g \notin A_{i}$ for all $i$, so $g \in B \setminus \cup_{i \in I} A_{i}$.
\endproof
\begin{corr}{Cantor's Theorem}
    If $\kappa \geq 2$ and $\lambda$ is infinite, then $\kappa^{\lambda} > \lambda$.
\end{corr}
\proof \hspace{0.1cm} \begin{align*}
    \lambda & = \sum_{\alpha < \lambda} 1              \\
            & < \prod_{\alpha<\lambda} 2 \tag{K\"onig} \\
            & =2^{\lambda}                             \\
            & \leq \kappa^{\lambda}.
\end{align*}
\endproof
\begin{corr}{}
    \[ \cf(2^{\lambda})>\lambda.\]
\end{corr}
\proof Let $f: \lambda \to 2^{\lambda}$, we will show that \[ |\cup f"\lambda| < 2^{\lambda}|.\]
Since for all $i \in I$, $f(i) < 2^{\lambda}$, \begin{align*}
    \cup f " \lambda & = \sum_{i < \lambda} f(i)        \\
                     & < \prod_{i< \lambda} 2^{\lambda} \\
                     & = (2^\lambda)^{\lambda}          \\
                     & = 2^{\lambda \cdot \lambda}      \\
                     & = 2^{\lambda}
\end{align*}
\endproof
\begin{corr}{}
    $2^{\aleph_{0}} \neq \kappa$ for any $\kappa$ of cofinality $\aleph_{0}$. (In partiular, $2^{\aleph_{0}} \neq \aleph_{\omega}$.)
\end{corr}
\begin{corr}{}
    $\kappa^{\cf(\kappa)}>\kappa$ for every infinite cardinality $\kappa$.
\end{corr}
\subsection{Cardinal Exponentiation}
\begin{fancydefine}{Generalized Continuum Hypothesis}
    The \textbf{GCH} says that fpr every cardinal $\kappa$, $2^{\kappa} = \kappa^{+}$, i.e. $2^{\aleph_{\alpha}}=\aleph_{\alpha+1}$.
\end{fancydefine}
\begin{fancythm}{}
    Assume that GCH holds, and $\kappa, \lambda$ are infinite cardinals. Then there are three cases for cardinal exponentiation \begin{enumerate}
        \item[(i)] $\kappa \leq \lambda$ means that $\kappa^{\lambda} = \lambda^{+}$.
        \item[(ii)] If $\cf(\kappa) \leq \lambda < \kappa$, then $\kappa^{\lambda}=\kappa^{+}$.
        \item[(iii)] If $\lambda < \cf(\kappa)$, then $\kappa^{\lambda}=\kappa$.
    \end{enumerate}
\end{fancythm}
Without GCH, we know much much less. the following thorem is essentially the only $\sfZFC$-provable restriction for regular cardinals.
\begin{fancythm}{}
    If $\kappa, \lambda$ are cardinals: \begin{itemize}
        \item[(i)] If $\kappa < \lambda$, then $2^{\kappa} \leq 2^{\lambda}$.
        \item[(ii)] $\cf(2^{\kappa})>\kappa$.
        \item[(iii)] If $\kappa$ is a limit cardinal, then $2^{\kappa} = (2^{<\kappa})^{\cf(\kappa)}$.
    \end{itemize}
\end{fancythm}
\begin{fancythm}{}
    Let $\kappa$, $\lambda$ be infinite cardinals, then \begin{enumerate}
        \item[(i)] If $\kappa \leq \lambda$ then $k^{\lambda}=2^{\lambda}$.
        \item[(ii)] If there exists some $\mu < \kappa$ such that $\mu^{\lambda}\geq \kappa$, then $\kappa^{\lambda} = \mu^{\lambda}$.
        \item[(iii)] If $\kappa < \lambda$ and $\mu^{\lambda}<\kappa$ for all $\mu<\kappa$, then we have two possibilites:
              \begin{enumerate}
                  \item[(a)] If $\cf(\kappa)>\lambda$ then $k^{\lambda} = \kappa$.
                  \item[(b)] If $\cf(\kappa) \leq \lambda$ then $\kappa^{\lambda} = \kappa^{\cf(\kappa)}$.
              \end{enumerate}
    \end{enumerate}
\end{fancythm}
\begin{fancythm}{Silver}
    Suppose that $\kappa$ is singular, $\cf(\kappa)>\aleph_{0}$, and $2^{\alpha}=\alpha^{+}$ for all $\alpha< \kappa$, then $2^{\kappa} = \kappa^{+}$.
\end{fancythm}
This says, essentially, that GCH cannot first break at a cardinal of countable cofinality.
\begin{remark*}{}
    It is consistent (relative to large cardinals) to have $2^{\aleph_{n}}=\aleph_{n+1}$ for all $n \in \omega$, but $2^{\aleph_{\omega}}=\aleph_{\omega+2}$.
\end{remark*}
\begin{fancythm}{Shelah}
    Suppose that $2^{\aleph_{n}}<\aleph_{\omega}$ for all $n \in \omega$, then $ 2^{\aleph_{\omega}}< \aleph_{\omega_{4}}.$
\end{fancythm}
There is a big open question: Can we improve this bound? Can we prove, as conjectured, that $2^{\aleph_{\omega}}<\aleph_{\omega_{1}}$.
\newlec
\section{Constructibility}
One of our aims is to prove that $\Con(\sfZF) \Rightarrow \Con(\sfZFC + \sfGCH)$. Constructibility will help us to do this. \gap
Recall that $V_{\alpha+1} = \calP(V_{\alpha})$. G\"odel's idea was to restrict this to \say{nice} subsets.
\begin{fancydefine}{Definable set}
    A set $x$ is said to be \textbf{definable} over $(\calM, \in)$ if there exists $a_{1}, \dots, a_{n} \in \calM$ and a formula $\varphi$ such that \[ x = \{ z \in \calM \, | \, (\calM, \in) \models \varphi(z,a_{1}, \dots, a_{n})\}.\]
    We can also write: \[ \Def(\calM) = \{ x \in \calM \, | \, x \text{ is definable over } \calM\}\]
\end{fancydefine}
Some observations about this definition:
\begin{itemize}
    \item $\calM \in \Def(\calM)$.
    \item $\calM \subseteq \Def(\calM) \subseteq \calP(\calM)$
\end{itemize}
However, this definition needs formalising as it stands.
\gap
\textbf{One method to formalise:} Code formulas by elements of $V_{\omega}$ using G\"odel codes. One then uses Tarski's satisfaction relation to define a formula $\Sat$ such that \[ \Sat(\calM,E, \lceil \varphi \rceil, x_{1}, \dots, x_{n}) \leftrightarrow (\calM, E) \models \varphi(x_{1},\dots,x_{n})\]
where $\lceil \varphi \rceil$ is the G\"odel code for $\varphi$. We won't use this definition though.\gap
\textbf{Another method to formalise:} \begin{fancydefine}{The Constructible Hierarchy}
    Define $L_{\alpha}$ by transfinite recursion as: \begin{itemize}
        \item $L_{0} = \emptyset$
        \item $L_{\alpha+1} = \Def(L_{\alpha})$
        \item $L_{\lambda} = \cup_{\alpha < \lambda} L_{\alpha}$ when $\lambda$ is a limit.
        \item $L = \cup_{\alpha \in \Ord} L_{\alpha}$.
    \end{itemize}
\end{fancydefine}
\begin{fancylem}{}
    For any ordinals $\alpha, \beta$: \begin{enumerate}
        \item[(i)] $\beta < \alpha \rightarrow L_{\beta} \subseteq L_{\alpha}$.
        \item[(ii)] $\beta < \alpha \rightarrow L_{\beta} \in L_{\alpha}$.
        \item[(iii)] $\Trans(L_{\alpha})$.
        \item[(iv)] $\alpha = \Ord \cap L_{\alpha}$.
        \item[(v)] $\Trans(L)$ and $\Ord \subseteq L$.
    \end{enumerate}
\end{fancylem}
\begin{fancydefine}{Inner model}
    Let $\calT$ be a set of axioms of $\calL_{in}$, and $W$ be a class. $W$ is called an \textbf{inner model} of $\calT$ if: \begin{enumerate}
        \item[(i)] $\Trans(W)$
        \item[(ii)] $\Ord \subseteq W$.
        \item[(iii)] $(T)^{W}$, that is, for every formula $\varphi$ in $T$, $\varphi^{W}$.
    \end{enumerate}
\end{fancydefine}
\begin{fancythm}{}
    $L$ is an inner model of $\sfZFC$.
\end{fancythm}
\proof \hspace{0.1cm}\begin{itemize}
    \item Extensionality and foundation: Since $L$ is transitive, $L$ satisfies these.
    \item Empty set: $\emptyset^{L} = \emptyset = L_{0} \in L$.
    \item Pairing: Given $a,b \in L$, have to prove $\{a,b\} \in L$. Fix $\alpha$ such that $a,b \in L_{\alpha}$. Then $\{a,b\} = \{ x \in \alpha \, : \, L_{\alpha} \models x = a \vee x=b\}$. So $\{a,b\} \in \Def(L_{\alpha})$.
    \item Unions: Fix $a \in L_{\alpha}$. So $\cup \alpha \subseteq L_{\alpha}$. Then \[ \cup a = \{ x \in L_{\alpha} \, | \, (L_{\alpha}, \in) \models \exists z. \, (z \in a \wedge x \in z) \} \in \Def(L_{\alpha}).\]
    \item Infinity: $\omega = \{ n \in L_{\omega} \, | \, (L_{\omega}, \in) \models n \in \Ord\}$.
\end{itemize}
These are all the \say{easy} axioms. Now for the harder ones: \begin{itemize}
    \item Separation: Let $\varphi$ be a formula, $a, \bar{u} \in L_{\alpha}$. Then we claim that $b= \{ x \in a \, | \, \varphi^{L}(x,bar{u})\} \in L$. Using the reflection theorem, we can find $\beta > \alpha$ such that \[ \sfZF \proves \forall x \in L_{\beta} \left( (\varphi)^{L}(x, \bar{u}) \leftrightarrow \varphi^{L}(x, \bar{u})\right)\] Moreover, $\varphi^{L_{\beta}}$ holds iff $(L_{\beta},\in) \models \varphi(x, \bar{u})$. Therefore \begin{align*}
               & \{x \in a \, | \, \varphi^{L}(x, \bar{u})\}                                                        \\
               & = \{x \in a \, | \, \varphi^{L_{\beta}}(x, \bar{u})\}                                              \\
               & = \{ x \in L_{\beta} \, | \, (L_{\beta} , \in) \models \varphi(x, \bar{u}) \} \in \Def(L_{\beta}).
          \end{align*}
    \item Replacement: It suffices to prove that if $\alpha \in L$ and $f: \alpha \to L$ is a function, then there exists $\gamma \in \Ord$ such that $f"a \subseteq L_{\gamma}$. Observe that for every $x \in a$, there exists $\beta \in \Ord$ such that $f(x) \in L_{\beta}$. So by replacement in $V$, then there exists an ordinal $\gamma$ such that for every $x \in a$, there exists $\beta \in \gamma$ such that $f(x) \in L_{\beta}$. Since $\beta \in \gamma \Rightarrow L_{\beta} \subseteq L_{\gamma}$, we have that $\forall x \in a. \, f(x) \in L_{\gamma}$.
    \item Power set: Suffices to prove that if $x \in L$, then $\calP(x) \cap L \in L$. Take $x \in L$, then using replacement in $V$ fix $\gamma \in \Ord$ such that $\calP(x) \cap L \subseteq L_{\gamma}$. Then \[\calP(x) \cap L = \{ z \in L_{\gamma} \, | \, (L_{\gamma}, \in) \models z \subseteq x \} \in \Def(L_{\gamma})\]
\end{itemize}
\endproof
\subsection{G\"odel functions}
Note for clarity: $(a,b,c) = (a,(b,c))$.
\begin{fancydefine}{G\"odel functions}\hspace{0.1cm}
    \begin{itemize}
        \item $\calF_{1}(x,y) = \{x,y\}$
        \item $\calF_{2}(x,y) = \cup x$
        \item $\calF_{3}(x,y) = x \setminus y$
        \item $\calF_{4}(x,y) = x \times y$
        \item $\calF_{5}(x,y) = \dom(x) = \{ \text{first} z \, | \, z \in x, \, z \text{ an ordered pair} \}$
        \item $\calF_{6}(x,y) = \ran(x) = \{ \text{second} z \, | \, z \in x, \, z \text{ an ordered pair} \}$
        \item $\calF_{7}(x,y) = \{(u,v,w) \, | \, (u,v) \in x, \, w \in y\}$
        \item $\calF_{8}(x,y) = \{(u,w,v) \, | \, (u,v) \in x \, w \in y\}$
        \item $\calF_{9}(x,y) = \{ (v,u) \in y \times x \, | \, u=v\}$
        \item $\calF_{10}(x,y) = \{ (v,u) \in y \times x \, | \, u \in v\}$.
    \end{itemize}
\end{fancydefine}
\begin{fancyprop}{}
    The following can all be written as a finite combination of $\calF_{1}$ to $\calF_{8}$: $\{x\}, \, x \cup y, \, x \cap y, \, (x,y), \, (x,y,z).$ For example $\{x\} = \calF_{1}(x,x)$.
\end{fancyprop}
\begin{fancyprop}{}
    For every $i \leq 10$, $z = \calF_{i}(x,y)$ can be written using a $\Delta_{0}$-formula.
\end{fancyprop}
\begin{fancylem}{G\"odel's normal form}
    For every $\Delta_{0}$ formiula $\varphi(x_{1},\dots, x_{n})$ with $\Fr(\varphi) \subseteq \{x_{1},\dots, x_{n}\}$, there is a term $\calF_{\varphi}$ built from the symbols $\calF_{1} - \calF_{10}$ such that \[ \sfZF \proves \forall a_{1}, \dots, a_{n}.\, \calF_{\varphi}(a_{1}, \dots, a_{n})=\{(x_{n},\dots, x_{1}) \in a_{n}\times \dots \times a_{1} \, | \, \varphi(x_{1},\dots,x_{n})\}\]
\end{fancylem}
\begin{remark*}\hspace{0.1cm}\begin{itemize}
        \item We reverse the ordering in the G\"odel normal form for technical reasons of being easier to work with.
        \item \begin{itemize}
                  \item $\calF_{2}$ will correspond to $\vee$
                  \item $x \cap y$ will correspond to $\wedge$
                  \item $\calF_{3}$ will correspond to $\neg$
                  \item $\calF_{9}, \calF_{10}$ are atomic.
              \end{itemize}
    \end{itemize}
    \item $\calF_{7}, \, \calF_{8}$ deal with ordered $n$-tuples. $(x_{1},x_{2},x_{3})$ formed using $x_{1}$ and $(x_{2},x_{3})$. It canno be formed using $(x_{1},x_{2})$ and $x_{3}$, so we have $\calF_{7}$ and $\calF_{8}$.
\end{remark*}
\begin{fancydefine}{Closed under G\"odel functions}
    A class $C$ is \textbf{closed under G\"odel functions} iff \[ \calF_{i}(x,y) \in C \text{ whenever } x,y \in C.\]
    Given a set $b$, we write $\cl(b)$ for the smallest set which contains $b$ as subset and which is closed under G\"odel functions.
\end{fancydefine}
\begin{fancydefine}{}
    Let $b$ be a set, define $\calD^{n}(b)$ inductively. \begin{itemize}
        \item $\calD^{0}(b) = b$.
        \item $\calD^{n+1}(b) = \{ \calF_{i}(x,y) \, | \, x,y \in \calD^{n}(b),\, i \leq 10 \}$.
    \end{itemize}
    Note that $\cl(b) = \cup_{n \in \omega} \calD^{n}(b)$.
\end{fancydefine}
\begin{fancylem}{}
    If $\calM$ is a transitive class, closed under G\"odel functions, then $\calM$ satisfies $\Delta_{0}$-separation.
\end{fancylem}
\proof Let $\varphi(x_{1},\dots, x_{n})$ be $\Delta_{0}$, and let $a,b_{1},\dots,b_{i-1},b_{i+1},\dots,b_{n} \in \calM$.\gap
Let $Y = \{ x_{i} \in a \, | \, \varphi(b_{1}, \dots, b_{i-1}, x_{i} b_{i+1}, \dots, b_{n})\}$. Let $\calF_{\varphi}$ be the formula from G\"odel's normal form, then for any $c_{1},\dots, c_{n} \in \calM$: \[ F_{\varphi}(c_{1},\dots,c_{n}) = \{(x_{n}, \dots, x_{1}) \in c_{n} \times \dots \times c_{1} \, | \, \varphi(x_{1},dots,x_{n})\} \in \calM.\]
So, since $\{b_{j}\} = \calF_{1}(b_{j},b_{j}) \in \calM$, $\calF_{\varphi}(\{b_{1}\}, \dots, \{b_{i-1}\}, a, \{b_{i+1}\}, \dots, \{b_{n}\}) \in \calM$. Then you can show $Y \in \calM$ by taking $\calF_{6}(\ran)$ $n-i$ times, and then taking $\calF_{5}$ ($\dom$). \endproof


\begin{fancythm}{}
    For every transitive set $\calM$, \[\Def(\calM) = \cl(\calM \cup \{ \calM \}) \cap \calP(\calM)\].
\end{fancythm}
\newlec
\proof $(\subseteq)$: Let $\varphi$ be a formula. Then $\varphi^{\calM}$ is $\Delta_{0}$. Therefore, there is a term $\calG$ built from $\calF_{1}, \dots, \calF_{10}$ such that for $a_{1},\dots, a_{n} \in \calM$:
\begin{align*}
    \{x \in \calM \, | \, (\calM, \in) \models \varphi(x,a_{1},\dots, a_{n})\} & = \{ x \in \calM \, | \, \varphi^{\calM}(x,a_{1},\dots,a_{n})\} \\
                                                                               & = \calG(\calM,a_{1},\dots, a_{n})                               \\
                                                                               & \in \cl(M \cup \{M\})                                           \\
\end{align*}
$(\supseteq)$: First, we have the following claim: \gap \textbf{Claim:} If $\calG$ is built from $\calF_{1},\dots, \calF_{10}$, then for any $x,a_{1},\dots,a_{n}$, the statements \begin{align*}
    x = \calG(a_{1},\dots,a_{n}) &  & x \in \calG(a_{1},\dots,a_{n})
\end{align*}
are $\Delta_{0}$. To prove this, suppose that $X,Y \in \calD^{k}(a_{1},\dots,a_{n})$. Then
\begin{align*}
    x=\calF_{1}(X,Y) \leftrightarrow \forall z \in x (z=X \vee z=Y) \wedge \exists w \in x ( w = X) \wedge \exists w' \in Y (w'=Y)
\end{align*}
$x \in \calF_{1}(X,Y) \leftrightarrow x \in \{X,Y\} \leftrightarrow x=X \vee x=Y$. Similar arguments work for $i \in \{2,3,\dots,10\}$.\gap
Assuming the claim, let $Z \in \cl(M \cup \{M\}) \cap \calP(M)$. Fix $\calG$ to be a term built from $\calF_{1},\dots,\calF_{10}$ such that \[ Z = \calG(\calM, a_{1},\dots, a_{n}).\]
Let $p$ bea $\Delta_{0}$ formula such that \[ x \in \calG(\calM, a_{1},\dots, a_{n}) \text{ iff } \varphi(x,\calM, a_{1}, \dots, a_{10}).\]
Then $\calG(\calM, a_{1}, \dots a_{n}) = \{ x \in M \, | \, \varphi(x,\calM,a_{1},\dots,a_{n})\}$. We need a formula $\psi$ such that \[ \psi^{\calM} \leftrightarrow \varphi(x,\calM, a_{1},\dots, a_{n}).\]
For example, define $\psi$ from $\varphi$ by the following replacements: \begin{enumerate}
    \item[(i)] $\exists v_{i} \in \calM$ is replaced by $\exists v_{i}$
    \item[(ii)] $\forall v_{i} \in \calM$ is replaced by $\forall v_{i}$.
    \item[(iii)] $v_{i} \in \calM$ replaced by $v_{i}=v_{i}$.
    \item[(iv)] $\calM = \calM$ replaced by $v_{0}=v_{0}$
    \item[(v)] $\calM \in \calM$, $\calM \in v_{i}$, $M =v_{i}$ replaced by $\neg (v_{0} = v_{0})$.
\end{enumerate}
Then \begin{align*} Z & = \calG(\calM,a_{1},\dots, a_{n})                                                          \\
                 & = \{ x \in \calM \, | \, (\calM, \in) \models \psi(x,a_{1},\dots,a_{n})\} \in \Def(\calM).
\end{align*}
\endproof
Now, recall lemma 8.9, that there is a G\"odel normal form for all $\Delta_{0}$ functions. We will prove this now:
\proof We will call a formula $\varphi$ a \textbf{termed-formula} (or a t-formula) if the conclusoon of the lemma holds for $\varphi$. We will only use the logical symbols $\vee, \wedge, \neg, \exists$. The only occurrence of existentials will be in formulas of the form: \[\varphi(x_{1}, \dots, x_{n}) \equiv \exists x_{m+1} \in x_{j} \psi(x_{1}, \dots, x_{m+1})\] where $j \leq m \leq n$. For example: \[ \varphi(x_{1},x_{2},x_{3},x_{4}) \equiv \exists x_{3} \in x_{1} (x_{1} \in x_{2} \wedge x_{3}=x_{1})\] is allowed, but \[ \exists x_{1} \in x_{2} \psi\] is not, since the containment of elements is going up, and \[ \exists x_{3} \in x_{1} (x_{3} \in x_{2} \vee \exists x_{4} \in x_{1} \psi)\] is not because $\psi$ needs to be a statement of at most $3$ variables.\gap
Note that every $\Delta_{0}$ formula is equivalent to one satisfying these assumptions. We allow for dummy variables, so $\varphi(x_{1},x_{2}) = x_{1} \in x_{2}$ and $\varphi(x_{1},x_{2},x_{3}) = x_{1} \in x_{2}$ are distinct. There are four sections of the proof \begin{enumerate}
    \item Logical points
    \item Propositional connectives
    \item Atomic formulas
    \item Existentially Bounded
\end{enumerate}
\textbf{Section 1: Logical Points} \begin{enumerate}
    \item[(a)] If $\sfZF \proves \varphi(\bar{x}) \leftrightarrow \psi(\bar{x})$ and $\varphi$ is a t-formula, then $\psi$ is a t-formula, since \[ \calF_{\psi}(\bar{a}) = \{ \bar{x} \in \bar{a} \, | \, \varphi(\bar{x})\} = \calF_{\varphi}(\bar{a})\]
    \item[(b)] For all $m,n$ if $\varphi(x_{1},\dots, x_{n}) \equiv \psi(x_{1},dots,x_{m})$ and $\psi$ is a t-formula, then so is $\psi$. There are a few cases: \begin{enumerate}
              \item[(i)] $n \geq m$. Prove by induction on $n$. If $n=m$ then trivial. $\varphi(x_{1},\dots,x_{n+1}) \equiv \psi(x_{1},\dots,x_{m})$ then $\varphi(x_{1},\dots, x_{n+1}) \equiv \theta(x_{1},\dots,x_{n})$ when $\theta$ is a $t$-formula. Then \[ \calF_{\varphi}(a_{1},\dots, a_{n},a_{n+1}) = a_{n+1} \times \calF_{\theta}(a_{1},\dots, a_{n}) = \calF_{4}(a_{n},\calF_{\theta}(a_{1},\dots,a_{n})).\] This is why we reverse the ordering, to make this section easier.
              \item[(ii)] If $n \leq m$, also by induction. $n=m$ is similarly trivial. If \[\varphi(x_{1},\dots,x_{n-1}) \equiv \psi(x_{1},\dots,x_{m})\] then \[ \varphi(x_{1},\dots, x_{n-1}) \equiv \theta(x_{1},\dots, x_{n})\] and $\{0\} = \{\calF_{3}(a_{1},a_{1})\} = \calF_{1}(\calF_{3}(a+_{1},a_{1}),\calF_{3}(a_{1},a_{1}))$. Then \begin{align*}
                        \calF_{\varphi}(a_{1},\dots,a_{n-1}) & = \{ (x_{n-1}, \dots, x_{1}) \in a_{n-1} \times \dots \times a_{1} | \varphi(x {1},\dots, x_{n-1})\}                              \\
                                                             & = \ran( \{(0,x_{n-1},\dots,x_{1})   \in \{0\} \times a_{n-1} \times \dots \times a_{1} \, | \, \theta(x_{1},\dots,x_{n-1},0) \} ) \\
                                                             & = \calF_{6}(\calF_{\theta}(a_{1},\dots, a_{n-1}, \{0\}), a_{1})
                    \end{align*}
                    and $\{0\}$ can be written using G\"odel functions so we're done.
          \end{enumerate}
    \item[(c)] If $\psi(x_{1},\dots,x_{n})$ is a t-formula, and $\varphi(x_{1},\dots,x_{n+1}) \equiv \psi(x_{1},\dots, x_{n-1}, x_{n+1}\setminus x_{n})$ then $\varphi$ is a t-formula. \gap
          Firstly, if $n=1$, then $\psi(x_{1})$ is a t-formula. Consider $\psi(x_{2}\setminus x_{1})$. Then \begin{align*}
              \calF_{\varphi}(a_{1},a_{2}) & = \{ (x_{2},x_{1}) \in a_{2}\times a_{1} \, | \, \psi(x_{2})\}                   \\
                                           & = \{ (x_{2},x_{1}) \, | \, x_{1} \in _{1} \wedge x_{2} \in \calF_{\psi}(a_{2})\} \\
                                           & = \calF_{\psi}(a_{2}) \times a_{1}                                               \\
                                           & = \calF_{4}(\calF_{\psi}(a_{2}),a_{1}).
          \end{align*}
          Then if $n>1$ \begin{align*}
              \calF_{\varphi}(a_{1}, \dots, a_{n+1}) & = \{ (x_{n+1},\dots, x_{1}) \in a_{n+1} \times \dots a_{1} \, | \, \\ &\hspace{1cm} x_{n} \in a_{n} \wedge (x_{n+1},x_{n-1},\dots,x_{1}) \in \calF_{4}(a_{1},\dots, a_{n-1}, a_{n+1}) \}\\ \hspace{0.1cm}\\
                                                     & = \calF_{8}(\calF_{4}(a_{1},\dots,a_{n-1},a_{n+1}),_{n})
          \end{align*}
    \item[(d)] If $\psi(x_{1},x_{2})$ is a t-formula, and $\varphi(x_{1},\dots,x_{n}) \equiv \psi(x_{n-1}\setminus x_{1}, x_{n}\setminus x_{2})$ then $\varphi$ is a t-formula. Assume that $n>2$, then \begin{align*}
              \calF_{\varphi}(a_{1},\dots,a_{n}) & = \{ (x_{n}, \dots, x_{1}) \in a_{n} \times \dots \times a_{1} \, | \, (x_{n}, x_{n-1}) \in \calF_{\psi}(a_{n-1},a_{n}) \} \\
                                                 & = \calF_{7}(\calF_{\psi}(a_{n-1},a_{n}),a_{n-2} \times \dots \times a_{1})
          \end{align*}
          \textbf{Section 2: Propsitional Connectives}
    \item[(e)] If $\varphi$ is a t-formula, so is $\neg \varphi$: \[ \calF_{\neg \varphi}(a_{1},\dots, a_{n}) = (a_{n} \times \dots \times a_{1}) \setminus \calF_{\varphi}(a_{1},\dots,a_{n}).\]
    \item[(f)] $\varphi$, $\psi$ are t-formulas, then so is $\varphi \vee \psi$: \[ \calF_{\varphi \vee \psi}(\bar{a}) = \calF_{\varphi}(\bar{a}) \cup \calF_{\psi}(\bar{a})\]
    \item[(g)] $\varphi$, $\psi$ are t-formulas, then so is $\varphi \wedge \psi$. Same as above but intersection instead of union.\\\hspace{0.1cm}\\
          \textbf{Section 3: Atomic formulas}
    \item[(h)] \newlec The formula $\varphi(x_{1},\dots,x_{n}) \equiv x_{i}=x_{j}$ is a $t$-formula for all $i,j\leq n$. There are n cases:
          \begin{enumerate}
              \item[(i)] $i=1$, $j=2$. Then: \begin{align*}
                        \calF_{9}(a_{1},a_{2}) & = \{(x_{2},x_{1}) \in a_{2} \times a_{1} \, | \, x_{1}=x_{2}\}
                    \end{align*}
                    So $\calF_{\varphi}$ isa formula using $\calF_{9}$ and (b) (which lets us add dummy variables).
              \item[(ii)] $j \geq i$. By induction. If $i=j$, then \begin{align*}
                        \calF_{\varphi}(a_{1},\dots,a_{n}) & = a_{n} \times \dots \times a_{1}                                                  \\
                                                           & = \{ (x_{n},\dots,x_{1}) \in a_{n} \times \dots \times a_{1} \, | \, x_{1}=x_{1}\}
                    \end{align*}
                    Suppose $j=i+1$, then let $\theta(x_{1}, \dots, x_{i+1}) = (x_{1}=x_{2})[x_{i}/x_{1}, x_{i+1}/x_{2}]$. This is a t-formula by case (i) and (d), substitution. Then we get $\calF_{\varphi}$ by adding dummy variables.\gap
                    Then for the general case, $\varphi \equiv x_{i} = x_{j+1}$. By (c), $\varphi(x_{1},\dots,x_{j+1}) \equiv (x_{i}=x_{j})[x_{j+1}/x_{j}]$ is a t-formula.\gap
              \item[(iii)] If $i > j$, then $x_{i}=x_{j}$ is logically equivalent to $x_{j}=x_{i}$ which is a t-formula by case (ii).
          \end{enumerate}
    \item[(i)] The statement $\varphi(x_{1},\dots,x_{n}) = x_{i} \in x_{j}$ is a t-formula for all $i,j\leq n$. Then $\calF_{\varphi}$ is formed using $\calF_{10}$ and (b) in the same way as in case (h).\gap
          If $\varphi \equiv x_{i} \in x_{i}$, then $\calF_{\varphi}(a_{1},\dots, a_{n}) = \emptyset = a_{1} \setminus a_{1}$. \gap
          For the general case, \[\psi(x_{1},\dots, x_{n+2}) \equiv (x_{i} = x_{n+1}) \wedge (x_{j}=x_{n+2}) \wedge (x_{n+1} \in x_{n+2})\]
          Then we have that $x_{n+1} \in x_{n+2}$ is a t-formula, since it is given by $(x_{1} \in x_{2})[x_{n+1}/x_{1},\, x_{n+2}/x_{2}]$. So $\psi$ is a t-formula. Then \begin{align*} \calF_{\varphi} & = \ran(\ran(\{x_{n+2}, \dots, x_{1} \in a_{j} \times a_{i} \times a_{n} \times \dots \times a_{1} \, | \, x_{i} = x_{n+1, \, x_{j}=x_{n+2}, x_{i} \in x_{j}} \} )) \\
                               & = \calF_{6}(\calF_{6}(\calF_{\psi}(a_{1},\dots, a_{n},a_{i},a_{j}),a_{1}),a_{1})
          \end{align*}
          \textbf{Section 4: Bounded quantifiers}
          Recall that the only occurrence of $\exists$ will be in formulas of the form: \[ \varphi(x_{1},\dots,x_{n}) \equiv \exists x_{m+1}\in x_{1}. \, \psi(x_{1},\dots,x_{m+1}),\] where $j \leq m \leq n$.
    \item[(j)] If $\psi(x_{1},\dots,x_{n+1})$ is a t-formula, then so is $\varphi(x_{1},\dots,x_{n}) \equiv \exists x_{n+1} \in x_{j}.\, \psi(x_{1},\dots,x_{n+1})$. Let $\theta(x_{1},\dots,x_{n+1}) \equiv x_{n+1} \in x_{1}$. Then $\theta \wedge \psi$ is a t-formula. Now \begin{align*}
              \calF_{\theta \wedge \psi}(a_{1},\dots,a_{n},\calF_{2}(a_{j},a_{j})) & = \calF_{\theta \wedge \psi}(a_{1},\dots,a_{n},\cup a_{j})                                                                    \\
                                                                                   & = \{(x_{n+1}, \dots, x_{1}) \, | \, x_{n+1} \in x_{j}, \, \forall k \leq n (x_{k} \in a_{k}), \, \psi(x_{1},\dots, x_{n+1})\}
          \end{align*}
          Then
          \begin{align*}
              \ran(\calF_{\theta \wedge \psi}) & = \{ (x_{n},\dots,x_{1}) \in a_{n} \times \dots \times a_{1} \, | \, \exists u.\, (u,x_{n},\dots,x_{1}) \in \calF_{\theta \wedge \psi}(a_{1},\dots,a_{n},\cup a_{j})\} \\
                                               & = \{ (x_{n},\dots,x_{1}) \in a_{n} \times \dots \times a_{1} \, | \, \exists x_{n+1} \in x_{j}. \, \psi(x_{1},\dots, x_{n+1})\}.
          \end{align*}
\end{enumerate}
\endproof
\section{Axiom of constructibility}
\begin{fancydefine}{Axiom of constructibility}
    The \textbf{axiom of constructibility} is the assertion \say{$V=L$} which is equivalent to $\forall x \exists \alpha \in \Ord. \, (x \in L_{\alpha})$.
\end{fancydefine}
Our aim is to show that being constructible is absolute.
\begin{fancylem}{}
    The statement $Z = \cl(M)$ is $\Delta_{1}^{\sfZF}$.
\end{fancylem}
\proof It is $\Pi_{1}$ from its definition as the smallest set closed under G\"odel functions: this comes from the statement \[ \forall W ((M \cup \{M\} \subseteq W \wedge (\forall x,y \in W. \, \bigwedge_{i\leq 10} \calF_{i}(x,y) \in W )) \rightarrow Z \subset W)\]
Then the statement is also $\Sigma_{1}$ from the inductive definition using $\calD^{n}$. The statement goes like: \begin{align*}
    \exists W. \, (\func(W) \wedge \dom(W)=W \wedge Z = \cup\ran(W) \wedge W(0)=M \wedge W(n) \subseteq W(n+1) \wedge (\forall x,y \in W(n).\, \wedge_{i \leq 10} \calF_{i}(x,y) \in W(n+1)) \wedge (\forall z \in W(n+1). \, \exists x,y \in W(n). \, \bigvee_{i \leq 10} \calF_{i}(x,y) = z))
\end{align*}
So $W$ is the function that maps $n$ to $\calD^{n}$, and we set $Z$ to be the union of the ranges of these functions.
\endproof
\begin{fancylem}{}
    The function $\alpha \to L_{\alpha}$ is absolute between transitive models of $\sfZF$.
\end{fancylem}
\proof Define $G: \Ord \times V \to V$ by \[ G(\alpha,x) = \begin{cases}
        \cl(x(\beta) \cup \{x(\beta)\}),  & \text{if $\alpha = \beta +1$, $x$ a function with domain $\beta$,} \\
        \cup_{\beta \in \alpha} x(\beta), & \text{$\alpha$ a limit,}                                           \\
        \emptyset                         & \text{o.w.}
    \end{cases}\]
Then $G$ is an absolute function. By transfinite recursion, we have a function $F$ such that \begin{align*}
    F : \Ord  & \to V                                  \\
    F(\alpha) & \mapsto G(x, F \upharpoonright \alpha)
\end{align*}
is absolute. Finally, $F(\alpha)=L_{\alpha}$ for all $\alpha$. \endproof
\begin{fancythm}{}\hspace{0.1cm}
    \begin{enumerate}
        \item[(i)] $L$ satisfies the axiom of constructibility.
        \item[(ii)] $L$ is the smallest inner model of $\sfZF$. I.e., if $M$ is an inner model, then $L \subseteq M$.
    \end{enumerate}
\end{fancythm}
\proof \hspace{0.1cm} \begin{enumerate}
    \item[(i)] Need to show that $(\forall x, \, \exists \alpha \in \Ord. \, (x \in L_{\alpha}))^{L}$. That is, $\forall x \in L, \, \exists \alpha \in \Ord. \, (x \in (L_{\alpha})^{L})$. Since the $L_{\alpha}$ hierarchy is absolute: \[ \forall x. \,(x \in L_{\alpha} \leftrightarrow x \in (L_{\alpha})^{L})\]
          Since $L$ contains every ordinal, if $x \in L$ then $x \in L_{\alpha}$ for some $\alpha$. Thus $x \in (L_{\alpha})^{L}$. So $L \models \alpha \in L \wedge x \in L_{\alpha}$.
    \item[(ii)] Let $M$ be an inner model of $\sfZF$. We construct $L$ inside $M$, $L^{M}$. Then by absoluteness, for every $\alpha \in M \cap \Ord$, $L_{\alpha} = (L_{\alpha})^{M}$. Thus $L_{\alpha} \subseteq M$ for every $\alpha \in M \cap \Ord = V \cap \Ord$. So $L \subseteq M$.
\end{enumerate}
\endproof
\section{The axiom of choice in $L$}
Our aim is to show that a \textbf{strong} version of the axiom of choice holds in $L$. There is a (definable) global well-order. The idea is to define partial well-orderings $<_{\alpha}$ on $L_{\alpha}$ such that $<_{\alpha+1}$ \textbf{end-extends} $<_{\alpha}$. That is, if $y \in L_{\alpha}$ and $x \in L_{\alpha+1}\setminus L_{\alpha}$, then $y <_{\alpha+1} x$. Then we will take $<_{L} = \cup_{\alpha} <_{\alpha}$.
\begin{fancythm}{}
    There exists a well-ordering of the class $L$.
\end{fancythm}
\proof For each $\alpha \in \Ord$ we will construct a well-ordering $<_{\alpha}$ on $L_{\alpha}$ such that if $\alpha < \beta$, then \begin{enumerate}
    \item[(i)] If $x<_{\alpha} y$, then $x <_{\beta} y$.
    \item[(ii)] If $x \in L_{\alpha}$ and $y \in L_{\beta} \setminus L_{\alpha}$ then $x <_{\beta} y$.
\end{enumerate}
In limit cases, we take unions. $<_{\lambda} = \cup_{\alpha < \lambda} <_{\alpha}$. \gap
Then the difficult part is just to define $<_{\alpha+1}$. The idea is to take the ordering on $L_{\alpha}$, then add $\{L_{\alpha}\}$ to the end oft his ordring, followed by elements of $\calD(L_{\alpha} \cup \{L_{\alpha}\}) \setminus L_{\alpha} \cup \{L_{\alpha}\}$. Then by elements of $\calD^{2}(L_{\alpha} \cup \{L_{\alpha}\}) \dots$.\gap
The details of how we do this are as follows, define $<_{\alpha+1}^{n}$ by: \begin{enumerate}
    \item[(1)] $<_{\alpha+1}^{0}$ is the well ordering of $L_{\alpha} \cup \{L_{\alpha}\}$ that end-extends $<_{\alpha}$ by adding $\{L_{\alpha}\}$ as a max element.
    \item[(2)] Suppose $<_{\alpha+1}^{n}$ is defined. Then $<_{\alpha+1}^{n+1}$ is the well-ordering of $\calD^{n+1}(L_{\alpha})$ satisfying: $x <_{\alpha+1}^{n} y$ iff $x <_{\alpha+1}^{n} y$ or $x \in \calD^{n}(L_{\alpha})$ and $y \in \calD^{n+1}(L_{\alpha})\setminus \calD^{n}(L_{\alpha})$, or $x,y \notin \calD^{n}(L_{\alpha})$ and (we choose an ordering on $\calD^{n+1}(L_{\alpha})$):
          \begin{enumerate}
              \item [(a)] The least $i \leq 10$ such that $\exists u,v \in \calD^{n}(L_{\alpha})$ such that $x=\calF_{i}(u,v)$ is less than the least such $j$ such that $y=\calF_{j}(u,v)$.
              \item[(b)] The least such $i$ is equal to the least such $j$, and the $<_{\alpha+1}^{n}$-least $u \in \calD^{n}(L_{\alpha})$ such that $\exists v \in \calD^{n}(L_{\alpha}) (x = \calF_{i}(u,v))$ is less than the $<_{\alpha+1}^{n}$-least $u' \in \calD^{n}(L_{\alpha})$ such that $\exists v' \in \calD^{n}(L_{\alpha})$ such that $y=\calF_{i}(u',v')$.
              \item[(c)] The least such $i$ is equal to the least such $j$, and the $<_{\alpha+1}^{n}$-least such $u'$s are equal and the least $v \in \calD^{n}(L_{\alpha})$... check moodle later today.
          \end{enumerate}
\end{enumerate}
\endproof
\begin{remark*}
    Just to give the main idea of the well-order of $L$, we assume that $<_{\alpha}$ is a well-order of $L_{\alpha}$, then we need to well-order $L_{\alpha+1}$, end-extending $<_{\alpha}$, and it only remains to figure out the order of $\Def(L_{\alpha})\setminus L_{\alpha} \cup \{L_{\alpha}\}$. For this, we can well order $\calD^{n+1}(L_{\alpha} \cup \{L_{\alpha}\})\setminus \calD(^{n}(L_{\alpha} \cup \{L_{\alpha}\})$, (and we start off by setting $D^{0}=\{L_{\alpha}\}$ to be the least element in $<_{\alpha+1}$ which is not in $L_{\alpha}$). Then for $x= \calF_{i}(u,v)$ and $y=\calF_{j}(u',v')$ in $D^{n+1}$, we order according to the least such G\"odel functions that give the ordering. First we check if $i<j$, then we check if $u< u'$, tjen we check if $v<v'$, and otherwise $x=y$.
\end{remark*}
\begin{fancylem}{}
    The relation $<_{L}$ is $\Sigma_{1}$-definable. Moreover, for every limit ordinal $\delta$ and $y \in L_{\delta}$, we have $x<_{L}y$ iff $x \in L_{\delta}$ and $L_{\delta} \models x<_{L}y$.
\end{fancylem}
\proof Example sheet 2.\endproof
Moreover, this gives the axiom of choice, because given $x \in L$, $<_{L}\upharpoonright x$ gives a well-order of $x$.
\section{GCH in $L$}
\begin{fancylem}{$\sfZFC$}
    \hspace{0.1cm} \begin{enumerate}
        \item[(i)] For all $n \in \omega$, $L_{n} = V_{n}$.
        \item[(ii)] If $\calM$ is infinite, $|\calM| = |\Def(\calM)|$,
        \item[(iii)] If $\alpha$ is an infinite ordinal, $|L_{\alpha}|=|\alpha|$.
    \end{enumerate}
\end{fancylem}
\proof These are relatively obvious. \endproof
\begin{fancylem}{G\"odel's condensation lemma}
    For every limit ordinal $\delta$, if $(\calM, \in) \preceq (L_{\delta}, \in)$, then there exists some $\beta \leq \delta$ such that ($\calM, \in) \equiv (L_{\beta},\in)$.
\end{fancylem}
\proof Let $\pi : (\calM, \in) \to (\calN, \in)$ be the Mostowski collapse of $\calM$, and set $\beta = \calN \cap \Ord$. Since $\calN$ is transtitive, $\beta \in \Ord$. We shall prove that $\beta \leq \delta$, and $\calN = L_{\beta}$.\gap
\begin{itemize}
    \item To prove $\beta \leq \delta$, suppose for a contradiction that $\delta< \beta$. Then $\delta \in \calN$. So $\pi^{-1}(\delta) \in \calM$. Now, since being an ordinal is absolute between transitive sets, $\calN \models \delta \in \Ord$. Thus, $\calM \models \pi^{-1}(\delta) \in \Ord$ (but we cannot say right now that $\pi^{-1}(\delta)$ is in fact an ordinal, since $\calM$ isn't transtitive). But $\calM \preceq L_{\delta}$, so $L_{\delta} \models \pi^{-1}(\delta) \in \Ord$. Since $L_{\delta}$ is transitive, $\pi^{-1}(\delta) \in \Ord \cap V$. \gap
          Also, $\calM \models x \in \pi^{-1}(\delta) \Leftrightarrow \calN \models \pi(x) \in \delta$, since $\pi : (\pi^{-1}(\delta) \cap \calM) \to \delta$ is an isomorphism. Therefore the order type, $\otp(\pi^{-1}(\delta) \cap \calM) = \delta$. Then let $f : \delta \to \pi^{-1}(\delta) \cap \calM$ be a strictly increasing enumeration. So for any $\alpha \in \delta$, \[\alpha \leq f(\alpha) < \pi^{-1}(\delta).\]
          So $\delta \leq \pi^{-1}(\delta)$ (since the union of things less than $\delta$ must be $\leq \delta$). On the other hand, $\pi^{-1}(\delta) \in \calM \preceq L_{\delta}$. So $\pi^{-1}(\delta) < \delta$. This gives us a contradiction.
    \item Now to prove that $\beta >0$. Since $L_{\delta} \models \exists x\, \forall y \in x (y \neq y)$, we must have $\calM \models \exists x \forall y \in x (y \neq y)$. Therefore $\calN$ believes this too. So $\emptyset \in \calN \cap \Ord =\beta$.
    \item We also prove that $\beta$ is a limit. Let $L_{\delta} \models \forall \alpha \in \Ord\,\, \exists x. \, (x = \alpha+1)$, then $\calN \models \forall \alpha \in Ord \,\, \exists x. \, (x=\alpha+1)$. Take $\alpha \in \beta$, then $\alpha \in \calN$, so (using absoluteness) $\alpha+1 \in \calN \cap \Ord = \beta$.
    \item Now, towards proving $\calN = L_{\beta}$, we prove $L_{\beta} \subseteq \calN$. We have $L_{\delta} \models \forall \alpha \in \Ord \,\, \exists y.\, (y = L_{\alpha})$. Then $\calN \models \forall \alpha \in \Ord\,\, \exists y.\, (y=L_{\alpha})$. Since the $L_{\alpha}$ hierarchy is absolute. $\forall \alpha \in \calN \cap \Ord = \beta$, $L_{\alpha}\in N$.
    \item Then finally we want to prove $\calN \subseteq L_{\beta}$. We have $L_{\delta} \models \forall x\, \exists y\, \exists z\, (y \in \Ord \wedge z = L_{y} \wedge x \in z)$. Then $\calN \models \forall x\, \exists y\, \exists z\, (y \in \Ord \wedge z = L_{y} \wedge x \in z)$ as well. Fix $a \in \calN$, find $\gamma \in \calN$ and $z \in \calN$ such that $\calN \models \gamma \in \Ord \wedge z=L_{\gamma} \wedge a \in z$. Then by absoluteness, $a \in L_{\gamma} \subseteq L_{\beta}$.
\end{itemize}
\endproof
\begin{fancythm}{}
    If $V=L$, then $2^{\aleph_{\alpha}}=\aleph_{\alpha+1}$ for every ordinal $\alpha$.
\end{fancythm}
\proof Assume that $V=L$. We shall show that $\calP(\omega_{\alpha}) \subseteq L_{\omega_{\alpha+1}}$. Then since $|L_{\omega_{\alpha+1}}| = \aleph_{\alpha+1}$, the theorem will follow. \gap
To do this, it suffices to show that if $X \subseteq \omega_{\alpha}$, then there exists $\gamma \in \omega_{\alpha+1}$ such that $X \in L_{\gamma}$. Take $X \subseteq \omega_{\alpha}$ and let $\delta > \omega_{\alpha}$ be a limit ordinal such that $X \in L_{\delta}$. Now let $\calM$ be an elementary submodel of $L_{\delta}$ such that $\omega_{\alpha} \subseteq \calM$, $X \in \calM$, $|\calM|=\aleph_{\alpha}$ (using downward L\"owenheim-Skolem). \gap
By G\"odel's condenstation lemma, take $\calN$ to be the Mostowski collapse. There is a limit ordinal $\gamma \leq \delta$ such that $\calN = L_{\gamma}$. Since $|\calN| = |\calM| = \aleph_{\alpha}$, $|L_{\gamma}|=\aleph_{\alpha}$, so $\gamma < \omega_{\alpha+1}$. Now $\omega_{\alpha} \subseteq \calM$, the collapsing map is the identity on $\omega_{\alpha}$. Then the map fixed $X$, so $X \in L_{\gamma}$, and we're done. \endproof
\begin{fancythm}{G\"odel}
    \[\Con(\sfZFC) \Rightarrow \Con(\sfZFC + V=L + \sfGCH)\]
\end{fancythm}
\proof We have shown that there is a definable class $L$ such that \[ \sfZF \proves (\sfZFC + V=L + \sfGCH)^{L}.\]
Suppose $\sfZFC + V=L + \sfGCH$ were inconsistent. Then we can fix $\varphi$ such that $\sfZFC + V=L + \sfGCH \proves \varphi \wedge \neg \varphi$. Then \begin{enumerate}
    \item[(i)] $\sfZF \proves (\sfZFC + V=L + \sfGCH)^{L}$.
    \item[(ii)] $\sf \proves (\varphi \wedge \neg \varphi)^{L}$.
    \item[(iii)] By relativisation, $\sfZF \proves \varphi^{L} \wedge (\neg \varphi)^{L}$.
    \item[iv] So, $\sfZF \proves \varphi^{L} \wedge \neg (\varphi^{L})$.
    \item[(v)] Therefore $\sfZF$ is inconsistent.
\end{enumerate}
\begin{fancylem}{Shepherdson}
    There is no class $W$ such that \[\sfZFC \proves \text{\say{$W$ is an inner model} and } (\neg \sfCH)^{W}\]
\end{fancylem}
\proof Omitted \endproof
\begin{fancydefine}{Club}
    Suppose $\Omega$ is either a regular cardinal, or $\Ord$, then $C \subset \Omega$ is said to be a \textbf{club} if it is: \begin{itemize}
        \item Closed: $\forall \gamma \in \Omega \sup(C \cap \gamma) \in C$.
        \item Unbounded: $\forall \alpha \in \Omega \exists \beta \in C (\beta > \alpha)$.
    \end{itemize}
\end{fancydefine}
\begin{fancydefine}{Stationary}
    A class $S \subseteq \Omega$ is \textbf{stationary} if for every club $C \subseteq \Omega$, $C \cap S \neq \emptyset$.
\end{fancydefine}
\subsection{$\diamond$ in $L$}
$\diamond$ is the statement that there is a single sequence of length $\omega_{1}$ which can approximate any subset of $\omega_{1}$.
\begin{fancydefine}{Approximate $\omega_{1}$}
    We say that a sequence $ \langle A_{\alpha} \, : \, \alpha \in \omega_{1}\rangle$ \textbf{approximates }$\omega_{1}$, if it is such that \begin{enumerate}
        \item[(i)] For all $\alpha \in \omega_{1}$, $A_{\alpha} \subseteq \alpha$.
        \item[(ii)] $\forall X \subseteq \omega_{1}$, $\{ \alpha \, : \, X \cap \alpha = A_{\alpha}\}$ is stationary.
    \end{enumerate}
\end{fancydefine}
\begin{fancylem}{}
    $\sfZF \proves \diamond \Rightarrow \sfCH$.
\end{fancylem}
\proof Let $\langle A_{\alpha} \, | \, \alpha \in \omega_{1}\rangle$ be a $\diamond$-sequence. Then $\forall X \subseteq \omega,\, \exists \alpha > \omega. \,\,\, X = A_{\alpha}$. Thus $\{A_{\alpha} \, | \, \alpha \in \omega_{1}, A_{\alpha} \subseteq \omega\} = \calP(\omega)$. \endproof
\begin{fancythm}
    If $V=L$, then $\diamond$ holds.
\end{fancythm}
Note that $\diamond$ is used in certain inductive constructions to build combinatorial objects (e.g. Suslin trees, more in notes).
\subsection{$\square$ in $L$}
\begin{fancydefine}{$\square_{\kappa}$}
    Let $\kappa$ be an uncountable cardinal, then $\square_{\kappa}$ is the assertion that there exists a sequence \[ (C_{\alpha} \, | \, \alpha \in \lim(\kappa^{+}))\] such that \begin{enumerate}
        \item[(i)] $C_{\alpha}$ is a club subset of $\alpha$.
        \item[(ii)] If $\beta \in \lim(C_{\alpha})$, then $C_{\beta} = C_{\alpha} \cap \beta$.
        \item[(iii)] If $\cf(\alpha) < \kappa$, then $|C_{\alpha}| < \kappa$.
    \end{enumerate}
\end{fancydefine}
\begin{fancythm}{Jensen}
    If $V=L$, then $\square_{\kappa}$ holds ofr every uncountable cardinal $\kappa$.
\end{fancythm}
\proof This is really really hard, so we won't prove it. \endproof
\begin{fancylem}{}
    If $\omega_{1}$ holds, then there exists a stationary set \[ S \subseteq \{ \beta \in \omega_{2} \, | \, \cf(\beta) = \omega\}\] such that for all $\alpha \in \omega_{2}$ with $\cf(\alpha)=\omega_{1}$, $S \cap \alpha$ is not stationary in $\alpha$.
\end{fancylem}
\begin{remark*}
    If $\kappa$ is a weakly compact cardinal, then every stationary subset of $\kappa$ reflects (i.e. $\exists \alpha \in \kappa$ such that $S \cap \alpha$ is stationary in $\alpha$). The claim that every stationary subset of $\{\beta \in \omega_{2} \, | \, \cf(\beta) = \omega\}$ reflects at a point of $\cf(\omega_{1})$ is equiconsistent with $\sfZFC + \exists \text{ Mahlo cardinal.}$
\end{remark*}
\newlec
\section{Forcing}
The idea is to widen our model of $\sfZFC$ to \say{add lots of reals.} But if we're working over $V$, then there's nothing to add. Instead, we'll work with countable transitive models (CTMs) of $\sfZFC$\footnote{Eventually we'll be able to remove this assumption, but first things first.}. If $\calM$ is a countable transitive model of $\sfZFC$, then we want to add the $\omega_{2}^{\calM}$-many reals to $\calM$. We want to do this in a \say{minimal} way. For example, we don't want to add any new ordinals to $\calM$. This gives us much more control over the model that we end up building. \gap To give some more intuition for what's happening, recall the argument that the sentence $\varphi \equiv \exists x (x^{2}=2)$ is independent of the axioms of fields. This involved starting with a model of $\neg \varphi$, and then extending to $\bbQ[\sqrt{2}]$ (i.e. extending it in a \say{minimal} way). The key thing about adding $\sqrt{2}$ is that we also have to add everything that's built from $\sqrt{2}$ and $\bbQ$ using the axioms of fields.\gap
There are some possible difficulties that could arise: \begin{itemize}
    \item Suppose that our countable transitive model (CTM) $\calM$, is of the form $L_{\gamma}$ where $\gamma \in \omega$ (this isn't actually possible, which is shown in ES 2, but suppose it were for now). Then $\gamma$ can be coded as a subset of $\omega$, $c$. Now $c \subseteq \omega$, so it can also be viewed as a real, so if we added $c$ to $\calM$, we would also add $\gamma = \Ord \cap \calM$.
    \item Enumerate all formulas as $\{\varphi_{n} \, | \, n \in \omega\}$. Let $r = \{n \, | \, \calM \models \varphi_{n}\}$ (I think this is kind of like $0^{\#}$ but for a CTM). Then we could add a truth predicate to $\calM$, so we would have a truth predicate for $\calM$ constructible inside $\calM$. This is problematic due to Tarski.
\end{itemize}
The main issues to overcome are the following: \begin{enumerate}
    \item We need a method to choose the $\omega_{2}^{\calM}$-many subsets of $\calM$.
    \item Having chosen them, we need to ensure the extension satisfies $\sfZFC$ (the hardest part).
    \item Why should $\omega_{1}^{\calM}$ and $\omega_{2}^{\calM}$ still be cardinals in the extension?
\end{enumerate}
We'll build these reals \textit{within} our model $\calM$. If $r$ is a real, then each of its finite decimal approximations is already in $\calM$. The issue is that within $\calM$, we do not actually \textit{know} what the real we want to add is. If we could say what it was, then using $\sfZFC$, it would already be in $\calM$. Instead, we add a \say{generic} real. \gap
To be generic, we don't want to specify any particular digits (i.e. beginning with $7$). It should, however, contain in its decimal expansion any finite sequence (e.g. \say{746}). \begin{fancydefine}{dense}
    We call a specification \textbf{dense} if any finite approximation can be extended to one satisfying the specification. \gap
    So beginning with a $7$ is not dense, since if $r$ doesn't begin with a $7$, we can't add things to the end to get it to begin with a $7$. However, having any finite sequence in your decimal expansion is dense, since you can always extend to add any finite sequence.
\end{fancydefine}
It will turn out that being generic corresponds to meeting all dense specifications.
\begin{remark*}
    When forcing, we will use countable transitive models. This means we don't get that $\Con(\sfZFC) \Rightarrow \Con(\sfZFC + \neg \sfCH)$, because we've added countable transitive models, but we will then use the reflection theorem to obtain the result we want.
\end{remark*}
\begin{remark*}
    The axiom of choice is not needed in the basic machinery, so we will primarily work over $\sfZF$, and state where choice is used.
\end{remark*}
\subsection{Partial Orders}
The chapter without which every set theory and mathematical logic course is incomplete.
\begin{fancydefine}{Pre-order}
    A \textbf{preorder} is a pair \[ (\bbP, \leq)\] such that \begin{itemize}
        \item $\bbP$ is nonempty.
        \item $\leq$ is a binary relation.
        \item $\leq$ is transitive ($p \leq q \wedge q \leq r \Rightarrow p \leq r$)
        \item $\leq$ is reflexive $(p \leq p)$.
    \end{itemize}
    A preorder is a \textbf{partial order} if $\leq$ is additionallty anti-symmetric $( p \leq q \wedge q \leq p \Rightarrow p=q$).
\end{fancydefine}
\begin{fancydefine}{Forcing poset}
    A \textbf{forcing poset} is a triple $(\bbP, \leq_{\bbP}, \bbone_{\bbP})$, where \begin{itemize}
        \item $(\bbP, \leq_{\bbP})$ is a pre-order
        \item $\bbone_{\bbP}$ is a maximal element.
    \end{itemize}
    Elements of $\bbP$ are called \textbf{conditions}. \begin{itemize}
        \item $q$ is \textbf{stronger} than $p$, if $q \leq_{\bbP} p$ also written that $q$ is an \textbf{extension} of $p$.
        \item $p$ and $q$ are \textbf{compatible}, written $p ||_{\bbP} q$ if there exists $r$ such that $r \leq_{\bbP} p,q$. Otherwise they are called \textbf{incompatible}, written $p \perp_{\bbP} q$.
    \end{itemize}
    In some texts forcing is written the other way round. This is called the Jerusalem notation.
\end{fancydefine}
Also note that $\bbP \in \calM$ abbreviates $(\bbP, \leq_{\bbP}, \bbone_{\bbP}) \in \calM$. (transitivity would give that $\bbone_{\bbP} \in \bbP$, but not necessarily $\leq_{\bbP}$, but whenever we talk about it in this course, it's fine).
\begin{fancydefine}{Separative}
    A pre-order is \textbf{separative} iff for all distinct $p,q$ exactly one of the following two conditions hold: \begin{itemize}
        \item Either $q \leq p$ and $p \nleq q$.
        \item (X)or $q > p$ and there is an $r\leq q$ such that $r \perp p$
    \end{itemize}
\end{fancydefine}
\begin{fancyprop}{}
    If $(\bbP, \leq)$ is a separative pre-order, then $(\bbP, \leq)$ is a partial order.
\end{fancyprop}
\proof Third (probably?) Example sheet.
\begin{fancyprop}{}
    Suppose that $(\bbP, \leq)$ is a pre-order. Define $p \sim q$ by \[ p \sim q \Leftrightarrow \forall r \in \bbP \, (r || p \leftrightarrow r ||q).\]
    Then there is a separative pre-order on $\bbP \setminus \sim$ such that \begin{itemize}
        \item $[p] \perp [q]$ iff $p \perp q$
        \item If $\bbP$ has a maximal element so does $\bbP/\sim$.
    \end{itemize}
\end{fancyprop}
\begin{ex}
    For sets $I$ and $J$, let $\Fn(I,J)$ denote the set of all finite partial functions from $I$ to $J$. That is \[ \Fn(I,J) = \{ p \, : \, |p|<\omega, \, p \text{ is a function, } \dom(P) \subseteq I, \, \ran(P)\subseteq J\}\]
\end{ex}
Let $\leq$ be the revesrse inclusion of $\Fn(I,J)$, so $q \leq p$ iff $q \supseteq p$. Also $\bbone_{\bbP} = \emptyset$. Then $(\Fn(I,J),\leq, \emptyset)$ is a forcing poset, and moreover the pre-order is separative.\gap
Note that when $\alpha$ is an ordinal, $\Fn(\alpha \times \omega, 2)$ is often written $\Add(\omega,\alpha)$.
\newlec
\begin{remark*}
    Note on reals: Let $\bbR$ be your favourite construction of the reals. There are explicit absolute bijections \begin{align*}
        f: \calP(\omega)      & \to {} ^{\omega}\omega \\
        g: {} ^{\omega}\omega & \to {} ^{\omega}2      \\
        h: {} ^{\omega}2      & \to \bbR               \\
    \end{align*}
    So, in $\calM \models \sfZFC$, knowledge of $\calP^{\calM}(\omega)$ gives knowledge of $(^{\omega}\omega)^{\calM}$, $(^{\omega}2)^{\calM}$, $\bbR^{\calM}$.
    So, in order to add a real to $\calM$ it suffices to add a function in $^{\omega}\omega$, or a subset of $\omega$. We can freely switch between these sets. In formal arguments, reals will normally be either subsets of $\omega$ or functions $\omega \to 2$.
\end{remark*}
Back to the forcing poset, $\Fn(I,J)$ (the collection of partial functions from $I$ to $J$).
\begin{fancydefine}{Chains and Antichains}
    Let $\bbP$ be a forcing poset \begin{itemize}
        \item A \textbf{chain} is a subset $C \subseteq \bbP$ such that $\forall p,q \in C (p \leq q \vee q \leq p)$.
        \item A \textbf{antichain} is a subset $A \subseteq \bbP$ such that $\forall p,q \in A. \,\, (p \perp q)$.
        \item An antichain is \textbf{maximal} if it is not a proper subset of any other antichain of $\bbP$.
        \item $\bbP$ has the countable chain condition, (CCC), if every antichain in $\bbP$ is countable.
    \end{itemize}
\end{fancydefine}
\begin{ex*}
    The set of functions $\{ \{(0,0), (1,n)\} \, | \, n \in \omega \}$ forms an antichain of length $\omega$ in $\Fn(2, \omega)$.
\end{ex*}
\begin{fancydefine}{Delta system/Root of a delta system}
    A family of sets $\calA$ forms a \textbf{Delta system with root }$R$ iff $X \cap Y = R$ for all $X \neq Y$ in $\calA$.
\end{fancydefine}
\begin{fancydefine}{Subsets of a set of cardinality $\theta$}
    Let $A$ be a set $\theta$ a cardinal. Then $[A]^{\theta}$ is the set of subsets of $A$ of size $\theta$: \[ [A]^{\theta} = \{ x \subseteq A \, | \, |x|=\theta\}\]. Naturally, $[A]^{<\theta}$ is defined as: \[[A]^{<\theta} = \{ x \subseteq A \, | \, |x|<\theta\}\]
\end{fancydefine}
Recall the theorem that if $\kappa$ is a regular cardinal, and $\calF$ is a family of sets with $|\calF|<\kappa$, then $|X|<\kappa$ for all $X \in \calF$, then $|\cup \calF| < \kappa$.
\begin{fancylem}{Delta System Lemma, $\sfZFC$}
    Let $\kappa$ be an uncountable regular cardinal, let $\calA$ be a family of finite sets with $|\calA| = \kappa$, then there exists $\calB \in [A]^{\kappa}$ such that $B$ forms a delta system.
\end{fancylem}
\proof To begin with, construct $C \in [A]^{\kappa}$ such that all elements of $C$ have the same size. By assumption $|X| < \aleph_{0}$ for all $X \in \calA$. So set $Y_{n} = \{X \in \calA \, | \, |X|=n\}$. If $|Y_{n}| < \kappa$ for each $n < \kappa$, then $|\calA| = |\cup Y_{n}| < \kappa$. This is a contradiction, so some $Y_{n}$ must have size $\kappa$. \gap
Now fix $n \in \omega$ such that $C = Y_{n}$ has size $\kappa$. We proceed inductively on $n$ to prove that if $C = \{X \, : \, |X|=n\}$ has size $\kappa$, then there is a $\calB \subseteq \calC$ of size $\kappa$ such that $\calB$ forms a Delta system.\\

Firstly, if $n=1$, then $C$ must be a family of pairwise disjoint singletons, so forms a Delta system. Now assume $n>1$, the claim holds for $n-1$. For each $p \in \cup C$ let $C_{p} = \{X \in C \, | \, p \in X\}$. Tjere are twp cases: \begin{itemize}
    \item Suppose $|C_{p}| =  \kappa$ for some $p \in \cup C$. Fix such a $p$, and set $\calD = \{ X \setminus \{p\} \, | \, X \in C_{p}\}$. Then $\calD$ has size $\kappa$ and each element of $\calD$ has size $n-1$. So by the inductibe hypothesis, w can find $\calE \in [\calD]^{\kappa}$ such that $\calE$ forms a Delta system with root $R$. Then $\{ Y \cup \{p\} \, | \, Y \in \calE\} \in [C]^{\kappa}$ forms a delta system with root $R \cup \{p\}$.
    \item Suppose $|C_{p}|< \kappa$ for all $p \in \cup C$. Since $\kappa$ is regular, for any set $S$ with $|S|<\kappa$, $\{X \in C \, | \, X \cap S \neq \emptyset \} = \bigcup_{p \in S} C_{p}$ has size less than $\kappa$.\gap
          Then there must exist $X \in C$ such that $X \cap S = \emptyset$. We recursively choose $X_{\alpha} \in C$ for $\alpha < \kappa$ such that \[X_{\alpha} \cap \bigcup_{\beta < \alpha} X_{\beta} = \emptyset.\] This means that $\{X_{\alpha} \, | \, \alpha \in \kappa\} \in [C]^{\kappa}$ is a Delta system with empty root.
\end{itemize}
\endproof
\begin{fancyprop}{}
    If $\kappa$ were either $\omega$ or a singular cardinal, then we \textit{could} find a family of finite sets with $|\calA| = \kappa$ and no $\calB \in [\calA]^{\kappa}$ forms a delta system.
\end{fancyprop}
\proof Exercise. \endproof
\begin{fancylem}{$\sfZFC$}
    $\Fn(I,J)$ has the countable chain condition iff $I = \emptyset$ or $J$ is countable.
\end{fancylem}
\proof Observe that if $I = \emptyset$, or $J= \emptyset$, then $\Fn(I,J)=\emptyset$, so trivially has the CCC.\gap
Now assume that $I,J \neq \emptyset$.
\begin{enumerate}
    \item[($\Rightarrow$)] Suppose $J$ were uncountable. Then fix some $i \in I$, and look at $\{ \{(i,j)\} \, | \, j \in J\}$, this forms an uncountable antichain.
    \item[($\Leftarrow$)] Suppose $J$ was countable, and let $\{p_{\alpha} \, | \, \alpha \in \omega_{1} \}$ be a collection of distinct elements of $\Fn(I,J)$. Let $\calA = \{ \dom(p_{\alpha}) \, | \, \alpha \in \omega_{1}\}.$ Then by the Delta System Lemma, there exists an uncountable set $\calB \in [\calA]^{\omega_{1}}$ with root $R \subseteq I$.\gap Since $R \subseteq \dom(p_{\alpha})$ for all $\dom(p_{\alpha}) \in \calB$, $R$ is finite. Since $J$ is countable, there are only countably many functions from $R$ to $J$. So, since $\calB$ is uncountable, we can find $\alpha \neq \beta$ such that $\dom(p_{\alpha})$ and $\dom(p_{\beta})$ are both in $\calB$, and they agree with each other on $R$ (i.e. $p_{\beta} \upharpoonright R = p_{\alpha} \upharpoonright R$). But then, since $R$ was a root of this system, this means that $\dom(p_{\alpha}) \cap \dom(p_{\beta})$. Therefore, $p_{\alpha} || p_{\beta}$ (since $p_{\alpha} \cup p_{\beta} = R$).
\end{enumerate}
\endproof
\subsection{Dense sets and Genericity}
\begin{fancydefine}{Dense/Open/Filter}
    Let $\bbP$ be a forcing poset. Then \begin{itemize}
        \item $D \subseteq \bbP$ is \textbf{dense} iff $\forall p \in \bbP \, \exists q \in D\,\,(q \leq p).$
        \item $D \subseteq \bbP$ is \textbf{open} iff $\forall p \in D \, \forall q \in \bbP \, (q \leq p \rightarrow q \in D)$.
        \item $\calG \subseteq \bbP$ is a \textbf{filter} on $\bbP$ iff
              \begin{enumerate}
                  \item[(i)] $\bbone_{\bbP} \in \calG$ (non-empty)
                  \item[(ii)] $\forall p,q \in \calG\,\, \exists r \in \calG \,\, (r \leq p \wedge r \leq q)$ (directed)
                  \item[(iii)] $\forall p, q \in \bbP \,\, (q \leq p \wedge q \in \calG \rightarrow p \in \calG)$ (upwards closed)
              \end{enumerate}
    \end{itemize}
\end{fancydefine}

\begin{fancydefine}{$\bbP$-generic over a model}
    Let $\bbP$ be a forcing poset. $\calG$ is said to be $\bbP$\textbf{-generic over }$\calM$ if $G$ is a filter on $\bbP$ and $G \cap D \neq \emptyset$ for every denset $D \subset \bbP$ such that $D \in \calM$.
\end{fancydefine}
\begin{fancylem}{Generic Filter existence}
    Let $\calM$ be an elementary countable set, and $\bbP \in \calM$ be a forcing poset. Then for all $p \in \bbP$, there is a filter $G \subseteq \bbP$ with $p \in G$ such that $\calG$ is $\bbP$-generic over $\calM$.
\end{fancylem}
\proof In our external universe, let $(D_{n} \, | \, n \in \omega)$ enumerate all dense subsets of $\bbP$ which lie in $\calM$. We inductively define $X \subseteq \bbP$, $X = \{q_{n} \, | \, n \in \omega\}$ as follows: \begin{itemize}
    \item $q_{0}=p$
    \item Given $q_{n}$, choose $q_{n+1} \in D_{n}$ such that $q_{n+1} \leq q_{n}$. Let $G = \{ r \in \bbP \, | \, \exists n. \, (q_{n} \leq r)\}$.
\end{itemize}
Then it is clear to see that $G$ is a filter:
\begin{itemize}
    \item Clear to see that $p \leq \bbone$, so $\bbone \in G$.
    \item If $r, r' \in G$, then $q_{n} \leq r$, $q_{m} \leq r'$. Then $q_{\max(m,n)} \leq r,r'$.
    \item If $s \leq r$ and $s \in G$, then $q_{n} \leq s \leq r$ gives $r \in G$. Then $G$ is a generic filter.
\end{itemize}
\newlec
\begin{fancydefine}{Minimal element}
    A $p \in \bbP$ is \textbf{minimal} if $\forall q \in \bbP. \, \neg(q<p)$.
\end{fancydefine}
\begin{fancylem}{}
    Let $\calM$ be a countable model of $\sfZF$ and $\bbP \in \calM$ is a separative partial order. Then either $\bbP$ has a minimal element, or for every $G \subseteq \bbP$ which is $\bbP$-generic over $\calM$, $G \notin \calM$.
\end{fancylem}
\proof Suppose $\bbP$ does not have a minimal element, and that $G \subseteq \bbP$ is generic. Let $\calF \in \calM$ be a filter, $F \subseteq \bbP$, then $D_{\calF} := \bbP \setminus \calF$ is in $\calM$, and is a dense set.\gap
Then $G \cap D_{\calF} \neq \emptyset$ fora ll filters $\calF \in \calM$. So $G \neq \calF$ and thus $G \neq \calM$.
Fix $p \in \bbP$. If $p \notin \calF$, then $p \in D_{\calF}$ Otherwise, suppose $p \in \calF$, since $p$ is not minimal, fix $q< p$. Then $p \nleq q$, so since $\bbP$ is separative, there is $r \leq p$ such that $r \perp q$. But all conditions in $\calF$ are compatible (WHY????), so one of $r$ and $q$ is not in $\calF$.
\endproof
\begin{fancyprop}{}
    For sets $I,J$ with $|I|\geq \omega$ and $|J| \geq 2$. $\Fn(I,J)$ is a separative partial order without a minimal element.
\end{fancyprop}
\begin{fancyprop}{$\sfZFC$}
    Suppose $\bbP \in \calM$ is a forcing poset, and $G \subseteq \bbP$, then TFAE: \begin{enumerate}
        \item[(i)] $\forall D \in \calM. \, (D \text{ dense in }\bbP \rightarrow G \cap D \neq \emptyset)$.
        \item[(ii)] $\forall p \in G\, \forall D \in \calM. \, (D \text{ is dense below }p \cap bbP \rightarrow G \cap D \neq \emptyset)$
        \item[(iii)] $\forall D \in \calM.\, (D \text{ is open dense in } \bbP \rightarrow G \cap D \neq \emptyset).$
        \item[(iv)]$\forall D \in \calM. \, (D \text{ is a maximal antichain in }\bbP \rightarrow G \cap D \neq \emptyset)$.
    \end{enumerate}
\end{fancyprop}
\proof Example Sheet 3 lol. \endproof
\section{Names}
\begin{fancydefine}{$\bbP$-names}
    Let $\bbP$ be a forcing poset, we define the class of $\bbP$\textbf{-names}, $\calM^{\bbP}$ recursively as follows: \begin{enumerate}
        \item[(i)] $\calM_{0}^{\bbP}= \emptyset$.
        \item[(ii)] $\calM_{\alpha + 1}^{\bbP} = \calP^{\calM}(\bbP \times \calM_{\alpha}^{\bbP})$
        \item[(iii)] $\calM_{\lambda}^{\bbP} = \bigcup_{\alpha < \lambda} \calM_{\alpha}^{\bbP}$
        \item[(iv)] $\calM^{\bbP} = \bigcup_{\alpha \in \Ord} M_{\alpha}^{\bbP}$.
    \end{enumerate}
    An element of $\calM^{\bbP}$ is called a $\bbP$\textbf{-name}, denoted by $\dot{x}$. Given $x \in \calM^{\bbP}$, $\ran(\dot{x}) = \{ \dot{y} \, | \, \exists p \in \bbP, \, \langle p, \dot{y}\rangle\}.$
\end{fancydefine}
\begin{remark*}
    Alternatively, by transfinite recursion on rank, the class of $\bbP$-names over $V$, $V^{\bbP}$ is defined in the following way: \gap
    If $\rank(\dot{x}) = \alpha$, then $x \in V^{\bbP}$ iff $\dot{x}$ is a relation and $\forall \langle p, \dot{y} \rangle \in \dot{x}$ $(p \in \bbP, y \in V^{\bbP} \cap V_{\alpha})$. Then $\calM^{\bbP}= V^{\bbP} \cap \calM$.
\end{remark*}
\begin{fancydefine}{$\bbP$-rank}
    The $\bbP$\textbf{-rank}, $\rank_{\bbP}(\dot{x})$ is the least $\alpha$ such that $\dot{x} \subseteq \bbP \times \calM_{\alpha}$.
\end{fancydefine}
\begin{fancydefine}{Interpretation of $\dot{x}$ by $G$}
    Let $\dot{x}$ be a $\bbP$-name, and $G \subseteq \bbP$. We define the \textbf{interpretation of }$\dot{x}$\textbf{ by }$G$ recursively as: \[ \dot{x}^{G} = \{\dot{y}^{G} \, | \, \exists p \in G. \, (\langle p, \dot{y} \rangle \in \dot{x})\}\]
\end{fancydefine}
\begin{fancydefine}{Forcing extension}
    The \textbf{forcing extension} of $\calM$ by $G$, $\calM[G]$ is \[ M[G] = \{ \dot{x}^{G} \, | \, \dot{x} \in \calM^{\bbP}\}.\]
\end{fancydefine}
\begin{ex*}\hspace{0.1cm}
    \begin{itemize}
        \item $\emptyset \in \calM \Rightarrow \emptyset^{G} = \emptyset$.
        \item $\dot{x} = \{ \langle p, \emptyset \rangle, \langle r, \{ \langle q, \emptyset \rangle \} \rangle \}$. If all of $p,q,r \in G$, then \begin{align*}
                  \dot{x}^{G} & = \{ ( \langle p, \emptyset \rangle)^{G}, (\langle r, \{ \langle q, \emptyset \rangle \} \rangle )^{G} \} \\
                              & = \{ \emptyset, \{(\langle q, \emptyset \rangle)\}^{G}\}                                                  \\
                              & = \{ \emptyset, \{ \emptyset \}\}
              \end{align*}
              On the other hand, if $p,r \in G$, then $\dot{x}^{G} = \emptyset$.\\
              If $r \in G, p,q \notin G$, then $\dot{x}^{G} = \{(\langle q, \emptyset \rangle)^{G} \} = \{\emptyset\}$. \\And finally, if $p \in G$ and $r \notin G$, then $\dot{x}^{G} = \{\emptyset\}$.
    \end{itemize}
\end{ex*}
\section{Generic Model Theorem}
\begin{fancythm}{Generic Model Theorem}
    Let $\calM$ be a countable transitive model of $\sfZF$. Let $\bbP$ be a forcing poset, and $G$ be a $\bbP$-generic filter. Then \begin{enumerate}
        \item[(i)] $\calM[G]$ is a transitive set.
        \item[(ii)] $|\calM[G]| = \aleph_{0}$.
        \item[(iii)] $\calM[G] \models \sfZFC$, and if $\calM \models \sfAC$ then $M[G] \models \sfAC$.
        \item[(iv)] $\Ord^{\calM} = \Ord^{\calM[G]}$
        \item[(v)] $\calM \subseteq \calM[G]$
        \item[(vi)] $\calM[G]$ is the smallest countable transitive model of $\sfZF$ such that $\calM \subseteq \calM[G]$ and $G \in \calM[G]$.
    \end{enumerate}
\end{fancythm}
This is a big theorem, that we will prove over multiple lectures, and in multipl small parts.
\subsection{Canonical names}
\begin{fancydefine}{Canonical name}
    Given $(\bbP, \leq, \bbone)$ ad set $x \in \calM$, we recursively define the \textbf{canonical name} of $x$, $\check{x}$ as \[ \check{x} = \{ \langle \bbone, \check{y} \rangle \, | \, y \in x \}.\] (pronounced $x$-check).
\end{fancydefine}
\begin{fancylem}{}
    If $\calM$ is a transitive model of $\sfZF$, $\bbP \in \calM$ and $\bbone \in G \subseteq \bbP$> Then \begin{enumerate}
        \item[(i)] $\forall x \in \calM (\check{x} \in \calM^{\bbP} \text{ and } \check{x}^{G} = x)$.
        \item[(ii)] $\calM \subseteq M[G]$.
        \item[(iii)] $\calM[G]$ is transitive.
    \end{enumerate}
\end{fancylem}
\proof (ii) follows from (i). (i) is proven inductively. We prove that $\check{x} \in \calM^{\bbP}$ by using the definition of $\bbP$-names by transfinite recursion. Next $\check{x}^{G} = \{ \check{y}^{G} \, | \, y \in x\} = x$. \gap
To prove that $\calM[G]$ is transitive, suppose that $x \in y$ and $y \in \calM[G]$. Then by definition, $y = \dot{y}^{G}$ for some $\dot{y} \in \calM^{\bbP}$. By construction, any element of $y$ is of the form $\dot{z}^{G}$. So in particular, we must have that $x = \dot{x}^{G}$ for some $\dot{x} \in \calM^{\bbP}$. So $\dot{x}^{G} \in \calM[G]$. \endproof
\begin{remark*}
    Even if $G \notin \calM$, we can still define a name for $G$ in $\calM$. From this, it follows that if $G \notin \calM$, then $\calM[G]\neq \calM$.
\end{remark*}
\begin{fancyprop}
    Let $\dot{G} = \{ \langle p, \check{p} \rangle \, | \, p \in \bbP\}$. Then $G = \dot{G}^{\bbP} \in M[G]$.
\end{fancyprop}
\proof
\begin{align*}
    \dot{G}^{G} & = \{ \check{p}^{G} \, | \, p \in G \} \\
                & = \{ p \, | \, p \in G \}             \\
    ^= G
\end{align*}
\endproof
\begin{fancydefine}{Unordered pairs/ordered pairs}
    Given $\dot{x},\dot{y} \in \calM^{\bbP}$, let $\up(\dot{x},\dot{y})=\{\langle \bbone, \dot{x} \rangle, \langle \bbone, \dot{y} \rangle \}$, and let $\ordp(\dot{x},\dot{y})= \up(\up(\dot{x},\dot{x}),\up{\dot{x},\dot{y}})$.
    Very similar to how we define kuratowski pairs.
\end{fancydefine}
\begin{fancyprop}{}
    For any $\dot{x}, \dot{y} \in \calM^{\bbP}$, and $G \subseteq \bbP$, $\bbone \in G$: \begin{align*}
        (\up(\dot{x},\dot{y}))^{G} & = \{ \dot{x}^{G}, \dot{y}^{G}\} \\
        (\ordp(\dot{x},\dot{y}))^{G} = \langle \dot{x}^{G}, \dot{y}^{G}\rangle
    \end{align*}
\end{fancyprop}
\proof Omitted \endproof
\subsection{First results towards showing Generic Model Theorem}
\begin{fancylem}{}
    Suppose $\calM$ is a transitive model of $\sfZF$, $\bbP \in \calM$ is a forcing poset. Then if $G \subseteq \bbP$ and $\bbone \in G$, then $\calM[G]$ is a transitive model of \begin{itemize}
        \item Extensionality
        \item Empty set
        \item Foundation
        \item Pairing
    \end{itemize}
\end{fancylem}
\newlec
\begin{fancylem}
    Suppose $\calM$ is a transitive model of $\sfZF$, $\bbP \in \calM$ is a forcing poset. Then if $G \subseteq \bbP$ and $\bbone \in G$, we have \begin{enumerate}
        \item[(i)] $\rank(x^{G}) \leq \rank(x)$ for all $x \in \calM^{\bbP}$.
        \item[(ii)] $\Ord^{\calM} = \Ord^{\calM[G]}$.
        \item[(iii)] $|\calM[G]| = |\calM|$.
    \end{enumerate}
\end{fancylem}
\proof \hspace{0.1cm} \begin{enumerate}
    \item[(i)] We prove this by induction: \begin{itemize}
              \item $\emptyset^{G}=\emptyset$ and $\rank(\emptyset^{G})=\rank(\emptyset)=0$.
              \item For the inductive step, we have: \begin{align*}
                        \rank(\dot{x}^{G}) & =\sup\{ \rank(u) + 1 \, | \, u \in \dot{x}^{G}\}                         \\
                                           & \leq \sup\{\rank(\dot{y}^{G}) + 1 \, | \, \dot{y} \in \ran(\dot{x})\}    \\
                                           & \leq \sup\{\rank(\dot{y})+1 \, | \, \dot{y} \in \ran(\dot{x})\} \tag{IH} \\
                                           & \leq \sup\{\rank(u)+1 \, | \, u \in \dot{x}\}                            \\
                                           & \leq \rank(\dot{x}).
                    \end{align*}
          \end{itemize}
    \item[(ii)] Since $\calM \subseteq \calM[G]$, and also being an ordinal is an absolute property, $\Ord^{\calM} \subseteq \Ord^{\calM[G]}$. For the reverse inclusion, take $\alpha \in \Ord^{\calM[G]}$ and fix $\dot{x} \in \calM^{\bbP}$ such that $\alpha = \dot{x}^{G}$. Then $\alpha = \rank(\alpha) \leq \rank(\dot{x})$. So, since $\calM$ is transitive, $\alpha \in \Ord^{\calM}$.
    \item[(iii)] Since any element of $\calM[G]$ is of the form $\dot{x}^{G}$ for some $\dot{x} \in \calM^{\bbP} \subseteq \calM \subseteq \calM[G]$. Therefore, \[ |\calM[G]| \leq |\calM^{\bbP}| \leq |\calM| \leq |\calM[G]|.\]
          This follows, since $\dot{x} = (\check{\dot{x}})^{G}$
\end{enumerate}
\endproof



\begin{corr}{}
    $M[G]$ models the axiom of infinity.
\end{corr}
\proof Omitted \endproof
\begin{fancylem}
    Suppose $\calM$ is a transitive model of $\sfZF$, $\bbP \in \calM$, $G \subseteq \bbP$, $\bbone \in G$. Then if $\calN$ is any other transitive model of $\sfZF$ containing $\calM$ as a definable class in $\calN$, and $G \in \calN$, then $M[G] \subseteq \calN$.
\end{fancylem}
\proof If $\calN \supseteq \calM$ is a transitive model of $\sfZF$ with $G \in \calN$, we carry out the construction of $\calM[G]$ in $\calN$. Namely, show that for all $\dot{x} \in \calM^{\bbP}$, $\dot{x}^{G} \in \calN$ by induction on ranks. \begin{itemize}
    \item First, sicne the empty set axiom holds in $\calN$, $\emptyset^{G} = \emptyset \in \calN$.\gap
          Moreover, since $\calM^{\bbP} = V^{\bbP} \cap \calM \subseteq \calV^{\bbP} \cap \calN = \calN^{\bbP}$. So if $\dot{x} \in \calM^{\bbP}$, then $\dot{x} \in \calN^{\bbP}$, and in particular $\dot{x} \in \calN$.
    \item For the inductive step, suppose that for every $\langle p, \dot{y}\rangle \in \dot{x}$, $\dot{y}^{G} \in \calN$:
          \begin{align*}
              \left( (\dot{x})^{G} \right)^{\calN} & = \left( \{\dot{y}^{G} \, | \, \exists p \in G. \, (\langle p, \dot{y} \rangle \in \dot{x}) \}\right)^{\calN}      \\
                                                   & = \left(\{ (\dot{y}^{g})^{\calN} \, | \, (\exists p \in G \langle p, \dot{y} \rangle \in \dot{x})\}\right)^{\calN} \\
                                                   & = \{ \dot{y}^{G} \, | \, \exists p \in G. \, \langle p, \dot{y} \rangle \in \dot{x} \}                             \\
                                                   & = \dot{x}^{G}
          \end{align*}
\end{itemize}
\endproof
\begin{fancylem}{}
    Suppose $\calM$, $\bbP$, $G$ as above. Additionally, assume that $G$ is a filter. Then $\calM[G]$ satisfies unions.
\end{fancylem}
\proof It suffices to prove that for all $a \in \calM[G]$, there is some $b \in \calM[G]$ such that $\cup a = b$. In order to show this, fix $\dot{a} \in \calM^{\bbP}$ such that $\dot{a}^{G} = a$ and let $\dot{b} = \{ \langle p, \dot{z} \rangle \, | \, \exists \langle q, \dot{y} \rangle \in \dot{a},$$\, \exists r \in \bbP (\langle r, \dot{z} \rangle \in \dot{y} \wedge p \leq r,q)\}$. \gap
    Observe that $\dot{b} \in \calM^{\bbP}$. Since $\dot{a}$ is a $\bbP$ name, any $adot{y} \in \ran(a)$ is a $\bbP$-name. Then $\dot{b}$ consists of pairs $\langle p, \dot{x}\rangle$ where $p \in \bbP$, $\dot{z} \in \ran{\dot{y}}$ for some $y \in \ran(\dot{a})$ Thus $\dot{z} \in V^{\bbP}$. Moreover, $b \in \calM$ since $dot{b} \subseteq \bbP \times \tcl(\dot{a})$.\gap
    Now we need to show that $\cup a \subseteq \dot{b}^{G}$. Take $w \in \cup a$. Then $w \in v$ for some $v \in a$. Since $\calM[G]$ is transitive, $w,v \in \calM[G]$. So for $\dot{y}, \dot{z} \in \calM^{\bbP}$. and conditions $q,r \in G$ such that $w= \dot{z}^{G}$, $v = \dot{y}^{G}$, $\langle q, \dot{y} \rangle \in \dot{a}$, $\langle r, \dot{x} \rangle \in \dot{y}$. \gap
    Now, since $G$ is a filter, by downwards directedness, fix $ p \leq q,r$, $p \in G$. But then $\langle p, \dot{x} \rangle \in \dot{b}$ and $w = \dot{z}^{G} \in \dot{b}^{G}$. \gap
    Now we want to show that $\dot{b}^{G} \subseteq \cup a$. Take $c \in \dot{b}^{G}$ and fix $\langle p, \dot{x} \rangle \in \dot{b}$ such that $p \in G$ and $\dot{z}^{G} = c$. \gap
    By definition, fix $\langle q, \dot{y} \rangle \in \dot{a}$ and $r \in \bbP$ such that $\langle r, \dot{z} \rangle \in \dot{y}$, and $p \leq q,r$. Since $G$ is a filter, it is upwards closed, so $q,r \in G$. Therefore $\dot{z}^{G} \in \dot{y}^{G}$, $\dot{y}^{G} \in \dot{a}^{G}$. So $c \in \dot{y}^{G}$ for some $\dot{y}^{G} \in a$.
    \endproof
    \subsection{Moivation for Genericity of Generic Model Theorem}
    Suppose $\bbP \in \calM$, $\calM \subseteq \calM[G]$. Then $\bbP, G \in \calM[G]$. If $\calM[G]$ models anything reasonable, $\bbP \setminus G \in \calM[G]$. So if we try to build a name for $\bbP \setminus G$, a natural name is $c = \{ \langle q, \check{p} \rangle  \, | \, p,q \in \bbP, \, p \perp q\}$. Then $\dot{c}^{G} = \{ p \, | \, \exists q \in G, \, p \perp q\}$. Now if $G$ is a filter, its elements are pairwise compatible, so $G \cap \dot{c}^{G} = \emptyset$. But we still need $G \cup \dot{c}^{G} = \bbP$.\gap
    Now, for each $p \in \bbP$, set $D_{p} = \{ q \in \bbP \, | \, p \perp q \vee q \leq p\}$. It is easy to check that $D_{p} \in \calM$ is dense. Now if $G$ were generic, we could fix $q \in G \cap D_{p}$. Then if $p \perp q$, $p \in \dot{c}^{G}$, and if $q \leq p$, then $p \in G$. Thus $G \cup \dot{c}^{G}=\bbP$.
    \begin{fancyprop}{}
        Suppose $\calM$ is a countable transitive model of $\sfZF$. Then there exists a $\bbP \in \calM$ and a (non-generic) filter $G \subseteq \bbP$ such that $\bbP \setminus G \notin \calM[G]$.
    \end{fancyprop}
    \proof Discussed in example sheet 3. \endproof
    \subsection{Forcing relation}
    To show separation, suppose $\varphi(x,y)$ were a formula, $a,b \in \calM^{\bbP}$, need to show that $C = \{ z \in \dot{a}^{G} \, | \, (\varphi(z,\dot{b}^{G}))^{\calM[G]}\} \in \calM[G]$. This is unclear even for $\varphi(x,y) \equiv x \notin y$. We will build a way to reason about when $\varphi$ holds from within $\calM$ without having to rely on $G$. \gap
    To do this, define a relation between conditions and names in $V^{\bbP}$, written $p \forces \varphi$. Its relativisation $(p \forces \varphi)^{\calM}$ will give us a way to work in $\calM$.\gap
    Our aim is to define $\forces$ such that: $ p \forces \varphi(\dot{u}) \Leftrightarrow \forall G \subseteq \bbP$ with $G$ generic, and $p \in G$, $\calM[G] \models \varphi(\dot{u}^{G})$.
    A naive definition might be to define $\langle p, \dot{x} \rangle \in \dot{y} \Rightarrow p \forces \dot{x} \in \dot{y}$.
    Why not $\Leftarrow$. Consider $\dot{x} = \{\langle p ,\varphi
\rangle \}$ where $p \neq \bbone$, Then $p \forces \emptyset \in \dot{x}$. Suppose $q \perp p$, then we will have $q \forces \dot{x} = \emptyset$, and $q \forces \dot{x} \in \check{1}$. But $\langle q, \dot{x} \rangle \notin \check{1} = \{ \langle \bbone, \emptyset \rangle \}$. (Note that generics meet dense sets, so suffices to consider dense sets).
    \begin{fancydefine}{$\bbP$-forcing language}
        For a forcing poset $\bbP$, the $\bbP$\textbf{-forcing language} $\calF \calL_{\bbP}$ is the class of logical formulas using the binary relation $\in$, and constant symbols from $V^{\bbP}$.
    \end{fancydefine}
    Recall that $D$ is dense below $p$ iff $\forall q \leq p \, \exists s \in D. \, s \leq q$. T
    \begin{fancydefine}{Forcing relation}
        Let $p \in \bbP$, and $\dot{x},\dot{y},\dot{u} \in V^{\bbP}$, we define the \textbf{forcing relation}, $p \forces \varphi(\dot{u})$ recursively as follows: \begin{itemize}
            \item $p \forces \varphi(\dot{u}) \wedge \psi(\dot{u})$ iff $p \forces \varphi(\dot{u})$ and $p \forces \psi(\dot{u})$.
            \item $p \forces \neg \varphi(\dot{u})$ iff tyhere is no $q \leq p$ such that $q \forces \varphi(\dot{u})$.
            \item $p \forces \exists x \varphi(x, \dot{u})$ iff the set $\{q \leq p \, | \, \exists \dot{x} \in V^{\bbP}, \, q \forces \varphi(\dot{x},\dot{u})\}$ is dense below $p$.
            \item $p \forces \dot{x} \subseteq \dot{y}$ iff for all $\langle q_{1}, \dot{z}_{1} \rangle \in \dot{x}$, $\{r \leq p \, | \, r \leq q_{1} \rightarrow \exists \langle q_{2}, \dot{z}_{2} \rangle \in \dot{y}. \, r \leq q_{2} \wedge r \forces \dot{z}_{1}=\dot{z}_{2}\}$ This is what was written but probably still check notes online for this section.
        \end{itemize}
        The idea of this inclusion definition is that for everything in the range of $\dot{x}$, there is something in the range of $\dot{y}$, and some $r \in G$ such that $r \forces \dot{z}_{1} = \dot{z}_{2}$, and $\langle q_{1}, \dot{z}_{1} \rangle \in \dot{x}$, $\langle q_{2}, \dot{z}_{2} \rangle \in \dot{y}$. So we want the $r$ to satisfy $r \leq p, q_{1}, q_{2}$.
        \begin{itemize}
            \item $p \forces \dot{x} \in \dot{y}$ iff $ \{ q \leq p \, | \, \exists \langle r, \dot{z} \rangle \in \dot{y} (q \leq r \wedge q \forces \dot{x} = \dot{z})\}$ is dense below $p$.
            \item $p \forces \varphi \wedge \psi$
            \item Note that $p \forces \dot{x} = \dot{y}$ is the same as $p \forces \dot{x} \subseteq \dot{y}$ and $p \forces \dot{y} \subseteq \dot{x}$.
        \end{itemize}
    \end{fancydefine}
    \newlec
    \begin{remark*}\hspace{0.1cm}
        \begin{itemize}
            \item Formally, $p \forces \dot{x} \subseteq \dot{y}$ and $p \forces \dot{x} \in \dot{y}$ are circular definitions, and thus are we define them recursively.
            \item All the clauses, except for the existential one, use only absolute notions. Being dense below a condition is absolute.
            \item When relativising to $\calM$, the key difference for the definition of forcing an existential is $(p \forces \exists x \varphi(x))^{\calM}$, the $\exists \dot{x} \in V^{\bbP}$ becomes $\exists \dot{x} \in \calM^{\bbP}$.
        \end{itemize}
    \end{remark*}
    \begin{fancyprop}{}
        For $p \in \bbP$, $\varphi \in \calF \calL_{\bbP}$, and $\dot{x}_{1}, \dots, \dot{x}_{n} \in V^{\bbP}$, the following are equivalent: \begin{enumerate}
            \item[(i)] $p \forces \varphi(\dot{x}_{1}, \dots, \dot{x}_{n})$.
            \item[(ii)] $\forall q \leq p. \, q \forces \varphi(\dot{x}_{1}, \dots, \dot{x}_{n})$.
            \item[(iii)] $\neq \exists q \leq p. \, q \forces \neg \varphi(\dot{x}_{1}, \dots, \dot{x}_{n})$
            \item[(iv)] $\{ r \, | \, r \forces \varphi(\dot{x}_{1}, \dots, \dot{x}_{n}) \}$ is dense below $p$.
        \end{enumerate}
    \end{fancyprop}
    \proof \hspace{0.1cm}
    \begin{enumerate}
        \item[(ii) $\Rightarrow$ (iii)] Suppose $\forall q \leq p$, $q \forces \varphi$, and $\exists q \leq p$, $q \forces \neg \varphi$. Then there is no $r \leq q$ such that $r \forces \varphi$. So $q \nVdash \varphi$.
        \item[(iii) $\Rightarrow$ (iv)] Suppose $\neg \exists q \leq p$, $p \forces \neg \varphi$. Take $q \leq p$, then $q \nVdash \neg p$. So fix $r \leq q$ such that $r \forces \varphi$. So the set is dense below $p$.
        \item[(i) $\Rightarrow$ (ii)] By induction on formula complexity: \gap \begin{itemize}
                  \item For atomic formulas, let $\square \in \{ \in, \subseteq \}$. Then $p \forces \dot{x} \square \dot{y}$ iff some set $A$ (from the definition of the forcing relation) is dense below $p$. Take $q \leq p$. Then $A$ is dense below $q$, since if $s \leq q \leq p$, there is some $r \leq s$ such that $r \in A$. Thus $q \forces \dot{x} \square \dot{y}$.
                  \item Suppose the claim holds for $\varphi$ and $\psi$. \gap
                        If $q \leq p$, then $p \forces \neg \varphi$ says $\neg \exists r \leq p. \, (r \forces \varphi)$. Therefore $\neg \exists r \leq q (r \forces \varphi)$. So $q \forces \neg \varphi$.\gap
                        If $p \forces \varphi \wedge \psi$, then $p \forces \varphi$ and $p \forces \psi$. So by our inductive hypothesis, $q \forces \phi$ and $q \forces \psi$. Therefore $q \forces \phi \wedge \psi$.\gap
                        If $p \forces \exists x. \, \varphi(x)$, then this says that some set $A$ is dense below $p$. This implies that $A$ is dense below $q$ (by our inductive hypothesis), so $q \forces \exists x. \, \varphi(x)$.
              \end{itemize}
        \item[(iv) $\Rightarrow$ (i)] As with the previous case, we proceed by induction on formula complexity. \begin{itemize}
                  \item We will first do atomic formulas. Let $\square \in \{\in, \subseteq\}$ again. To prove that $p \forces \dot{x} \square \dot{y}$, we need to show that some set $A$ is dense below $p$. But now, we have that $\{ r \, | \, r \forces \dot{x} \square \dot{y}\}$ is dense below $p$ (which it is, by assumption). Fix some $q \leq p$, then there is some $r \leq q$ such that $r \forces \dot{x} \square \dot{y}$. So there is some $s \leq r \leq q \leq p$ such that $s \in A$. Thus $A$ is dense below $p$.
                  \item The proof for $p \forces \exists x \varphi(x)$ is the same.\gap
                        Next, suppose that $\{r \, | \, r \forces \varphi \wedge \psi\}$ is dense below $p$. Since $r \forces \varphi \wedge \psi$ iff $r \forces \psi$ and $r \forces \varphi$. By the inductive hypothesis, $p \forces \varphi$ and $p \forces \psi$.\gap
                        Finally, suppose $\{r \, | \, r \forces \neg \varphi\}$ is dense below $p$. Fix $q \leq p$, and suppose for a contradiction that $q \forces \varphi$. Then by (i) $\Rightarrow$ (iii), there is no $r \leq q$ such that $r \forces \neg \varphi$, this contradicts that $\{r \, | \, r \forces \neg \varphi\}$ is dense below $p$.
              \end{itemize}
    \end{enumerate}
    \endproof
    \begin{fancyprop}{}
        For any $\bbP$, $p,q \in \bbP$, and name $\dot{a}, \dot{b} \in V^{\bbP}$: \begin{enumerate}
            \item[(i)] $p \forces \dot{a}=\dot{a}$
            \item[(ii)]$\langle q, \dot{b}\rangle \in \dot{a}$ and $p \leq q$ then $p \forces \dot{b} \in \dot{a}$.
            \item[(iii)] If $\calM$ is a transitive model of $\sfZF$ and $\bbP \in \calM$, then for any $\varphi$, $\psi$, we have $\{ \langle q, \dot{x} \rangle \, | \, \langle q, \dot{x} \rangle \in \dot{a} \wedge (q \forces \varphi(\dot{x}))^{\calM} \} \in \calM$. Also, we have that $\{q \, | \, q \in \bbP \wedge q \forces (\psi(\dot{a}))^{\check{\calM}} \} \in \calM$. So there are two ways of relativising, and they both work.
            \item[(iv)] $p \forces \varphi \vee \psi$ iff $\{ q\leq p \, | \, q \forces \varphi \text{ or } q \forces \psi\}$.
            \item[(v)] $p \forces \varphi \rightarrow \psi$ iff for all $\dot{x} \in V^{\bbP}$, $p \forces \varphi(\dot{x})$.
            \item[(vi)] $p \forces \forall \dot{x} \varphi(\dot{x})$ iff for all $\dot{x} \in V^{\bbP}$, $p \forces \varphi(\dot{x})$.
            \item[(vii)] For every $\varphi$, $\{ p \in \bbP \, | \, p \forces \varphi \vee p \forces \neg \varphi\}$ is a dense open set.
            \item[(viii)] There is no $p$ and formula $\varphi$ such that $p \forces \varphi$ and $p \forces \neg \varphi$.
        \end{enumerate}
    \end{fancyprop}
    \proof Example sheet 4 (probably). \endproof
    \newlec
    \begin{fancythm}{Forcing Theorem}
        Suppose $\calM$ is a transitive model of $\sfZF$, $\bbP \in \calM$, $\varphi(u)$ is a formula, and $G$ is $\bbP$-generic over $\calM$. Then, for any $\dot{x} \in \calM^{\bbP}$: \begin{enumerate}
            \item[(i)] If $p \in G$ and $(p \forces \varphi(\dot{x}))^{\calM}$, then $\calM[G] \models \varphi(\dot{x}^{G})$.
            \item[(ii)] If $\calM[G] \models \varphi(\dot{x}^{G})$ then $\exists p \in G. \, (p \forces \varphi(\dot{x}))^{\calM}$.
        \end{enumerate}
    \end{fancythm}
    \proof
    Proceed by induction on the complexity of formulas. Note that we need to work with $(p \forces \varphi(\bar{v}))^{\calM}$, i.e. everything is relatizied, but since relativization and parameters are only important for statements with existential quantifiers, we will suppress them if there aren't any existential quantifiers.

    Let $\Psi(\varphi)$ be the claim that for any $\dot{x} \in \calM^{\bbP}$: \begin{enumerate}
        \item[(i)] If $p \in G$ and $(p \forces p(x))^{\calM}$, then $\calM[G] \models \varphi(\dot{x}^{G})$.
        \item[(ii)] If $M[G] \models \varphi(\dot{x}^{G})$ then $\exists p \in G (p \forces \varphi(x))^{\calM}$.
    \end{enumerate}
    \begin{itemize}
        \item So $\Psi(\varphi) \Rightarrow \Psi(\neg \varphi)$.
              \begin{enumerate}
                  \item[(i)] Suppose $p \in G$, $p \forces \neg \varphi$. Suppose for a contradiction $\calM[G] \models \varphi$, i.e. $\varphi^{\calM[G]}$ holds. Then by $\Psi(\varphi)$, fix $q \in G$ such that $q \forces \varphi$. Since $G$ is a filter, fix $r \leq p,q$. So $r \forces \varphi$, contradicting $p \forces \neg \varphi$. Thus $\neg(\varphi^{\calM[G]})$. So $(\neg \varphi)^{\calM[G]}$. That is, $M[G] \models \neg \varphi$.
                  \item[(ii)] Suppose that $M[G] \models \neg \varphi$, and let $D = \{ p \in \bbP \, : \, p \forces \varphi \vee p \forces \neg \varphi\}$. $D$ is dense, since if $q \nVdash \varphi$, then by definition there s some $p \leq q$ such that $p \forces \neg \varphi$. Then $p \in D$. So fix $p \in G \cap D$. If $p \forces \varphi$ then, by $\Psi(\varphi)$, $M[G]\models \varphi$. So $p \forces \neg \varphi$.
              \end{enumerate}
        \item $\Psi(\varphi) \wedge \Psi(\psi) \Rightarrow \Psi(\varphi \wedge \psi)$.
              \begin{enumerate}
                  \item[(i)] Suppose $p \forces \varphi \wedge \psi$. Then $p \forces \varphi$ and $p \forces \psi$. Also $p \in G$. Since $\Psi(\varphi)$ and $\Psi(\psi)$ hold, $M[G] \models \varphi$ and $M[G] \models \psi$. So $M[G] \models \varphi \wedge \psi$.
                  \item[(ii)] Suppose $M[G] \models \varphi \wedge \psi$. Then $M[G] \models \varphi$ and $M[G] \models \psi$. Since $\Psi(\varphi)$ and $\Psi(\psi)$ hold, fix $p_{1},p_{2} \in G$ such that $p_{1} \forces \varphi$ and $p_{2} \forces \psi$. Since $G$ is a filter, fix $r \leq p_{1}, p_{2}$. Then $r \forces \varphi \wedge \psi$.
              \end{enumerate}
        \item $\Psi(\varphi(\dot{x})) \Rightarrow \Psi(\exists x \varphi(x)).$
              \begin{enumerate}
                  \item Suppose $(p \forces \exists x \varphi(x))^{\calM}$, $p \in G$. Let \begin{align*}
                            D & = (\{ q \leq p \, | \, \exists \dot{x} \in V^{\bbP} ( q \forces \varphi(\dot{x}))\})^{\calM}.      \\
                              & = \{ q \leq p \, | \, \exists \dot{x} \in \calM^{\bbP} (q \forces \varphi(x))^{\calM} \} \in \calM
                        \end{align*}
                        By definition, $D$ is dense. So since $G$ is generic, fix $q \in G \cap D$. Then fix $\dot{x} \in \calM^{\bbP}$ such that $(q \forces \varphi(\dot{x}))^{\calM}$. So since $\Psi(\varphi(\dot{x}))$, $M[G] \models \varphi(\dot{x}^{G})$. Therefore $M[G] \models \exists x. \, \varphi(x)$.
                  \item Suppose $M[G] \models \exists x. \, \varphi(x)$. Fix $\dot{x} \in \calM^{\bbP}$. such that $M[G] \models \varphi(\dot{x}^{G})$. Since $\Psi(\varphi(x))$, there eixsts $p \in G$ such that $(p \forces \varphi(\dot{x}))$. Thus $\{ q \leq p \, | \, (q \forces \varphi(\dot{x}))^{\calM}\}$ is dense. Then, by definition, $(p \forces \exists x. \, \varphi(x))^{\calM}$.
              \end{enumerate}
        \item $\Psi(x=y)$ We will prove next lecture, and we will assume it for now.
        \item $\Psi(x \in y)$ \begin{enumerate}
                  \item[(i)] Suppose $p \forces \dot{x} \in \dot{y}$, $p \in G$. Let $D = \{ q \leq p \, | \, \exists \langle r, \dot{z}\rangle \in \dot{y} (q \leq r \wedge q \forces \dot{x} = \dot{z})\}$. By definition, $D$ is dense. Fix $q \in G \cap D$. Since $q \in D$, fix $\langle r, \dot{z} \rangle \in \dot{y}$ such that $q \leq r$ and $q \forces \dot{x} = \dot{z}$. Since $q \in G$, assuming that $\Psi(x=y)$, $\calM[G] \models \dot{x}^{G} = \dot{z}^{G}$. Since $G$ is a filter, $q \leq r$, $r \in G$, and so $\dot{z}^{G} \in \dot{y}$. So $M[G] \models \dot{x}^{G} \in \dot{y}^{G}$.
                  \item[(ii)] Suppose $M[G] \models \dot{x}^{G} \in \dot{y}^{G}$. Then fix $\langle r, \dot{z} \rangle \in \dot{y}$ such that $r \in G$ and $\dot{z}^{G} = \dot{x}^{G}$. Now by $\Psi(x=y)$, there is some $q \in G$ such that $q \forces \dot{x} = \dot{z}$. Since $G$ is a filter, fix $p \in G$ such that $p \leq q,r$. Finally, $p \forces \dot{z} \in \dot{y}$ and $p \forces \dot{x}=\dot{z}$, so for all $s \leq p$, $s \leq r$, and $s \forces \dot{x} = \dot{z}$. This gives that $D$ is dense below $p$. hence $p \forces \dot{x} \in \dot{y}$.
              \end{enumerate}
    \end{itemize}
    \endproof

    \begin{corr}{}
        Suppose $\calM$ is a countable transitive model of $\sfZF$, $\bbP \in \calM$, and $\varphi(u)$ is a formula. Then for any name $\dot{x}$: \[ (p \forces \varphi(\dot{x}))^{\calM} \Leftrightarrow \text{for any $\bbP$-generic filter $G$ over $M$ with $p \in G$, $\calM[G] \models \varphi(\dot{x}^{G})$}\]
    \end{corr}
    \proof \hspace{0.1cm}
    \begin{enumerate}
        \item[($\Rightarrow$)] This is just clause (i) of the Forcing Theorem.
        \item[($\Leftarrow$)] Suppose $(p \nVdash \varphi(\dot{x}))^{\calM}$. Then there is some $q \leq p$ such that $(q \forces \neg \varphi(\dot{x}))^{\calM}$. Now, let $G$ be a $\bbP$-generic filter over $\calM$ such that $q \in G$. Then, since $G$ is upwards closed, $p \in G$. However, $M[G] \models \varphi$, but since $q \in G$, then by the Forcing Theorem, we get $M[G] \models \neg \varphi$. This leads to a contradiction, since then there is some $p$ which forces $\varphi$ and $\neg \varphi$ which contradicts part (viii) of proposition 15.16.
    \end{enumerate}
    \endproof
    \begin{fancydefine}{Star forcing, $\forces^{*}$}
        Suppose $\calM$ is a countable transitive model of $\sfZF$, $\bbP \in \calM$, $\dot{x}_{1}, \dots \dot{x}_{n} \in \calM^{\bbP}$, $p \in \bbP$, $\varphi(v_{0},\dots, v_{n})$ a formula. Then define $\forces_{\bbP,\calM}^{*}$ by $p \forces_{\bbP,\calM}^{*} \varphi(\dot{x}_{1},\dots,\dot{x}_{n})$ iff $M[G] \models \varphi(\dot{x}_{1}^{G},\dots,\dot{x}_{n}^{G})$ for all $G \subseteq \bbP$ such that $p \in G$ and $G$ is a $\bbP$-generic filter.
    \end{fancydefine}
    \newlec
    \begin{fancylem}{}
        If $\calM$ is a transitive model of $\sfZF$, $\bbP \in \calM$, $G \subseteq \bbP$ is generic, then $\calM[G] \models \text{Sep}$.
    \end{fancylem}
    \proof Let $\varphi(x,v)$ be a formula, $\Fr(\varphi) \subseteq \{x,v\}$. Then it suffices to prove that for any $a,v \in \sfM[G]$, $b = \{ x \in a \, | \, \calM[G] \models \varphi(x,v)\} \in \calM[G]$. Fix names $\dot{a}, \dot{v} \in \calM^{\bbP}$ such that $\dot{a}^{G}=a$ and $\dot{v}^{G}=v$. \gap
    Any element of $\dot{a}^{G}$ is of the form $\dot{x}^{G}$ when $\langle p, \dot{x} \rangle \in \dot{a}$ and $p \in G$. Thus \[b = \{ \dot{x}^{G} \, | \, \exists p. \, (\langle p, \dot{x} \rangle \in \dot{a} \wedge p \in G \wedge M[G] \models \varphi(\dot{x}^{G},\dot{a}^{G})\}\].
    Let $\dot{b} = \{ \langle p, \dot{x} \rangle \, | \, \langle p, \dot{x} \rangle \in dot{a} \wedge (p \forces \varphi(\dot{x}, \dot{v}))^{\calM}\} \in \calM$. Thus $\dot{b}^{G} \in \calM[G]$. \gap
    So \begin{align*}
        x \in \dot{b}^{G} \Leftrightarrow & \text{there is some $\bbP$-name $\dot{x} \in \calM$ such that $\dot{x}^{G}=x$, $\langle p, \dot{x} \rangle \in \dot{a}$, $(p \forces \varphi(\dot{x},\dot{v}))^{\calM}$.} \\
        \Leftrightarrow                   & x \in \dot{a}^{G} \text{ and } M[G] \models \varphi(x,v)                                                                                                                  \\
        \Leftrightarrow                   & x \in b
    \end{align*}
    \endproof
    \begin{fancylem}{}
        $\calM$ is a transitive model of $\sfZF$, $\bbP \in \calM$, $G \subseteq \bbP$ is generic. Then $M[G] \models \text{ collection}$. It's easier to show that it models collection than replacement.
    \end{fancylem}
    Let $\varphi(x,y,v)$ be a formula with $\Fr(\varphi) \subseteq \{x,y,v\}$. Fox $a,v \in \calM[G]$ with names $\dot{a}, \dot{v} \in \calM^{\bbP}$.\gap
    Suppose that $\calM[G] \models \forall x \in a \, \exists y, \, \varphi(x,y,v)$. \gap
    Then we claim that $\exists b \in \calM[G]$ such that $\calM[G] \models \forall x \in a \, \exists y \in b. \, \varphi(x,y,v)$. To prove this claim, let $C = \{ \langle p, \dot{x} \rangle \, | \, p \in \bbP, \, \dot{x} \in \ran(\dot{a}), \, \exists \dot{y} \in \calM^{\bbP} (p \forces \varphi(\dot{x},\dot{y},\dot{v}))^{\calM}\}.$ Now, for all $\langle p, \dot{x} \rangle \in C$. $\exists \dot{y} \in \calM^{\bbP} (p \forces \varphi(\dot{x},\dot{y},\dot{v}))^{\calM}$. \gap
    Note that there may be class many $\dot{y}$ such that $p \forces \varphi(\dot{x},\dot{y},\dot{v})$. Using collection in $\calM$, there is a set $B \in \calM$, $B \subseteq \calM^{\bbP}$ such that $\forall \langle p, \dot{x} \rangle \in C, \, \exists dot{y} \in B. \, (p \forces \varphi(\dot{x}, \dot{y}, \dot{v}))^{\calM}$. Finally, set $\dot{b} = \{\langle \bbone, \dot{y} \rangle \, | \, \dot{y} \in B\} \in \calM^{\bbP}$. Now to show that $\dot{b}^{G}$ suffices, fix $x \in a$. We can find $\langle q, \dot{x} \rangle \in \dot{a}$ such that $q \in G$ and $\dot{x}^{G} =x$. By our assumption, $\calM[G] \forces \exists y \, \varphi(x,y,v)$. So fix $\dot{z}^{G}$ such that $\calM[G] \models \varphi(x, \dot{z}^{G}. v)$. \gap
    Then by the forcing theorem, fix $p \in G$ such that $(p \forces \varphi(\dot{x}, \dot{z}, \dot{v}))^{\calM}$. From this, it follows that $\langle p, x \rangle \in C$. So we can fix $\dot{y} \in B$ such that $(p \forces \varphi(\dot{x},\dot{y},\dot{v}))^{\calM}$. Therefore $\langle \bbone, \dot{y} \rangle \in \dot{b}$. Since $\bbone \in G$, $\dot{y}^{G} \in \dot{b}^{G}$. Finally, by the forcing theorem, $\calM[G] \models \dot{y}^{G} \in \dot{b}^{G} \wedge \varphi(\dot{x}^{G} \dot{y}^{G}, v)$.
    Note that we haven't used power set here, so if $\calM \models \sfZF^{-}$, then $\calM[G] \models \sfZF^{-}$.\endproof
    \newlec
    It only remains to show that if $\calM \models \sfZFC$, then $\calM[G] \models \text{ power set}$.\gap
    \textbf{Recall:} \begin{enumerate}
        \item[(a)] for all $\langle q_{1}, z_{1} \rangle \in \dot{x}$, $\{r \leq p \, | \, r \leq q_{1} \rightarrow \exists \langle q_{2}, z_{2}\rangle \in y. \, (r \leq q_{2} \wedge r \forces \dot{z}_{1} = \dot{z}_{2}) \}$ is dense below $p$.
        \item[(b)] For all $\langle q_{2}, \dot{z}_{2}\rangle \in \dot{y}$, $\{ r \in p \, | \, r \leq q_{2} \rightarrow \exists \langle q_{1}, \dot{z}_{1} \rangle \in \dot{x}. \, (r \leq q_{1} \wedge r \forces \dot{z}_{1} = \dot{z}_{2})\}$.
    \end{enumerate}
    Then, continuing to prove the forcing theorem, it remains to prove:
    \begin{fancylem}
        For any $\dot{x}, \dot{y} \in \calM^{\bbP}$, \begin{enumerate}
            \item[(i)] If $p \in G$ and $(p \forces \dot{x} = \dot{y})^{\calM}$, then $\calM[G] \models \dot{x}^{G} = \dot{y}^{G}$.
            \item[(ii)] If $\calM[G] \models \dot{x}^{G} = \dot{y}^{G}$, then $\exists p \in G (p \forces \dot{x} = \dot{y})^{\calM}$.
        \end{enumerate}
    \end{fancylem}
    \proof We proceed by transfinite induction on the pair $(\dot{x},\dot{y})$ when ordered lexicographically. \begin{enumerate}
        \item[(i)] Suppose $p \forces \dot{x}=\dot{y}$, $p \in G$. To prove that $\calM[G] \models \dot{x}^{G} = \dot{y}^{G}$, it suffices to show that $\calM[G] \models \dot{x}^{G} \subseteq \dot{y}^{G}$, and then we obtain equality by symmetry.\gap
              Any element of $\dot{x}^{G}$ must be of the form $\dot{z}_{1}^{G}$ where $\langle q_{1}, \dot{z}_{1}\rangle \in \dot{x}$, and $q_{1} \in G$. Since $G$ is a filter, fix $s \in G$ such that $s \leq p, \dot{q}_{1}$. Since $s \leq p$, $s \forces \dot{x} = \dot{y}$. So \[\{ r \leq s \, | \, r \leq q_{1} \rightarrow \exists \langle q_{2}, \dot{z}_{2} \rangle \in y. \, ( r \leq q_{2} \wedge r \forces \dot{z}_{1} = \dot{z}_{2})\} \cap G \neq \emptyset.\] Fix $r$ in this intersection. Then $r \leq s \leq q_{1}$. So fix $\langle q_{2}, \dot{z}_{2}\rangle \in \dot{y}$ such that $r \leq q_{2} \wedge r \forces \dot{z}_{1} = \dot{z}_{2}$. Since $G$ is generic, $q_{2} \in G$. Therefore $\dot{z}_{1}^{G} \in \dot{y}^{G}$. Then by induction, since $r \in G$, $\calM[G] \models \dot{z}_{1}^{G} = \dot{z}_{2}^{G}$. Thus, $\dot{z}_{1}^{G} \in \dot{y}^{G}$. So we have proven that $\calM \models \dot{x}^{G} \subseteq \dot{y}^{G}$.
        \item[(ii)] Suppose $\calM[G] \models \dot{x}^{G} = \dot{y}^{G}$. Then define $D$ to be the set of conditions $r \in \bbP$ such that at least one of the following hold \begin{enumerate}
                  \item[(0)] $r \forces \dot{x} = \dot{y}$.
                  \item [(a')]
                        $\exists \langle q_{1}, \dot{z}_{1} \rangle \in \dot{x}. \, (r \leq q_{1} \wedge \forall \langle q_{2}, \dot{z}_{2} \rangle \in \dot{y}, \, \forall s \in \bbP ((s \leq q_{2} \, \wedge$ $ s \forces \dot{z}_{1} = \dot{z}_{2}) \rightarrow s \perp r)$.
                  \item[(b')]$\exists \langle q_{2}, \dot{z}_{2} \rangle \in \dot{y} (r \leq q_{2} \, \wedge \, \forall \langle q_{1}, \dot{z}_{1}\rangle \in \dot{x}\, \forall s \in \bbP (( s \leq q_{1} \, \wedge $ $ s \forces \dot{z}_{1} = \dot{z}_{2}) \rightarrow s \perp r))$.
              \end{enumerate}
              Note by separation in $\calM$ and absoluteness results, $D \in \calM$. Then we claim: \gap
              \textbf{Claim:} $D$ is dense. \gap To show this, fix $p \in \bbP$ and suppose $p \nVdash \dot{x} = \dot{y}$. Then at least one of (a) or (b) in the definition of $\forces$ for equality fails. Suppose it were (a), then we can fix $\langle q_{1}, \dot{z}_{1}\rangle \in \dot{x}$ such that $\{ r \leq p \, | \, r \leq q_{1} \rightarrow \exists \langle q_{2}, \dot{z}_{2} \rangle \in \dot{y} ( r \leq q_{2} \, \wedge \, r \forces \dot{z}_{1} = \dot{z}_{2}) \}$ is not dense below $p$. This means we can fix $s \leq p$ such that \[\forall r \leq s. \, (r \leq q_{1} \, \wedge \, \forall \langle q_{2}, \dot{z}_{2} \rangle \in \dot{y}. \, \neg(r \forces \dot{z}_{1} = \dot{z}_{2} \, \wedge r \, \leq q_{2})) \tag{$*$}\]. Observe that this gives $s \leq q_{1}$. Now, if we fix $\langle q_{2}, \dot{z}_{2} \rangle \in \dot{y}$, $r \leq q_{2}$, and $r \forces \dot{z}_{1} = \dot{z}_{2}$, then it must be the case that $r \perp s$, as any common extension of $s$ and $r$ would contradict $(*)$. Thus, $s \leq p$ and $s$ satisfies $(a')$. \gap
              Since $D$ is dense, $D$ is dense below $p \in G$, $G$ is $\bbP$-generic, fix $r \in G \cap D$. We shall show $r$ satisfies (0), which finishes the proof. Suppose for contradiction that $r$ satisfies (a'), then we can fix $\langle q_{1}, z_{1} \rangle$ satisfying the statement of (a'). Since $r \in G$ and $r \leq q_{1}$, we must have that $q_{1} \in G$ by upwards closure of filters. Therefore $\calM[G] \models z_{1}^{G} \in \dot{x}^{G} = \dot{y}^{G}$. So we can fix $\langle q_{2}, \dot{z}_{2}\rangle \in y$ such that $q_{2} \in G$ and $\calM[G] \models \dot{z}_{1}^{G} = \dot{z}_{2}^{G}$. By the inductive hypothesis, fix $p' \in G$ such that $p' \forces \dot{z}_{1} = \dot{z}_{2}$. Now since $G$ is a filter, and $p',q_{2} \in G$, we can find $s \in G$ so that $s \leq p',q_{2}$. Then $s \forces \dot{z}_{1} = \dot{z}_{2}$, so by (a'), $s \perp r$. But now $s \in G$ and $r \in G$ so they must be compatible. Contradiction. Therefore (a') doesn't hold for $r$. By symmetry, neither can (b'). Then since $r \in D$ and one of (0), (a'), and (b') must hold, (0) must hold.
    \end{enumerate}
    \endproof
    Recall from last lecture that if $\calM \models \sfZF^{-}$, and $G$ is $\bbP$-generic, then $\calM[G] \models \sfZF^{-}$.
    \begin{fancylem}{}
        If $\calM \models \sfZF$, then $\calM[G] \models \sfZF$.
    \end{fancylem}
    \proof It only remains to prove power set. For this, it suffices to show that if $a \in \calM[G]$, then $\calP(a) \cap \calM[G] = \{ x \in \calM[G] \, | \, x \subseteq a\} \subseteq b$ for some $b \in \calM[G]$.\gap
    Fix $a \in \calM[G]$ with corresonding name $\dot{a}$. Then define \[ S = \{ x \in \calM^{\bbP} \, | \, \ran(x) \subseteq \ran(\dot{a})\} = \calP(\bbP \times \ran(a))^{\calM}.\]
    Let $\dot{b} = \{ \langle \bbone, \dot{x} \rangle \, | \, x \in S\} \in \calM^{\bbP}$. Let $c \in \calP(a) \cap \calM[G]$, we need to show that $c \in \dot{b}^{G}$. Let $\dot{c} \in \calM^{\bbP}$ be such that $c = \dot{c}^{G}$. Then let \[ \dot{x} = \{ \langle p, \dot{z} \rangle \, | \, \dot{z} \in \ran(\dot{a}) \wedge (p \forces \dot{z} \in \dot{c})^{\calM}\} \in S.\]
    \textbf{Claim:} $\dot{x}^{G} = \dot{c}^{G} = c$.
    \begin{enumerate}
        \item[($\dot{x}^{G} \subseteq c$)] To show this, fix $\dot{z}^{G} \in \dot{x}^{G}$. Then by definition, we can fix $p \in G$ such that $\langle p, \dot{z} \rangle \in x$. From this it follows that $z \in \ran(\dot{a})$ and $p \forces \dot{z} \in \dot{c}$. Since $p \in G$, by the Forcing Theorem, $\calM[G] \models \dot{z}^{G} \in c$.
        \item[($c \subseteq \dot{x}^{G}$)] Since $\calM[G] \models c \subseteq \dot{a}^{G}$, we know every element of $c$ is of the form $\dot{z}^{G}$, for some pair $\langle q, \dot{z}\rangle \in a$ with $q \in G$. Also, if $\calM[G] \models \dot{z}^{G} \in c$, then by the Forcing theorem, we can fix $p \in G$ which forces this, so $p \forces \dot{z} \in c$. Thus $\langle p, z \rangle \in x$. So $\dot{z}^{G} \in \dot{x}^{G}$.
    \end{enumerate}
    \endproof
    \begin{fancylem}{}
        $\calM \models \sfZFC^{-} \Rightarrow \calM[G] \models \sfZFC^{-}$. (And obviously also for $\sfZFC$, by the above).
    \end{fancylem}
    \proof To show this, we will show well-ordering holds. It suffices to show that any $a \in \calM[G]$ can be well-ordered in $\calM[G]$. (There is definitely one in the universe, since $a$ is countable, but we need to show that it exists in $\calM[G]$). Using well-ordering in $\calM$, list the elements of $\ran(a)$ as $\{ \dot{x}_{\alpha} \, : \, \alpha < \delta \}$ for some $\delta \in \Ord$. \gap
    Then let $\dot{f} = \{ \langle \bbone, \op(\check{\alpha},x_{\alpha}) \rangle \, | \, \alpha < \delta \} \in \calM^{\bbP}$. Then in $\calM[G]$, $\dot{f}^{G} = \{ \langle \alpha, \dot{x}_{\alpha}^{G}\rangle \, | \, \alpha < \delta\}$. So $\dot{f}^{G}$ is a function with domain $\delta$ and $a \subseteq \ran(zdot{f}^{G})$. Therefore, define a well-order on $a$ by saying $x < y$ if and only if \[\min\{ \alpha < \delta \, | \, \dot{f}^{G}(\alpha) = \dot{x}\} < \min\{ \alpha < \delta \, | \, \dot{f}^{G}(\alpha) = y\}.\]
    Which finishes the proof. \endproof
    \begin{remark*}\hspace{0.1cm}
        \begin{enumerate}
            \item $\dot{f}^{G}$ may not be injective, we could have $\dot{x}_{\alpha}^{G} = \dot{x}_{\beta}^{G}$.
            \item $\ran(\dot{f}^{G})$ may not equal $a$. (Elements of $\dot{a}$ are $\langle p, x_{\alpha} \rangle$, if $p \notin G$, we may not have $x_{\alpha}^{G} \in a$).
            \item For Power set it sufficed to find a set of names which contained enough names to represent all \underline{possible} subsets of $a$.
            \item $\calM[G] \models \varphi$ should be considered as a ternary relation, it is a statement about $\calM$, about $G$, and about $\varphi$. We could have two generic filters $G$ and $H$, such that $\calM[G] \models \varphi$ and $\calM[H] \models \neg \varphi$.
            \item We were careful to write $(p \forces \varphi)^{\calM}$ with the relativisation to $\calM$. However, we will drop this relativisation when it is clear.
        \end{enumerate}
    \end{remark*}
























    \newlec
    \begin{fancylem}{}
        Let $\varphi$ and $\psi$ be two formulas in our forcing language. Let $\calM$ be a countable transitive model of $\sfZFC$, and let $\bbP \in \calM$. Then for any $p \in \bbP$, and $x \in \calM^{\bbP}$: \begin{enumerate}
            \item[(i)] If $\sfZF \proves \forall v. \, (\varphi(v) \rightarrow \psi(v))$ then \[ (p \forces \varphi(\dot{x}))^{\calM} \Rightarrow (p \forces \psi(\dot{x}))^{\calM}\]
            \item[(ii)] If $\sfZFC \proves \forall v. \, (\varphi(v) \leftrightarrow \psi(v))$, then \[ (p \forces \varphi(\dot{x}))^{\calM} \Leftrightarrow (p \forces \psi(\dot{x}))^{\calM}\]
        \end{enumerate}
    \end{fancylem}
    \proof Suppose $\sfZFC \proves \forall v. \, (\varphi(v) \rightarrow \psi(v))$. and $(p \forces \varphi(\dot{x}))^{\calM}$. Since $\calM$ is countable, let $G$ be a $\bbP$-generic filter over $\calM$ such that $p \in G$. By the forcing theorem, $\calM[G] \models \varphi(\dot{x}^{G})$. Then since $\calM[G] \models \sfZFC$, we must have $\calM[G] \models \psi(\dot{x}^{G})$. By the forcing theorem, since true for all generics containg $p$, $(p \forces \psi(\dot{x}))^{\calM}$. \endproof
    \section{Forcing and Independence Results}
    We want to show $\Cons(\sfZFC + V \neq L)$ assuming $\Cons(\sfZFC)$. Before we do that, we need one theorem:
    \begin{fancythm}{}
        Let $\calM$ be a countable transitive model of $\sfZFC$, then there is a countable transitive model $N \supseteq M$ of $\sfZFC$ such that $N \models \sfZFC + V \neq L$.
    \end{fancythm}
    \proof Let $\calM$ be a countable transitive model of $\sfZFC$, and $\bbP \in \calM$ be any atomless forcing poset (i.e. having no minimal element). The obvious example of this is $\Fn(\omega, 2)$. Since $\calM$ is countable, let $G$ be $\bbP$-generic over $\calM$. Since $\bbP$ is atomless, $G \notin \calM$. So $M \subsetneq M[G]$ is a model of $\sfZFC$. By the generic model theorem, $\Ord \cap \calM = \Ord \cap \calM[G]$. Therefore, $L_{\Ord \cap \calM} = \calL_{\Ord \cap M[G]} = L^{\calM} \subseteq \calM \subsetneq \calM[G]$. Therefore $M[G] \neq L_{\Ord \cap \calM[G]} = L^{\calM[G]}$. So $(V \neq L)^{\calM[G]}$. \endproof
    This doesn't yet give us the consistency proof, since we had to assume the existence of a countable transitive model. Now we prove: \begin{fancythm}{}
        $\Con(\sfZFC) \Rightarrow \Con(\sfZFC + V \neq L)$. Therefore $\sfZFC \notproves V=L$.
    \end{fancythm}
    \proof Assume that $\sfZFC + V \neq L$ gives rise to a contradiction. Then, from a finite set of axioms $\Gamma$ of $\sfZFC + V \neq L$, we can find $\psi$ such that $\Gamma \proves \psi \wedge \neg \psi$. By following the previous proofs, there is a finite set of axioms $\Lambda$ of $\sfZFC$ such that: $\sfZFC \proves$ if there is a countable transitive model for $\Lambda$, then there is a countable transitive model for $\Gamma + V \neq L$. \gap
$\Lambda$ should be sufficient to: \begin{itemize}
        \item Prove the basic properties of forcing and constructibility.
        \item Prove the necessary absoluteness facts (e.g. \say{finite}, \say{partial order}).
        \item Prove facts about forcing (the forcing theorem)
        \item Finally, to be such that $\sfZFC$ can prove that if $\calM$ is a countable transitive model of $\Lambda$, $\bbP \in \calM$. $G$ is $\bbP$-generic over $\calM$, then $\calM[G] \models \Gamma$.
    \end{itemize}
    Finally, by reflection, since $\Lambda$ is finite, and a subset of $\sfZFC$, there is a countable transitive model of $\Lambda$. Therefore there is a countable transitive model of $\Gamma + V \neq L$. But $\Gamma \proves \psi \wedge \neg \psi$. So $\calN \models \psi \wedge \neg \psi$. So $\psi \wedge \neg \psi)^{\calN}$. So by relativisation, $\sfZFC \proves \psi^{\calN} \wedge \neg \psi^{\calN}$. So $\neg \Con(\sfZFC)$. \endproof
    \begin{remark*}
        Someone showed that this can be done in Isabelle (very recently). This is called \say{The formal verification of the ctm approach to forcing.} \gap
        To get $\Con(\sfZFC) \Rightarrow \Con(\sfZFC + \neg CH)$ takes $\sfZC$ plus 21 instances of replacement.
    \end{remark*}
    \subsection{Cohen Forcing}
    Fix $\calM$ to be a countable transitive model of $\sfZFC$. Recall that for $I,J \in \calM$: \begin{itemize}
        \item $\Fn(I,J)= \{ p \, | \, |p| < \omega, \, p \text{ is a function}, \dom(p) \subseteq I, \, \ran(p) \subseteq J\}$.
        \item Also $(\Fn(I,J),\supseteq,\emptyset)$ is a forcing poset.
        \item $\Fn(I,J) \in M$.
        \item $\Fn(I,J)$ has the countable chain condiion iff $I \neq \emptyset$ or $J$ is countable.
        \item The sets $D_{i} = \{ q \in \Fn(I,J)\, | \, i \in \dom(q)\}$, and $R_{j} = \{ q \in \Fn(I,J)\, | \, j \in \ran(q)\}$ are dense for all $i \in I$, $j \in J$.
    \end{itemize}
    Now, suppose that $G \subseteq \Fn(I,J)$ ws $\Fn(I,J)$-generic over $\calM$. Since $G$ is a filter, if $p,q \in G$, then $p \cap q \in G$. So if $p, q \in G$, then $p$ and $q$ agree on the intersection of their domains. \gap
    Since wherever they overlap, they agree, we can take $f_{G} = \cup G$. Then $f_{G}$ is a function with $\dom(f_{G}) \subseteq I$, and $\ran(f_{G}) \subseteq J$. Note that a name for $f_{G}$ is \[ \dot{f} = \{ \langle p, \op(\check{i}, \check{j})\rangle \, | \, p \in \bbP, \langle i, j \rangle \in p\}.\]
    Since $D_{i}$ and $R_{j}$ are dense for all $i,j$ $G \cap D_{i} \neq emptyset$ %WHAT IS C_R?
    so $i \in \dom(f_{G})$. Thus
    \begin{fancyprop}{}
        If $G \subseteq \Fn(I,J)$ is $\Fn(I,J)$-generic over $\calM$ and $I,J \neq \emptyset$, then $\calM[G] \models f_{G} : I \to J$ is a surjection.
    \end{fancyprop}
    \proof Omitted. \endproof
    \begin{fancyprop}{}
        Suppose that $I,J$ are non-empty sets, at least one of which is infinite. Then $|\Fn(I,J)| = \max\{|I|,|J|\}$.
    \end{fancyprop}
    \proof Each $p \in \Fn(I,J)$ is a finite function, so $\Fn(I,J) \subseteq {} ^{<\omega}(I \times J)$. Then: \begin{align*}
        |\Fn(I,J)| & \leq |(I \times J)^{< \omega}| \\
                   & = |I \times J|                 \\
                   & = \max\{|I|,|J|\}.
    \end{align*}
    For the reverse, fix $i_{0} \in I$, and $j_{0} \in J$. Then $\{ \langle i_{0}, j \rangle \, | \, j \in J\} \cup \{ \langle i, j_{0} \rangle \, | \, i \in I\}$ is a collection of $|I \cup J|$ many distinct elements of $\Fn(I,J)$. Then $|I \cup J| = \max\{|I|, |J|\} \leq |\Fn(I,J)|$. \endproof
    To provide a model where $\sfCH$ fails, consider $\Fn(\omega_{2}^{\calM} \times \omega, 2)$. Consider $f_{G} : \omega_{2}^{\calM} \times \omega \to 2$. Let $g_{\alpha} : \omega \to 2$. be defined by \[ f_{G}(\alpha, n) = g_{\alpha}(n).\]
    In order to get that $\calM[G] \models \sfZFC + \neg \sfCH$, it remains to:
    \begin{itemize}
        \item Prove that for $\alpha \neq \beta$, $g_{\alpha} \neq g_{\beta}$.
        \item Prove that $\omega_{1}^{\calM[G]}=\omega_{1}$ and $\omega_{2}^{\calM[G]} = \omega_{2}^{\calM}$.
    \end{itemize}
    It will turn out that the countable chain condition (ccc) guarantees that all cardinals in $\calM$ remain cardinals in $\calM[G]$. This is not trivially true, consider the following examples: \begin{ex*}
        Let $\kappa$ be an uncountable cardinal in $\calM$. Consider $\Fn(\omega, \kappa)$. Then in $\calM[G]$, $f_{G} : \omega \to \kappa$ is a surjection. Therefore $(\kappa \text{ is countable})^{\calM[G]}$. The reason this broke was because $\Fn(\omega, \kappa)$ does not have the countable chain condition.
    \end{ex*}
    \newlec
    \begin{fancydefine}{Preserving cardinals/cofinalities}
        Let $\bbP \in \calM$ be a forcing poset. Then we say that: \begin{enumerate}
            \item[(i)] $\bbP$ \textbf{preserves cardinals} iff for every generic filter $G \subseteq \bbP$ on $\calM$ \[ (\text{$\kappa$ is a cardinal})^{\calM} \Leftrightarrow (\text{$\kappa$ is a cardinal})^{\calM[G]}\]
                  for all $\kappa \in \Ord$.
            \item[(ii)] $\bbP$ \textbf{preserves cofinalities} iff for every genric filter $G \subseteq \bbP$ over $\calM$ \[ \cf^{\calM}(\gamma)= \cf^{\calM[G]}(\gamma)\] for all limit ordinals $\gamma \in \Ord \cap \calM$.
        \end{enumerate}
    \end{fancydefine}
    \begin{remark*}\hspace{0.1cm}
        \begin{itemize}
            \item Being a cardinal is $\Pi_{1}$ definable, so downwards absolute.
            \item Finite and $\omega$ are absolute.
        \end{itemize}
        Also, remember that we showed the forcing $\Fn(\omega, \kappa)$ collapses $\kappa$, so there are forcings which do \textit{not} preserve cardinals.
    \end{remark*}
    \begin{fancylem}{}
        Suppose $\bbP \in \calM$ is our forcing poset. Then \begin{enumerate}
            \item[(i)] $\bbP$ preserves cofinalities iff for all $\bbP$-generic filters $G$ for all limit cardinals $\beta$ with $\omega < \beta < \Ord \cap \calM$, \[ (\text{$\beta$ is regular})^{\calM} \rightarrow (\text{$\beta$ is regular})^{\calM[G]}.\]
            \item[(ii)] If $\bbP$ preserves cofinalities, then $\bbP$ preserves cardinals.
        \end{enumerate}
    \end{fancylem}
    \proof \hspace{0.1cm}
    \begin{enumerate}
        \item[(i, $\Rightarrow$)] Suppose $\bbP$ preserves cofinalities, and $G$ is $\bbP$-generic. Fix $\beta$, $\omega < \beta < \Ord \cap \calM$, a limit. Then if $(\text{$\beta$ is regular})^{\calM}$, $\beta = \cf^{\calM}(\beta) = \cf^{\calM[G]}(\beta)$. So if $\beta$ is regular in $\calM$ then it is regular in $\calM[G]$.
        \item[(i, $\Leftarrow$)] Let $\gamma$ be a limit ordinal with $\omega < \gamma < \Ord \cap \calM$. Let $\beta = \cf^{\calM}(\gamma)$. Then $\beta$ is regular in $\calM$, since the cofinality of an ordinal is always regular. Let $f \in \calM$, $f: \beta \to \gamma$ be strictly increasing and cofinal. Then if $\beta$ is uncountable in $\calM$, then $(\text{$\beta$ is regular})^{\calM[G]}$ holds, by assumption. Otherwise, $\beta= \omega$. Then $(\beta = \omega)^{\calM[G]}$ so $(\text{$\beta$ is regular})^{\calM[G]}$. Finally, since $f \in \calM$, $f \in \calM[G]$, so there is a strictly increasing cofinal map from $\beta$ to $\gamma$ in $\calM[G]$. Thus, $\cf^{\calM[G]}(\gamma) = \cf^{\calM[G]}(\beta) = \beta = \cf^{\calM}(\gamma)$, so cofinalities are preserved as required.
        \item[(ii)] Suppose $\bbP$ preserves cofinalities and let $\kappa$ be a cardinal in $\calM$. One of three cases occur, \begin{enumerate}
                  \item $\kappa \leq \omega$, so $(\kappa \leq \omega)^{\calM[G]}$, so $(\text{$\kappa$ is a cardinal})^{\calM[G]}$.
                  \item $\kappa$ is regular. Then by (i), $(\text{$\kappa$ is regular})^{\calM[G]}$. In particular, regular implies cardinal, so $(\text{$\kappa$ is a cardinal})^{\calM[G]}$.
                  \item $(\text{$\kappa$ is singular})^{\calM}$. Then $\kappa$ is the supremum of a set of regular cardinals. So since $\bbP$ preserves regular cardinals, every element of this set is regular in $\calM[G]$. Therefore $\kappa$ is the supremum of a set of cardinals in $\calM[G]$, and thus a cardinal in $\calM[G]$ (by a question on Example Sheet 1).
              \end{enumerate}
    \end{enumerate}
    \endproof
    \begin{fancylem}{}
        Let $A,B, \bbP \in \calM$. Then if $(\text{$\bbP$ has the ccc})^{\calM}$, and $G$ is $\bbP$-generic over $\calM$. Then for any function $f \in \calM[G]$ with $f : A \to B$, there is a fucntion $F \in \calM$ with $F : A \to \calP^{\calM}(B)$ such that for all $a \in A$, $f(a) \in F(a)$, $|F(a)| \leq \aleph_{0}$.
    \end{fancylem}
    \proof Suppose that $f : A \to B$ is our function in $\calM[G]$. Since $A, B$ are in $\calM$, let $\check{A}, \check{B}$ be the canonical names. Let $\dot{f}$ be a name for $f$. Then, by the forcing theorem, there is a $p \in G$ such that \[ p \forces \dot{f} : \check{A} \to \check{B}\] is a function. \gap
    Define $F : A \to \calP^{\calM}(B)$ by $F(a) = \{ b \in B \, | \, \exists q \leq p. \, (q \forces \dot{f}(\check{a}) = \check{b}) \}$. Then we prove that $F(a)$ has the desired properties. \gap
    \textbf{Claim:} $F \in \calM$. \gap
    This is fine, since by the definability of the forcing relation, we have that $F(a) \in \calM$ for all $a \in \calM$. So, since $A \in \calM$, $F = \{ \langle a, F(a) \rangle \, | \, a \in A\} \in M$ by replacement. \gap
    \textbf{Observe:} $F$ is a function.\gap
    \textbf{Claim:} $f(a) \in F(a)$.\gap Suppose that $\calM[G] \models f(a) = b$ for $b \in B$. Then by the forcing theorem, we can fix $q \in G$ such that $q \forces \dot{f}(\check{a})=\check{b}$. Since $G$ is a filter, fix $r \leq p, q$ with $r \in G$. Then $r$ forces that $\dot{f}$ is a function, and that $\dot{f}(\check{a})= \check{b}$. So $b \in F(a)$. \gap
    \textbf{Claim:} $|F(a)| \leq \aleph_{0}$.\gap
    Working in $\calM$, using the axiom of choice in $\calM$. For each $b \in F(a)$, choose $q_{b} \leq p$ such that $q_{b} \forces \dot{f}(\check{a})=\check{b}$. We shall show that $q_{b} \perp q_{c}$ for $b \neq c$, and then the countable chain condition will entail immediately that $|F(a)| \leq \aleph_{0}$. \gap
    Suppose not, then $r \leq q_{b},q_{c}$ since they are compatible. Then $r \forces \dot{f} : \check{A} \to \check{B} \text{ is a function}\, \wedge \, \check{b} \neq \check{c} \, \wedge \, f(\check{a})=\check{b} \, \wedge \, f(\check{a})=\check{c}$. Now let $H$ be a generic filter with $r \in H$. Then $r \leq p$ and $\calM[H] \models f: A \to B \text{ is a function} \, \wedge \,$ $f(a)=b \, \wedge \, f(a)=c \, \wedge b \neq c$. This gives a contradiction.
    \endproof
    \begin{fancythm}{}
        If $\bbP \in \calM$ and $(\bbP \text{ has ccc})^{\calM}$. then $\bbP$ preserves cofinalities and hence cardinals.
    \end{fancythm}
    \proof Using the previous lemma, it suffices to show that if $\omega < \beta < \Ord \cap \calM$, and $\beta$ is a limit, then if $(\text{$\beta$ is regular})^{\calM}$ then $(\text{$\beta$ is regular})^{\calM[G]}$. Suppose for a contradiction that $\beta$ is regular in $\calM$ but singular in $\calM[G]$. \gap Then in $\calM[G]$ fix a cofinal map $f : \alpha \to \beta$ for some $\alpha < \beta$. So in $\calM$ there is some function $F : \alpha \to \calP^{\calM}(\beta)$ such that for all $\gamma \in \alpha$, $f(\gamma) \in F(\gamma)$, and $|F(\gamma)| \leq \aleph_{0}$. Then let \[ X = \bigcup_{\gamma < \alpha} F(\gamma).\]
    Then we have that $X \subseteq \beta$, and $X$ is the union of less than $\beta$-many countable sets. So $X \neq \beta$. But, $f$ was a cofinal function, so: \[ \beta \bigcup_{\alpha < \gamma} f(\alpha). \] But this means $\beta$ is contained in $X$ by the definition of our approximation map $F$. Therefore $\beta = X$. But this is a contradiction. \endproof

    \subsection{The Failure of $\sfCH$}
    \begin{fancythm}{}
        Fix $\alpha < \Ord \cap \calM$ and let $\kappa = (\aleph_{\alpha})^{\calM}$. Let $\bbP = \Fn(\kappa \times \omega, 2)$, and let $G$ be $\bbP$ generic over $\calM$. Then $\calM[G]$ contains a $\kappa$-length sequence of distinct elements of $2^{\omega}$. Hence, \[ \calM[G] \models \sfZFC + 2^{\aleph_{0}} \geq \kappa = \aleph_{\alpha}.\]
    \end{fancythm}
    \proof Let $f = \cup G \in \calM[G]$. Then we have already shown earlier that $f$ is a function from $\kappa \times \omega$ onto $2$. For $\beta < \kappa$, let $g_{\beta} : \omega \to 2$, $g_{\beta}=f(\beta,n)$. \gap
    \textbf{Claim:} $\alpha \neq \beta$, $g_{\alpha} \neq g_{\beta}$. Then define $E_{\alpha,\beta} \in \calM$ by \[ E_{\alpha. \beta} = \{ q \in \bbP \, | \, \exists n. \, (\langle \beta, n\rangle, \langle \alpha, n \rangle \in \dom(q) \wedge q(\langle \beta, n \rangle) \neq q(\langle \alpha, n \rangle))\}\]
    Then to prove that $E_{\alpha, \beta}$ is dense, fix $p \in \bbP$. Since $p$ is finite, tere is some $m$ such that $\langle \beta, m \rangle, \langle \alpha, m \rangle \notin \dom(p)$. \gap
    Define $q \leq p$ by $q : \dom(p) \cup \{\langle \beta, m \rangle, \langle \alpha, m \rangle \} \to 2$ by \begin{itemize}
        \item $p(z) = q(z)$ if $z \in \dom(p)$.
        \item $q(\langle \beta, m \rangle) = 1$.
        \item $q(\langle \alpha, m \rangle) =0$.
    \end{itemize}
    Then $q \in E_{\alpha, \beta}$. Since $G$ is $\bbP$-generic, fix $q \in G \cap E_{\alpha, \beta}$. Then: \begin{align*}
        g_{\beta}(m) & = f(\beta, m) = q(\langle \beta, m \rangle)                       \\
                     & \neq q(\langle \alpha, m \rangle) = f(\alpha, m) = g_{\alpha}(m).
    \end{align*}
    for the chosen $m$.\gap
    Finally, since $\bbP$ has the ccc in $\calM$, $\bbp$ preserves cardinals, so $\kappa = (\aleph_{\alpha}^{\calM[G]})$.
    \newlec
    To recap the previous lecture, we had that if $\calM$ were a countable transitive model of $\sfZFC$, and $\alpha < \Ord \cap \calM$, and $(\kappa = \aleph_{\alpha})^{\calM}$, $\bbP = \Fn(\kappa \times \omega, 2)$. Then letting $G$ be $\bbP$-generic, and $\cup G : \kappa \times \omega \to 2$. \gap Then for $\gamma < \kappa$, we defined $g_{\gamma} : \omega \to 2$ so that $g_{\gamma}(n) = \cup G(\gamma, n)$. Then if $\gamma \neq \gamma'$, $g_{\gamma} \neq g_{\gamma'}$. So $\calM[G] \models 2^{\aleph_{0}}\geq \kappa = \aleph_{\alpha}$. Taking $\alpha=2$ gives us a model of $\sfZFC$ where the continuum hypothesis fails.
    \begin{fancythm}{}
        $\Con(\sfZFC) \Rightarrow \Con(\sfZFC + \neg \sfCH)$.
    \end{fancythm}
    \proof We won't rewrite this proof, since it follows from the proof last lecture. \endproof
    \begin{fancydefine}{Cohen reals}
        The $g_{\gamma}$'s that we added to $\calM$ are called \textbf{Cohen reals}.\gap
        In particular, we say that $c : \omega \to 2$ is a Cohen real over $\calM$ iff \[ \exists H. \, (H \text{ is $\Fn(\omega, 2)$-generic over $\calM$ and $c = \cup H$.})\]
    \end{fancydefine}
    \begin{fancyprop}{K\"onig}
        $2^{\aleph_{0}} \neq \kappa$ for any $\kappa$ with cofinality $\aleph_{0}$.
    \end{fancyprop}
    \proof This was proven on the first example sheet. \endproof
    Moreover, under $\sfGCH$, for any $\kappa$, we have that $\cf(\kappa) \neq \omega$ iff $\kappa^{\omega}=\kappa$.\gap
    In the proof of the power sex axiom in $\calM[G]$, we showed that given $a \in \calM^{\bbP}$, a name for its power set was $\calP(\bbP \times \ran(a))$. But this is a very large object, and we want to get a better bound than this.
    \begin{fancythm}{}
        Let $\calM$ be a (countable) transitive model of $\sfZFC$, assume that $(\kappa = \aleph_{\alpha} \wedge \kappa^{\omega}=\kappa)^{\calM}$. Let $\bbP = \Fn(\kappa \times \omega, 2)$ and let $G$ be $\bbP$-generic over $\calM$. Then $\calM[G] \models 2^{\aleph_{0}}=\aleph_{\alpha}=\kappa$.
    \end{fancythm}
    \proof We already have $\calM[G] \models \sfZFC + \kappa=\aleph_{\alpha}\leq 2^{\aleph_{0}}$. \gap
    To show the other direction, let $\dot{x}$ be a name for a subset of $\omega$. For $n \in \omega$, let $E_{\dot{x},n} = \{ p \in \bbP \, | \, (p \forces \check{n} \in \dot{x}) \vee (p \forces \check{n} \notin \dot{x}) \}$. Then $E_{\dot{x},n}$ is dense in $\bbP$\footnote{Because if we take any condition $q$, then either $q$ forces that $\check{n} \in \dot{x}$, or $q$ does not force it. If $q$ does not force it, that means it must be dense below that there is a condition that forces its negation, from the definition.}. Next, for each $n \in \omega$, choose a maximal antichain $A_{\dot{x},n} \subseteq E_{\dot{x},n}$ (this is possible by ES3Q48). \gap
    Now, define $\dot{z}_{x} = \bigcup\limits_{n \in \omega} \{ \langle p, \check{n} \rangle \, | \, p \in A_{\dot{x},n} \wedge p \forces \check{n} \in \dot{x} \}$. We call such a name \say{nice.} Then: \gap
    \textbf{Claim:} $\bbone \forces x = \dot{z}_{x}$. \gap
    To show this, it suffices to prove that for every $n \in \omega$, the set $D_{\dot{x},n} = \{ q \in E_{\dot{x},n} \, | \, (q \forces \check{n} \in x) \Leftrightarrow (q \forces \check{n} \in \dot{z}_{x})\}$ is dense. Fix $n \in \omega$ and $p \in \bbP$. Then since $E_{\dot{x},n}$ is dense, we can fix some condition $p_{0} \leq p$ such tat $p_{0} \in E_{\dot{x},n}$. Now because $A_{\dot{x},n}$ is a maximal antichain, we can fix our condition $q_{0} \in A_{\dot{x},n}$ such that $p_{0} \, || \, q_{0}$, i.e. such that $p_{0}$ and $q_{0}$ are compatible. Now fix $r \leq p_{0}, q_{0}$. We show that $r \in D_{\dot{x},n}$.
    \begin{itemize}
        \item If $r \forces \check{n} \in \dot{x}$, then $q_{0} \forces \check{n} \in \dot{x}$. Thus $\langle q_{0}, \check{n} \rangle \in \dot{z}_{x}$. So $r \forces \check{n} \in \dot{z}_{x}$.
        \item If $r \forces \check{n} \in \dot{z}_{x}$, then by definition, $\{ s \leq r \, | \, \exists \langle q_{1}, \check{m} \rangle \in \dot{z}_{x}. \, ( s \leq q_{1} \wedge s \forces \check{m} = \check{n})\}$ is dense below $r$. This can only happen if there is some $\langle q_{1}, \check{n} \rangle \in \dot{z}_{x}$ such that $r \, || \, q_{1}$. So this $q_{1}$ must be in the antichain, but also we had that $q_{0}$ is in the antichain. However, they are both compatible with $r$. This is only possible if $q_{0} = q_{1}$, so $\langle q_{0}, \check{n} \rangle \in \dot{z}_{x}$. Thus $q_{0} \forces \check{n} \in x$. So, since $r \leq q_{0}$, $r \forces \check{n} \in \dot{x}$. Thus $r \in D_{\dot{x},n}$.
    \end{itemize}
    We have shown that every subset of $\omega$ has a \say{nice} name. Computing the number of nice names: \begin{itemize}
        \item $|\bbP| = \kappa$.
        \item Since $\bbP$ has the ccc, every antichain is countable.
        \item Thus, there are at most $(\kappa^{\omega} \times \omega)^{\omega} = \kappa$\footnote{$\kappa^{\omega}$ is the number of antichains, multiplied by $\omega$ because our conditions are of the form $\langle p, \check{n}\rangle$, and then all raised to the $\omega$, since we have to do this for each $n \in \omega$.} many \say{nice} names.
    \end{itemize}
    So $\calM[G] \models 2^{\aleph_{0}}=\kappa$. \endproof
    \subsection{Precise Values of the Continuum}
    \begin{corr}{}
        $\Con(\sfZFC) \Rightarrow \Con(\sfZFC + 2^{\aleph_{0}}=\aleph_{2}$). But becuse we worked in such generality, it could be $\aleph_{12}$, it could be $\aleph_{\omega^{3}+2}$, even $\aleph_{\omega_{1}}$.
    \end{corr}
    \begin{corr}{}
        The following are equiconsistent: \begin{itemize}
            \item $\sfZFC + \exists \text{ a weakly inaccessible cardinal}$
            \item $\sfZFC + \exists \text{ a strongly inaccessible cardinal} + \sfGCH$.
            \item $\sfZFC + 2^{\aleph_{0}} \text{ is weakly inaccessible}$.
            \item $\sfZFC + \exists \kappa, \, \text{weakly inaccessible and not strongly inaccessible.}$
        \end{itemize}
    \end{corr}
    \proof To go from the first to the second, just go into $L$. This was shown in an example sheet. To go from the third to the fourth, $2^{\aleph_{0}}$ exhibits such a $\kappa$. To go from the fourth to the first is immediate. All that's left is going from the second to the third, \say{can we have the continuum being weakly inaccessible?} And the answer is yes, we can, just take the forcing $\Fn(\kappa \times \omega,2)$ where $\kappa$ is the strongly inaccessible. Then $2^{\aleph_{0}}$ gives a weakly inaccessible cardinal. \endproof
    \begin{remark*}
        When trying to build models of $\sfZFC + 2^{\aleph_{0}} = \kappa$, we often assume GCH just for convenience (i.e. we work in $L$ or something like $L$) because we usually only care about consistency results anyway.
    \end{remark*}
    \begin{ex*}
        What happens when we consider $\Fn(\aleph_{\omega} \times \omega, 2)$. Let $G$ be $\Fn(\aleph_{\omega} \times \omega, 2)$-generic. Then $\calM[G] \models 2^{\aleph_{0}} \geq \aleph_{\omega}$. Then we want to get an upper bound, which we know (by K\"onig's lemma) cannot be $\aleph_{\omega}$. Assuming GCH, \[ \cf(\kappa) = \omega \Rightarrow \kappa^{\omega} = \kappa^{+}.\]
        Thus \begin{align*}
            \calM[G] \models \aleph_{\omega} & \leq 2^{\aleph_{\omega}} \\
                                             & \leq \aleph_{\omega}^{+}
        \end{align*}
        So $\calM[G] \models 2^{\aleph_{0}} = \aleph_{\omega+1}$. So this is what happens in this case.\gap
        Another remark on GCH is that we can have $2^{\aleph_{0}} < \aleph_{\omega}$, but $\aleph_{\omega}^{\aleph_{0}} = \aleph_{\omega+1}^{\aleph_{0}}= \aleph_{\omega+2}$. This is hard to prove and requires the existence of certain large cardinals.
    \end{ex*}
    If $\calM \models 2^{\aleph_{0}} = \aleph_{\alpha} > \aleph_{\beta}$, and $\Fn(\aleph_{\beta} \times \omega, 2)$. and $G$ is generic on this set. Then the continuum is greater than the number of reals we add, so we have $\calM[G] \models 2^{\aleph_{0}} = \aleph_{\alpha}$.
    \begin{remark*}
        The following are equiconsistent: \begin{itemize}
            \item $\sfZFC + \exists \text{ measurable} + \sfCH$
            \item $\sfZFC + \exists \text{ measurable} + \neg \sfCH$.
        \end{itemize}
        And the same for I0-I3, and it still holds even if we replace $\sfCH$ with $\sfGCH$.
    \end{remark*}
    \begin{remark*}
        $\sfZFC + $ Proper Forcing Axiom implies $2^{\aleph_{0}} = \aleph_{2}$.
    \end{remark*}
    \newlec
    \section{Generalized Cohen Forcing}
    Suppose $\sfZFC + \sfCH + 2^{\aleph_{1}} = \aleph_{3}$. Our first idea might be to have be our forcing poset $\Fn(\omega_{3} \times \omega_{1},2)$. The issue is that this doesn't give us $\sfCH$.
    \begin{fancyprop}{}
        Suppose $\calM$ is a transitive model of $\sfZFC + \sfGCH$, and we have $\kappa$ so that $(\kappa=\aleph_{\alpha} \wedge \kappa^{\omega}=\kappa)^{\calM}$. Let $\bbP = \Fn(\kappa \times \omega, 2)$, and let $G$ be $\bbP$-generic. Then for any cardinal $\lambda$ in $\calM$ such that $\aleph_{0} \leq \lambda < \kappa$ in $\calM[G]$, then \[ 2^{\lambda} = \begin{cases}
                \kappa,    & \text{if $\cf(\kappa) > \lambda$}     \\
                \kappa^{+} & \text{if $\cf(\kappa) \leq \lambda$.}
            \end{cases}\]




        Then there is a natural bijection between $\omega_{3} \times \omega$ and $\omega_{3} \times \omega_{1}$. From this, it will follow that in $\calM[G]$: \[ 2^{\aleph_{0}}=2^{\aleph_{1}}=\aleph_{3}.\]
    \end{fancyprop}
    \proof Omitted? or maybe we're about to do it. \endproof %Flag
    \begin{fancydefine}{$\Fn_{\kappa}$}
        Let $I,J$ be sets and $\kappa$ a regular cardinal. Then \begin{itemize}
            \item $\Fn_{\kappa}(I,J)= \{ p \, | \, |p| < \kappa \wedge \dom(p) \subseteq I \wedge \ran(p) \subseteq I\}$.
            \item Maximal element is $\emptyset$.
            \item The ordering is given by $q \leq p$ iff $p \subseteq q$.
        \end{itemize}
    \end{fancydefine}
    \begin{remark*}
        \hspace{0.1cm}
        \begin{itemize}
            \item $\Fn_{\omega}(I,J) = \Fn(I,J)$.
            \item $\Fn_{\kappa}(I,J)$ is not absolute for $\kappa> \omega$. Moreover, if $\calM$ is a CTM, then $\Fn_{\kappa}(I,J) \notin \calM$, so we need to consider $(\Fn_{\kappa}(I,J))^{\calM}$.
            \item For $\kappa > \omega$, $\Fn_{\kappa}(I,J)$ does not have the ccc. (of course, as long as $I,J \neq \emptyset$).
            \item If $G$ is $\Fn_{\kappa}(I,J)$-generic over $\calM$, then $f=\cup G$ is a function from $I$ to $J$.
        \end{itemize}
        :et $\bbP = \Fn_{\kappa}(\lambda \times \kappa,2)$, and let $\lambda \geq \kappa$, $\kappa$ is regular, and $\lambda^{\kappa}=\lambda$. By a similar argument to the $\omega$ case, \begin{align*}
            f = \cup G, &  & h_{\alpha} : \kappa \to 2 &  & h_{\alpha}(\beta) = f(\alpha, \beta)
        \end{align*}
        This is a sequence of $\lambda$-many distinct functions from $\kappa$ to $2$, by the exact same \say{nice} names argument from earlier, in $\calM[G]]$ there are precisely $\lambda$-many functions from $\kappa$ to $2$. So $\calM[G] \models 2^{\kappa} = \lambda$.
    \end{remark*}
    Our aim is to show that $\Fn_{\kappa}(I,J)$ preserves cardinals.
    \begin{fancydefine}{$\kappa$ chain condition}
        This is an analogue to the countable chain condition, we say that $\bbP$ has the $\kappa$\textbf{ chain condition} iff every antichain has cardinality less than $\kappa$
    \end{fancydefine}
    \begin{remark*}
        The ccc is the $\aleph_{1}$-cc.
    \end{remark*}
    \begin{fancydefine}{Preserving cofinalities $\geq \kappa$}
        We say that $\bbP$ \textbf{preserves cofinalities} $\geq \kappa$ iff for every $G$, $\cf^{\calM}(\gamma) = \cf^{\calM[G]}(\gamma)$ for all limit ordinals $\gamma \in \Ord \cap \calM$ with $\cf^{\calM}(\gamma)\geq \kappa$.
    \end{fancydefine}
    \begin{fancylem}{}
        Suppose that $\bbP$ is our forcing poset in $\calM$, $(\kappa \text{ is a regular cardinal})^{\calM}$, then \begin{enumerate}
            \item[(i)] $\bbP$ preserves cofinalities $\geq \kappa$ iff for all $\bbP$-generic filters $G$, for all limit ordinals $\beta$ with $\kappa \leq \beta < \Ord \cap \calM$: \[ (\beta \text{ is regular})^{\calM} \Rightarrow (\beta \text{ is regular})^{\calM[G]}.\]
            \item[(ii)] If $\bbP$ presrves cofinalities above $\kappa$, then $\bbP$ preserves cardinals above $\kappa$.
        \end{enumerate}
    \end{fancylem}
    \proof Omitted, it might be on the example sheet. \endproof
    \begin{fancylem}{}
        Let $A$, $B$, $\bbP \in \calM$, $(\kappa \text{ is regular})^{\calM}$, $(\bbP \text{ has the $\kappa$-cc})$, $G$ is $\bbP$-generic over $\calM$. Then for every $f: A \to B$ in $\calM[G]$, there is a function $F : A \to \calP(B)$ in $\calM$ such that for all $a \in A$: \begin{itemize}
            \item $f(a) \in F(A)$
            \item $(|F(a)|\leq \kappa)^{\calM}$
        \end{itemize}
    \end{fancylem}
    \proof Omitted. \endproof
    \begin{fancythm}{}
        If $\bbP \in \calM$, $(\kappa \text{ is regular})^{\calM}$ and $(\bbP \text{ has the $\kappa$-cc})^{\calM}$. Then $\bbP$ preserves cofinalities $\geq \kappa$ and hence cardinals $\geq \kappa$.
    \end{fancythm}
    \proof Omitted \endproof
    From the examples sheet, for any infinite cardinal $\kappa$, $\Fn_{\kappa}(I,J)$ has the $(|J|^{<\kappa})^{+}$-cc. In particular $\Fn_{\kappa}(\lambda \times \kappa, 2)$ has the $(2^{<\kappa})^{+}$-cc. We won't prove this in lecture, it will be on the examples sheet, but we will prove a weaker version that gives us what we need.
    \begin{fancylem}{}
        Let $\kappa$ be a regular cardinal, $(2^{<\kappa} = \kappa)^{\calM}$, $(1 \leq |J| \leq 2^{<\kappa})^{\calM}$, then $\bbP = \Fn_{\kappa}(I,J)^{\calM}$ has the $\kappa^{+}$-cc.
    \end{fancylem}
    \proof If $I = \emptyset$ then trivial, so assume $I \neq \emptyset$. Let $W$ be an antichain in $\bbP$, To show $|W| \leq \kappa$, we construct chains $\langle A_{\alpha} \, | \, \alpha < \kappa\rangle$ in $I$, and $\langle W_{\alpha} \, | \, \alpha < \kappa \rangle$ such that: \begin{itemize}
        \item[(i)] For $\alpha < \beta < \kappa$, $A_{\alpha} \subseteq A_{\beta} \subseteq I$, $W_{\alpha} \subseteq W_{\beta} \subseteq W$.
        \item[(ii)] For $\gamma$ a limit, $A_{\gamma} = \cup_{a < \gamma} A_{\alpha}$, $W_{\gamma} = \cup_{\alpha < \gamma} W_{\alpha}$.
        \item[(iii)] $W = \cup_{\alpha < \kappa} W_{\alpha}$
        \item[(iv)] For all $\alpha < \kappa$, $|W_{\alpha}|, \, |A_{\alpha}| \leq \kappa$.
    \end{itemize}
    Assuming we can do this, by the regularity of $\kappa^{+}$, $|W| \leq \kappa$. \gap
    How do we build this? Well, we'll start by setting $A_{0} = W_{0} = \emptyset$. Then assume that $A_{\alpha}, W_{\alpha}$ are defined. Then for each $p \in \bbP$ with $\dom(p) \subseteq A_{\alpha}$, using AC, choose $q_{p} \in W$ such that $p = q_{p} \upharpoonright A$ (if it exists, if it doesn't we just move on and everything is fine). \gap
    Then $W_{\alpha+1} = W_{\alpha} \cup \{ q_{p} \, | \ \dom(p) \subseteq A_{\alpha}\}$. If $\dom(p)$ were contained in $A_{\beta}$, with $\beta < \alpha$, then the $q_{p}$ we pick here should be the same as the $q_{p}$ that was picked at stage $\beta$. \gap Then define $A_{\alpha+1} = \cup \{ \dom(q) \, | \, q \in W_{\alpha+1}\}$. \gap Finally, set $A_{\gamma} = \cup_{\alpha< \gamma} A_{\alpha}$, $W_{\gamma} = \cup_{\alpha< \gamma} W_{\alpha}$ when $\gamma$ is a limit, and then define $A = \cup_{\alpha < \kappa} A_{\alpha}$.\gap
    \textbf{Claim:} $W = \cup_{\alpha < \kappa} W_{\alpha}$. \gap
    By construction, $\cup_{\alpha < \kappa} W_{\alpha} \subseteq W$. Fix $q \in W$. Firstly, $\dom(q) \cap A \neq \emptyset$. Otherwise, take $q_{1} \in W_{1}$. Then $\dom(q_{1}) \subseteq A$. So if $\dom(q_{1}) \cap \dom(q) = \emptyset$, then $\dom(q_{1}) \, || \, \dom(q)$, contradicting $q_{1}, q \in W$. \gap
    Since $\dom(q) \cap A = \emptyset$ and $|\dom(q)| < \kappa$, $\dom(q) \cap A = \dom(q) \cap A_{\alpha}$ for some $\alpha$. \gap
    Now define $p = q \upharpoonright A_{\alpha}$. Then there is some $q' \in W_{\alpha+1}$ such that $q' \upharpoonright A_{\alpha}=p$. Since $\dom(q') \subseteq A$, so $q \, || \, q'$. Since $W$ is an antichain, this is only possible if $q = q'$.\gap
    \textbf{Claim:} For all $\alpha < \kappa$, $|W_{\alpha}|, \, |A_{\alpha}| \leq \kappa$. Proven by induction. \begin{itemize}
        \item The limit cases follow by regularity.
        \item If $|W_{\alpha}| \leq \kappa$, then $|A_{\alpha}| \leq \kappa$ because $q$ lie in $\Fn_{\kappa}(I,J)$, the $\kappa$ is the important part here.
        \item Suppose $|W_{\alpha}| \leq \kappa$. Then since every $q$ that is added to $W_{\alpha}$ is chosen from some $p \in \bbP$, $\dom(p) \subseteq A_{\alpha}$, so \[ |W_{\alpha+1}| \leq |W_{\alpha}| + |\{ p \in \bbP \, | \, \dom(p)\subseteq A_{\alpha}\} | \]
        \item Since $|A_{\alpha}| \leq \kappa$, and $|\dom(p)| < \kappa$ %What?
              $|[A_{\alpha}]^{<\kappa}| \leq \kappa^{<\kappa} = 2^{<\kappa} = \kappa$.
    \end{itemize}
    So we have $\bbP$ = $\Fn_{\kappa}(\lambda \times \kappa, 2)$, $\calM[G] \models 2^{\kappa} = \lambda$ and all cardinals $\geq \kappa^{+}$ are preserved.
    \newlec
    \subsection{Closure and Distributivity}
    \begin{fancydefine}{$\kappa$-closed}
        A poset $\bbP$ is $\kappa$\textbf{-closed} iff $\forall \delta < \kappa$, every decreasing sequence $\langle p_{\alpha} \, | \, \alpha < \delta \rangle$ in $\bbP$ has a lower bound. \gap
        I.e. for all $\alpha < \beta < \delta$, if $p_{\beta}<p_{\alpha}$ then $\exists q, \, \forall \alpha < \delta.\, q \leq p_{\alpha}$
    \end{fancydefine}
    \begin{fancydefine}{$<\kappa$-closed}
        $\bbP$ is $<\kappa$-closed iff it is $\lambda$-closed for all $\lambda < \kappa$ cardinals. That is, any decreasing sequence of conditions $\langle p_{\alpha} \, | \alpha < \gamma \rangle$ has a lower bound.
    \end{fancydefine}
    \begin{fancydefine}{$<\kappa$-distributive}
        $\bbP$ is $<\kappa$\textbf{-distributive} iff for every $\gamma < \kappa$, the intersection of $<\kappa$-many open dense sets is open and dense. Remember that: \begin{itemize}
            \item $D$ is dense if $\forall p in \bbP \exists q \in D. \, q \leq p$.
            \item $D$ is open if $\forall p \in D, \, \forall q \in \bbP. \, (q \leq p \rightarrow q \in D)$.
        \end{itemize}
    \end{fancydefine}
    \begin{fancylem}{}
        $\bbP$ is $<\kappa$-closed $\Rightarrow$ $\bbP$ is $<\kappa$-distributive.
    \end{fancylem}
    \proof Omitted \endproof
    \begin{fancylem}
        If $(\kappa \text{ is regular})^{\calM}$ then $\Fn_{\kappa}(I,J)^{\calM}$ is $<\kappa$-closed.
    \end{fancylem}
    \proof Omitted \endproof
    \begin{fancythm}{}
        Let $A, B, \bbP \in \calM$, $(\kappa \text{ is a cardinal})^{\calM}$, $(|A| < \kappa)^{\calM}$, $(\bbP \text{ is $\kappa$-distributive})^{\calM}$. Let $G$ be generic, then any function $f \in \calM[G]$ with $f : A \to B$ has $f \in \calM$.
    \end{fancythm}
    \proof Suffices to prove the statement for $A= \delta$ when $\delta < \kappa$. Suppose $\calM[G] \models f : \delta \to \beta$. By the forcing theorem, fix $p \in G$ such that $p \forces f : \check{\delta} \to \check{\beta}$. $\dot{f}$ is a name for $f$.\gap
    For each $\alpha < \delta$, let $D_{\alpha} = \{ q \leq p \, | \, \exists x \in B. \, q \forces \dot{f}(\check{\alpha}) = \check{x}\}$. This is open and dense, so since $\bbP$ is $\kappa$-distributive, we can take $D = \cap_{\alpha < \kappa} D_{\alpha}$ is open and dense below $p$ (open is a technical condition which makes things work, the key point is that it is dense). So fix $q \in D \cap G$. To argue this is in $\calM$, for each $\alpha < \delta$, choose $x_{\alpha} \in \beta$ such that $q \forces f(\check{\alpha}) = \dot{x}_{\alpha}$. Now define \begin{align*}
        g : \delta & \to \beta          \\
        \alpha     & \mapsto x_{\alpha}
    \end{align*}
    Then $g \in \calM$. But for any $\alpha < \delta$, we have \[q \forces \dot{f}(\check{\alpha})=\check{x}_{\alpha} = \check{g}(\check{\alpha})\]. So $\calM[G] \models f=g$. So $f \in \calM$.
    \endproof
    \begin{fancythm}{}
        Let $I, J, \kappa \in \calM$, suppose $(\kappa \text{ is regular})^{\calM}$, $(2^{<\kappa}=\kappa)^{\calM}$, and $(|J| \leq \kappa)^{\calM}$. Then $\Fn_{\kappa}(I,J)^{\calM}$ preserves cofinalities and hence cardinals.
    \end{fancythm}
    \proof Recall, it suffices to show that for every limit ordinal $\beta < \Ord \cap \calM$, if $\beta$ is regular in $\calM$, then $\beta$ is regular in $\calM[G]$. There are two cases to consider: \begin{enumerate}
        \item[($\beta > \kappa$)] Since $(|J| \leq \kappa = 2^{<\kappa})^{calM}$., $\Fn_{\kappa}(I,J)^{\calM}$ has the $\kappa^{+}$-cc. So $\Fn_{\kappa}(I,J)$ oreserves all cardinals and cofinalities above $\kappa^{+}$. In particular, if $(\beta \text{ is regular})^{\calM}$ then $(\beta \text{ is regular})^{\calM[G]}$.
        \item[$(\beta \leq \kappa)$] We will show that if $(\beta \text{ is singular})^{\calM[G]}$, then $(\beta \text{ is singular})^{\calM}$. In $\calM[G]$, fix $\delta < \beta$, $f@ \delta \to \beta$ cofinal. Since $\calM$ is transitive, $\delta \in \calM$. Since $\bbP$ is $\kappa$-closed, $f \in \calM$. So $(\beta \text{ is singular})^{\calM}$. Then we are done by the contrapositive.
    \end{enumerate}
    \endproof
    \begin{fancythm}{}
        Suppose $(\kappa, \lambda \text{ are cardinals})^{\calM}$, that $\aleph_{0} \leq \kappa \leq \lambda$, and that $(\kappa \text{ is regular})^{\calM}$, $(2^{<\kappa}=\kappa)^{\calM}$, $(\lambda^{\kappa} = \lambda)^{\calM}$. Let $\bbP = \Fn_{\kappa}(\lambda \times \kappa, 2)$, $G$ be $\bbP$-generic over $\calM$. Then $\bbP$ preserves cardinals and $\calM[G] \models 2^{\kappa} = \lambda$.
    \end{fancythm}
    \proof Comes from the above \endproof
    \subsection{Fixing multiple values of the continuum}
    \begin{fancythm}{}
        Suppose there is a countable transitive model of $\sfZFC + \sfGCH$, then there is a countable transitive model of $\sfZFC$ satisfying any of the following statements: \begin{enumerate}
            \item[(i)] $\sfCH + 2^{\aleph_{1}}=\aleph_{3}$
            \item[(ii)] $2^{\aleph_{0}}=2^{\aleph_{1}}=\aleph_{5} + 2^{\aleph_{2}}=\aleph_{\omega+5}$.
            \item[(iii)] For a fixed $n \in \omega$, for all $m \leq n$, $2^{\aleph_{m}}=\aleph_{2m+3}$.
        \end{enumerate}
    \end{fancythm}
    \proof Let $\calM$ be a ctm of $\sfZFC + \sfGCH$. Then \begin{enumerate}
        \item[(i)] Let $\bbP = \Fn_{\aleph_{1}}(\omega_{3} \times \omega_{1},2)$. If $G$ is $\bbP$-generic, then \[ \calM[G] \models 2^{\aleph_{1}}=\aleph_{3}.\] Also, as $\bbP$ is $\omega_{1}$-closed, it doesn't add new functions from $\omega$ to $2$, so $(\sfCH)^{\calM[G]}$.
        \item[(ii)] First, force $\bbP_{0} = \Fn_{\aleph_{2}}(\omega_{\omega+5} \times \omega_{2}, 2)^{\calM}$. Let $G_{0}$ be $\bbP_{0}$-generic. By closure, $(2^{<\aleph_{1}} = \aleph_{1})^{\calM[G_{0}]}$, $(\aleph_{5}^{\aleph_{1}}=\aleph_{5})^{\calM[G_{0}]}$. Now let $\bbP_{1} = \Fn_{\aleph_{0}}(\omega_{5} \times \omega, 2)^{\calM[G_{0}]}$. Let $G_{1}$ be $\bbP_{1}$-generic, then $\calM[G_{1}] \models 2^{\aleph_{0}}=2^{\aleph_{1}} =\aleph_{5}$ as well as $2^{\aleph_{2}}=\aleph_{\omega+5}$. Recall a result that if we make the continuum have size (e.g.) $aleph_{5}$, and we are beginning from a model of $\sfGCH$ then for any cardinal $ \aleph_{0} \leq \lambda <\kappa$ in $\calM$, in $\calM[G]$ we have $2^{\lambda} = \kappa$ or $2^{\lambda}=\kappa^{+}$ depending on $\cf(\kappa)$.
        \item[(iii)] Similar.
    \end{enumerate}
    \endproof
    \begin{remark*}\hspace{0.1cm}
        \begin{itemize}
            \item It is necessary to start at the largest cardinals and work backwards for cardinal arithmetic to work.
            \item Our iterative approach works for any \textit{finite} number of cardinals.
            \item  See later how to get forcings of the sort $2^{\aleph_{n}}=2^{\aleph_{2n+3}}$ for all $n \in \omega$, i.e. infinitary forcings.
        \end{itemize}
    \end{remark*}
    \begin{fancyprop}{}
        Suppose $\calM$ is a countable transitive model of $\sfZFC$, $\calM \models 2^{\aleph_{0}}= \aleph_{\alpha}$ for $\alpha \geq 1$. Let $\bbP = \Fn_{\aleph_{1}}(\kappa \times \aleph_{1},2)$. Then if we let $G$ be $\bbP$-generic, in $\calM[G]$ we get a model of $\sfCH$, and all cardinals $\delta$ of $\calM$ with $(\aleph_{1} \leq \alpha \leq \aleph_{\alpha})^{\calM}$. \gap
        In particular $\aleph_{\alpha}^{\calM} \neq \aleph_{\alpha}^{\calM[G]}$. This is because we can encode subsets of $\aleph_{0}$ into the forcing which we are using here.
    \end{fancyprop}
    \proof Example sheet 4 (I think).\gap
    \newlec
    \textbf{Observation:} If $(\delta \text{is a cardinal})^{\calM}$, and $(\delta > |\bbP|)^{\calM}$, then $(\delta \text{is a cardinal})^{\calM[G]}$. Observe also that $\bbP$ has the $|\bbP|^{+}$-cc.\gap
    Recall that $p \forces \exists x \varphi(x)$ iff $\{q \leq p \, | \, \exists \dot{x} \in V^{\bbP}. \, q \forces \varphi(x)\}$ is dense below $p$. In most cases, the witness $\dot{x}$ doesn't depend on the filter. For example, $p \forces \exists x (\dot{a} \in x \wedge \dot{b} \in x)$. We don't need the generic to find a name as $\dot{x} = \op(\dot{a},\dot{b})$.
    \begin{fancylem}{The mixing lemma, $\sfZFC$}
        Suppose $(p \forces \exists x \varphi(x))^{\calM}$, then $\exists \dot{x} \in \calM^{\bbP}$ such that $(p \forces \varphi(\dot{x}))^{\calM}$.
    \end{fancylem}
    \proof Since the set $\{ q \leq p \, | \, \exists x \in \calM^{\bbP}. \, q \forces \varphi(x)\}$ is dense below $p$, it contains a maximal antichain, $D$ (this requires choice). Now, for each of our conditions $q \in D$, choose $\dot{x}_{q} \in \calM^{\bbP}$ such that $q \forces \varphi(\dot{x}_{q})$. Without loss of generality, we may assume that if $\langle r, \dot{y} \rangle \in \dot{x}_{q}$, then $r \leq q$. This is because if $r \nleq q$ then there are two cases: \begin{itemize}
        \item $r \perp q$, but if this is the case, then $q \forces \dot{x}_{q}=\dot{x}_{q} \setminus \langle r, \dot{y} \rangle$.
        \item If $r || q$, then define $\dot{x}_{q}' = (\dot{x}_{q} \setminus \langle r,\dot{y} \rangle) \cup \{ \langle s, \dot{y} \rangle \, | \, s \leq r,q\}$. Again, $q \forces \dot{x}_{q} = \dot{x}_{q}'$.
    \end{itemize}
    Now if $q,q' \in D$ and $q \neq q'$, then $q \perp q'$ so $q' \forces \dot{x}_{q} = \emptyset$. Next, let $\dot{x} = \cup \{ \dot{x}_{q} \, | \, q \in D \}$. Then if $q \in D$, $q \forces \dot{x} = \dot{x}_{q}$. By the Forcing theorem, $q \forces \varphi(\dot{x})$ (since it forces it for $\dot{x}_{q}$, and that $\dot{x} = \dot{x}_{q}$). \gap
    It remains to prove that $p \forces \varphi(\dot{x})$. Suppose not, then we can fix an $r \leq p$ such that $r \forces \neg \varphi(\dot{x})$. Since $D$ is a \textit{maximal} antichain, we can fix $q \in D$ such that $q || r$. But now if we take $s \leq q,r$, we have $s \forces \varphi(\dot{x})$, since it is below $q$, and $s \forces \neg \varphi(\dot{x})$, since it is below $r$. Contradiction, so $p \forces \varphi(\dot{x})$. \endproof
    \textbf{Question:} Given cardinals $\kappa \leq \lambda$ in $\calM$, can we find a generic $G$ (for some $\bbP$) such that $\calM[G] \models \lambda = \kappa^{+}$. \gap
    First, observe that in order to find such a generic, then $(\lambda \text{ must be regular})^{\calM}$. If $f : \alpha \to \lambda$ cofinal, $\alpha < \lambda$, $f \in \calM$. Then $f \in \calM[G]$, so $\cf^{\calM[G]}(\lambda) \leq  \cf^{\calM[G]}(\alpha) \leq |\alpha|^{\calM[G]} < \lambda$. However, this is the only restriction.
    \begin{fancythm}{}
        Let $(\kappa \text{ be regular})^{\calM}$, $(\delta > \kappa)$, $(\delta \text{ is a cardinal})^{\calM}$. Let $(\lambda = \delta^{+})^{\calM}$. Let $G$ be $\Fn_{\kappa}(\kappa, \delta)$-generic over $\calM$. Then in $\calM[G]$: \begin{itemize}
            \item $|\delta| = \kappa$.
            \item Every cardinal $\alpha \leq \kappa$ in $\calM$ remains a cardinal.
            \item If $\delta^{<\kappa} = \delta$, then every cardinal $\alpha$ such that $\alpha > \delta$ in $\calM$ remains a cardinal in $\calM[G]$.
        \end{itemize}
        In particular, if $\delta^{<\kappa}= \delta$, then $\calM[G] \models \lambda = \kappa^{+}$.
    \end{fancythm}
    \proof $\cup G : \kappa \to \delta$ is a surjection, so $|delta| = |\kappa|$. So there are no cardinals between $\delta$ and $\lambda$. \gap
    Since $\kappa$ is regular, $\Fn_{\kappa}(\kappa, \delta)$ is $<\kappa$-closed, so every cardinal $\alpha \leq \kappa$ is preserved. \gap
    And finally, if $\delta^{<\kappa}=\delta$, then $|\Fn_{\kappa}(\kappa, \delta)|=\delta$. So $\Fn_{\kappa}(\kappa,\delta)$ has the $\delta^{+}$-cc. (i.e. $\lambda$-cc). Thus every cardinal $\alpha > \delta$ (in particular, $\lambda$) is preserved. \endproof
    The other case is that $\lambda$ is weakly inaccessible. Here we use a forcing callued the L\'evy collapse.
    \begin{fancydefine}{$\Col(\kappa, < \lambda)$}
        Let $\lambda > \kappa$ be infinite ordinals, and let $\Col(\kappa, <\lambda)$ consist of all functions $p$ such that: \begin{itemize}
            \item $p$ is a partial function from $\kappa \times \lambda$ to $\lambda$.
            \item $|\dom(p)| < \kappa$.
            \item $p(\alpha, \beta) < \beta$ for each $(\alpha, \beta) \in \dom(p)$.
        \end{itemize}
        And we say $q \leq p$ iff $q$ extends $p$ as a function
    \end{fancydefine}
    \begin{fancythm}{}
        Let $\kappa$ be regular, and suppose $\lambda > \kappa$ is weakly inaccessible. Let $G$ be $\Col(\kappa, <\lambda)$-generic over $\calM$. Then in $\calM[G]$: \begin{itemize}
            \item Every ordinal $\beta$ with $\kappa \leq \beta < \lambda$ has cardinality $\kappa$.
            \item Every cardinal $\leq \kappa$ and $\geq \lambda$ remains a cardinal.
        \end{itemize}
        Hence $\calM[G] \models \lambda = \kappa^{+}$.
    \end{fancythm}
    \proof For each $\beta < \lambda$, define $G_{\beta} : \kappa \to \beta$ by $G_{\beta}(\alpha) = G(\alpha,\beta)$. So if $\kappa \leq \beta < \lambda$, $\calM[G] \models |\beta| = |\kappa|$. Next $\Col(\kappa, <\lambda)$ is $<kappa$-closed, so it preserves cardinals $\leq \kappa$. Finally ,$|\Col(\kappa, <\lambda)| = \lambda$, so $\Col(\kappa, < \lambda)$ has the $\lambda$-cc, and preserves cardinals $\geq \lambda$\footnote{Actually, we only have $\lambda^{+}$-cc right now, but there is an argument which shows it preserves cardinals $\geq \lambda$.} \endproof
    \begin{remark*}
        $\lambda$ is weakly compact iff it is inaccessible and satisfies something called the \textbf{tree property}.
    \end{remark*}
    \textbf{Claim:} If $G$ is $\Col(\aleph_{0}, <\lambda)$-generic, then we have that $\calM[G] \models \aleph_{1}\text{ has the tree property}$.
    \begin{remark*}
        This shows that $\lambda$ being a limit cardinal is not absolute between transitive models $\calM$ and $\calN$, \textit{even if} $\lambda$ being a cardinal is!
    \end{remark*}
    The final question we want to address in this course is: Is the following consistent: \[ \sfZFC + \forall n \in \omega. \, 2^{\aleph_{n}} = \aleph_{2n+3}.\]
    \begin{remark*}
        We have the consistency of $\sfZFC + \forall n \leq k. \, 2^{\aleph_{n}}=2^{\aleph_{2n+3}}$, for any $k \in \omega$, because we can do arbitrarily finite forcing, but we can't yet do infinite forcings.
    \end{remark*}
    The answer to this question (spoiler) is yes! We use Easton's forcing. Then our follow-up is: What are the restrictions on the function $F : \Card \to \Card$, where $F(\aleph_{\alpha})=2^{\aleph_{\alpha}}$?\gap
    \newlec
    \begin{fancythm}{}
        Let $\kappa$ be regular, $\lambda > \kappa$ is strongly inaccessible, $G$ a $\Col(\kappa, <\lambda)$-generic. Then $\calM[G] \models \lambda = \kappa^{+}$.
    \end{fancythm}
    \begin{remark*}
        Suppose $\lambda$ were weakly inaccessible, and $2^{\aleph_{0}}>\lambda$, then $\Col(\aleph_{1},<\lambda)$ has an antichain of length $2^{\aleph_{0}}$.
    \end{remark*}
    \proof For $A \subseteq \omega$, define $p_{A} : \{w\} \times [\omega, \omega + \omega)$ by: \[ p_{A}(\alpha, \omega + n) = \begin{cases}
            0, & n \in A    \\
            1, & n \notin A
        \end{cases}\]
    Then if $A \neq B$, then $p_{A} \perp p_{B}$.
    From last time, we know that all cardinals $\delta$ with $\kappa < \delta < \lambda$ are collapsed. We need to prove that $\lambda$ is a cardinal, which follows from the $\lambda$-cc. Given $p \in \Col(\kappa, <\lambda)$, let $\sp(p)= \{ \beta \, : \, \exists \alpha. \, (\alpha, \beta) \in \dom(p)\}$. So $|\sp(p)| < \kappa$.\gap
    Let $W$ be an antichain. Consider chains $\langle A_{\alpha} \, | \, \alpha < \kappa\rangle$ and $\langle W_{\alpha} \, | \, \alpha < \kappa\rangle$ such that: \begin{enumerate}
        \item for $\alpha < \beta < \kappa$, $A_{\alpha} \subseteq A_{\beta}$, $W_{\alpha} \subseteq W_{\beta}$.
        \item If $\gamma$ is a limit $\cup_{\alpha} A_{\alpha} = A_{\gamma}$, $\cup_{\alpha < \gamma} W_{\alpha} = W_{\gamma}$.
        \item $W = \cup_{\alpha < \kappa} W_{\alpha}$
        \item For all $\alpha < \kappa$, $|A_{\alpha}|, |W_{\alpha}| < \lambda$.
    \end{enumerate}
    Assuming this, by regularity of $\lambda$, $|W| = |\cup_{\alpha<\kappa} W_{\alpha}| < \lambda$. \gap
    To find these sets, set $A_{0},W_{0}=\emptyset$. Assume $A_{\alpha}, B_{\alpha}$. Then for $p \in \Col(\kappa, <\lambda)$ with $\sp(p) \subseteq A_{\alpha}$, using the axiom of choice, choose $q_{p} \in W$ such that $p = q_{p} \upharpoonright (\kappa \times \sp(p))$ (if it exists).\gap
    Then $W_{\alpha+1} = \{ q_{p} \, | \, \sp(p) \subseteq A_{\alpha}\}$. And $A_{\alpha+1}= \cup\{\sp(q) \, | \, q \in W_{\alpha}\}$. \gap
    \textbf{Claim:} $W = \cup_{\alpha < \kappa} W_{\alpha}$.\gap
    The proof of this is the same as the proof in the $\Fn_{\kappa}(I,J)$ proof. \gap
    \textbf{Claim:} For $\alpha < \kappa$, $|W_{\alpha}|,|A_{\alpha}|<\lambda$.\gap
    We prove this by induction on $\alpha$.\begin{itemize}
        \item The limit cases follow by regularity.
        \item If $|W_{\alpha+1}| < \lambda$, then $|A_{\alpha+1}| \leq \kappa \cdot \lambda = \lambda$, since it is a union of less than lambda many things each of size less than $\kappa$.
        \item If $|W_{\alpha}| < \lambda$, then since every $q$ that is added to $W_{\alpha}$ is chosen from some $p \in \Col(\kappa, <\lambda)$ with $\sp(p) \subseteq A_{\alpha}$. So $|W_{\alpha+1}| \leq |A_{\alpha}|^{<\kappa}$. Since $\lambda$ is strongly inaccessible, this must be less than $\lambda$.
    \end{itemize}
    \endproof
    \begin{corr}{}
        If $\Con(\sfZFC + \exists \text{inacc. cardinal.})$, then $\Con(\sfZFC + \aleph_{1}^{V} \text{ is inaccessible in $L$})$.
    \end{corr}
    \proof If we start with $V=L$, and let $G$ be $\Col(\omega_{1},<\lambda)$-generic. Then $\calM[G] \models \lambda = \aleph_{1}$, so $\calM[G] \models (\lambda \text{ is inacc.})^{L}$. \endproof
    \begin{remark*}
        If $V \models \sfZFC + \kappa \text{ is measurable}$, then $\aleph_{1}^{V}$ is inaccessible in $L$.
    \end{remark*}
    \subsection{Product Forcing}
    \begin{fancydefine}{Product Order}
        Suppose we have two forcing notions $(\bbP, \leq_{\bbP})$, and $(\bbQ, \leq_{\bbQ})$, then we can define the \textbf{product order}: $\leq$ on $\bbP \times \bbQ$, by \[ \langle p_{1}, q_{1} \rangle \leq \langle p_{2}, q_{2} \rangle \text{ iff } p_{1} \leq_{\bbP} p_{2} \text{ and } q_{1} \leq_{\bbQ} q_{2}\]
    \end{fancydefine}
    \begin{fancydefine}{Projections}
        Given a $\bbP \times \bbQ$-generic filter, $G$, over $\calM$. Let $G_{0} = \{ p \in \bbP \, | \, \exists q \in \bbQ. \, (p,q) \in G \}$, and $G_{1} = \{q \in \bbQ \, | \, \exists p \in \bbP. \, (p,q) \in G \}$. These are called the \textbf{projections} of $G$.
    \end{fancydefine}
    \begin{fancylem}{}
        Suppose $\calM$ is a transitive model of $\sfZFC$, and $\bbP, \bbQ \in \calM$. Let $G \subseteq \bbP$ and $H \subseteq \bbQ$. Then TFAE: \begin{enumerate}
            \item $G \times H$ is $\bbP \times \bbQ$-generic over $\calM$.
            \item $G$ is $\bbP$-generic over $\calM$, and $H$ is $\bbQ$-generic over $\calM[G]$.
            \item $H$ is $\bbQ$-generic over $\calM$, and $G$ is $\bbP$-generic over $\calM[H]$.
        \end{enumerate}
        Moreover, if (i) holds, then $\calM[G \times H] = \calM[G][H] = \calM[H][G]$.
    \end{fancylem}
    \proof The first part is on example sheet 4. The idea is to prove first that they are filters, and then prove genericity, for both (i) implies (ii) and (ii) implies (i). To prove the last claim: \gap
    Recall the last part of the Generic Model Theorem which said that if $\calM$ is transitive, $\calM \models \sfZF$, $\bbP \in \calM$, $G \subseteq \bbP$, $\bbone \in G$. Then if $\calN$ is a transitive model of $\sfZF$, $\calM \subseteq \calN$ is a definable class in $\calN$, then $\calM[G] \subseteq \calN$. \gap
    Since $M \subseteq \calM[G][H]$, and $G \times H \in \calM[G][H]$. Then $\calM[G \times H] \subseteq \calM[G][H]$.\gap
    For the other direction, $G \in \calM[G \times H]$, and $\calM \subseteq \calM[G \times H]$, so $\calM[G] \subseteq \calM[G \times H]$, and then since also $H \in \calM[G \times H]$, $\calM[G][H] \subseteq \calM[G \times H]$.
    \endproof
    \section{Easton Forcing}
    ** Proof Non-examinable **.\gap
    Return to the model of $2^{\aleph_{0}}=\aleph_{3}$ and $2^{\aleph_{1}}=\aleph_{5}$. We started with $\calM \models \sfZFC + \sfGCH$, $\Fn(\omega_{3} \times \omega, 2)^{\calM}$, $G_{0}$ generic over this. And $\Fn_{\omega_{1}}(\omega_{5} \times \omega_{1}, 2)^{\calM[G_{0}]}$. Then $\calM[G_{0}][G_{1}] \models \sfCH$. So we did it the other way around:\gap
    We took $\bbP_{0} = \Fn_{\omega_{1}}(\omega_{5} \times \omega_{1},2)^{\calM}$, and had $G_{0}$ generic over that. Then we took $\bbP_{1} = \Fn(\omega_{3} \times \omega, 2)^{\calM{G_{0}}}$, and had $G_{1}$ generic over that. Then $\calM[G_{0}][G_{1}]$ gives us the model we want. \gap
    However, $\bbP_{0}$ is $<\omega_{1}$ closed, so it doesn't add new sequences of length $\omega$. Thus $\bbP_{1} = \Fn_{\omega_{3} \times \omega,2}^{\calM}$, the smaller forcing isn't changed in $\calM[G_{0}]$. So we can define the forcing $\bbP_{0} \times \bbP_{1}$ over $\calM$, and $G_{0} \times G_{1}$ is $\bbP_{0} \times \bbP_{1}$-generic over $\calM$. \gap
    This leads to an obvious candidate for $2^{\aleph_{n}}=\aleph_{2n+3}$: \[ \bbP = \prod_{n < \omega} \Fn_{\omega_{n}}(\omega_{2n+3} \times \omega_{n}, 2).\]
It turns out that this works.
\begin{fancythm}{Easton (Set version)}
    Let $\calM$ be a countable transitive model of $\sfZFC + \sfGCH$. Let $S$ be a set of regular cardinals in $\calM$, and let $F : S \to \Card \cap \calM$ be a function in $\calM$ such that for all $\kappa \leq \lambda$ in $S$, the following hold: \begin{enumerate}
        \item[(i)] $F(\kappa) > \kappa$.  (restriction comes from Cantor)
        \item[(ii)] $F(\kappa) \leq F(\lambda)$. (restriction comes from montonicity)
        \item[(iii)] $\cf(F(\kappa)) > \kappa$. (restriction comes from K\"onig)
    \end{enumerate}
    Then there is a generic extension $\calM[G]$ of $\calM$ such that $\calM$ and $\calM[G]$ have the same cardinals, and for all $\kappa \in S$, $\calM \models 2^{\kappa}=F(\kappa)$.
\end{fancythm}
\begin{remark*}
    This can be generalised by a virtually identical proof to the class of regular cardinals. This, however, needs a \say{class-forcing}, so $\bbP$ needs to be a proper class. The big issue with class forcings is that when $\bbP$ is a proper class, it is nontrivial to show $\calM[G] \models \sfZFC$, and in fact it sometimes does not model $\sfZFC$. For example $\Fn(\Ord \times \omega, 2)$ makes $2^{\aleph_{0}}$ a proper class. Or $\Fn(\omega,\Ord)$ which adds a surjection from $\omega$ to $\Ord$. \gap
    In fact, $\forces$ may not be defined.
\end{remark*}








\end{document}